% Les cours à distance : les problèmes liés à l'utilisation du numérique — Allemand Tle A — Cours signé OnBuch+
% Programme officiel MINESEC. Compiler avec : tectonic AL20-les-cours-a-distance-les-problemes-lies-a-l-utilisation-du-n.tex
\newcommand{\DOCMATIERE}{Allemand}
\newcommand{\DOCNIVEAU}{Tle A}
\newcommand{\DOCMODULE}{Chapitre 4 : Média et communication}
\newcommand{\DOCLECON}{AL20}
\newcommand{\DOCTITRE}{Les cours à distance : les problèmes liés à l'utilisation du numérique}
% OnBuch+ — préambule « cours signés » ENRICHI (compilé avec tectonic / XeLaTeX)
% Système visuel « Pop » d'OnBuch+ : fond crème (jamais blanc), bordures encre
% épaisses, ombres « dures » décalées, titres Archivo Black en capitales.
% Version MATHÉMATIQUES : graphes (pgfplots/tikz), boîtes pédagogiques
% (définition, propriété, méthode, attention, le savais-tu, exercice résolu,
% exercice, TP, bilan), page de garde et signature OnBuch+ ancrée en bas.
%
% Macros attendues dans main.tex AVANT \input{preamble} :
%   \DOCMATIERE  \DOCNIVEAU  \DOCMODULE  \DOCLECON (numéro)  \DOCTITRE
\documentclass[11pt]{article}

\usepackage{fontspec}
\usepackage[ngerman,french]{babel} % français = langue principale ; l'allemand via \all{..} / textebox
\usepackage[a4paper,margin=2cm,top=3.4cm,bottom=2.8cm,headheight=1.4cm,footskip=1.3cm]{geometry}
\usepackage{amsmath,amssymb,amsfonts}
\usepackage{mathtools} % \xmapsto, \coloneqq... souvent utilisés par les modèles
\usepackage{cancel} % simplifications barrées (bilans de forces)
% \abs{z} (module, valeur absolue) et \norm{u} (norme), utilisés par les modèles
\providecommand{\abs}{}\let\abs\relax\DeclarePairedDelimiter{\abs}{\lvert}{\rvert}
\providecommand{\norm}{}\let\norm\relax\DeclarePairedDelimiter{\norm}{\lVert}{\rVert}
\usepackage[table]{xcolor} % [table] : fournit aussi \rowcolors (alternance de lignes)
\usepackage{pagecolor}
\usepackage{tikz}
\usetikzlibrary{arrows.meta,positioning,calc,shapes.geometric,shapes.misc,shapes.arrows,decorations.pathmorphing,decorations.pathreplacing,decorations.markings,patterns,shadows,mindmap,trees,fit,backgrounds,matrix,intersections,angles,quotes}
\usepackage{pgfplots}
\usepackage[version=4]{mhchem} % équations nucléaires \ce{^{235}_{92}U}
\usepackage[european,siunitx]{circuitikz} % schémas électriques
\pgfplotsset{compat=1.18}
\usepgfplotslibrary{fillbetween,groupplots} % aires entre courbes (calcul intégral)
\usepackage{enumitem}
\usepackage{fancyhdr}
% [most] charge toutes les bibliothèques tcolorbox (listings, minted, raster,
% documentation...) et gonfle inutilement la mémoire TeX sur un document long
% (~300 boîtes) : on ne charge que ce qui est réellement utilisé.
\usepackage[skins,breakable]{tcolorbox}
\usepackage{graphicx}
\usepackage{colortbl}
\usepackage{booktabs}
\usepackage{tabularx}
\usepackage{diagbox} % case d'en-tête diagonale des tableaux à double entrée
% Colonnes extensibles centrée / gauche / droite (C, L, R), souvent utilisées par les modèles
\newcolumntype{C}{>{\centering\arraybackslash}X}
\newcolumntype{L}{>{\raggedright\arraybackslash}X}
\newcolumntype{R}{>{\raggedleft\arraybackslash}X}
\usepackage{array}
\usepackage{multicol}
\usepackage{multirow}
\usepackage{siunitx}
% parse-numbers=false : accepte aussi \SI{2,5 \times 10^{-3}}{...}
\sisetup{locale=FR,output-decimal-marker={,},group-minimum-digits=4,parse-numbers=false}
\DeclareSIUnit\ohm{\ensuremath{\symup{\Omega}}} % U+2126 absent de la police
\DeclareSIUnit\division{div} % réglages d'oscilloscope (ms/div, V/div)
\DeclareSIUnit\tr{tr} % tours (vitesse de rotation, ex. \SI{1500}{\tr\per\minute})
\DeclareSIUnit\molar{M} % concentration molaire (ex. \SI{3}{\molar}, \SI{1}{\micro\molar})
\DeclareSIUnit\an{an} % année (le modèle confond parfois avec \year, un compteur TeX) — ex. \SI{5}{\kilo\meter\per\an}
\DeclareSIUnit\FCFA{FCFA} % franc CFA (ex. \SI{500}{\FCFA\per\kilo\gram})
\DeclareSIUnit\degreeBrix{\textdegree Brix} % teneur en sucre d'un sirop/jus (réfractométrie)
\DeclareSIUnit\month{mois} % le modèle utilise parfois \month, un compteur TeX, comme unité de durée
\usepackage{qrcode}
% \degree : commande fréquemment utilisée par les modèles pour un angle ou une
% température en dehors de \SI{}{} (non fournie par les paquets chargés).
\providecommand{\degree}{^{\circ}}
% \qed (fin de démonstration), utilisé par les modèles sans amsthm
\providecommand{\qed}{\hfill$\square$}
% \barre : double barre des valeurs interdites dans un tableau de variations (tkz-tab)
\providecommand{\barre}{\ensuremath{\|}}
% environnement proof (amsthm non chargé) utilisé par les modèles
\ifdefined\proof\else\newenvironment{proof}[1][Démonstration]{\par\noindent\textit{#1.}\ }{\qed\par}\fi
\usepackage{lastpage}
\usepackage{hyperref}
\hypersetup{hidelinks}

% ============================================================
% Palette « Pop » OnBuch+ — mêmes hex que la classe P (plus_ui.dart)
% ============================================================
\definecolor{poporange}{HTML}{F59321}
\definecolor{popdark}{HTML}{E07A0C}
\definecolor{popcream}{HTML}{FFF6E8}
\definecolor{popink}{HTML}{1C1714}
\definecolor{poppurple}{HTML}{7A5AE0}
\definecolor{popgreen}{HTML}{2E9E6A}
\definecolor{poppink}{HTML}{E0457B}
\definecolor{popblue}{HTML}{2F7DE1}
\definecolor{popgold}{HTML}{C98A12}
\definecolor{popmuted}{HTML}{8A7F76}
\definecolor{popline}{HTML}{EADFD2}
\definecolor{popgray}{HTML}{8A7F76}
\colorlet{popgrey}{popgray}
% Alias tolérants (le modèle nomme parfois les couleurs différemment)
\colorlet{popwhite}{white}
\colorlet{popblack}{popink}
\colorlet{popred}{poppink}
\colorlet{popyellow}{popgold}
% Teintes claires (fonds de tableaux/encadrés) — le modèle nomme parfois
% les couleurs avec un suffixe L (ex. popblueL) sans les avoir définies.
\colorlet{poporangeL}{poporange!20}
\colorlet{popdarkL}{popdark!20}
\colorlet{poppurpleL}{poppurple!20}
\colorlet{popgreenL}{popgreen!20}
\colorlet{poppinkL}{poppink!20}
\colorlet{popblueL}{popblue!20}
\colorlet{popgoldL}{popgold!20}
\colorlet{popredL}{popred!20}
\colorlet{popgrayL}{popgray!20}
% Teintes claires pour fonds de tableaux / zones
\colorlet{popblueL}{popblue!10!white}
\colorlet{poporangeL}{poporange!14!white}
\colorlet{popgreenL}{popgreen!12!white}
\colorlet{poppurpleL}{poppurple!12!white}
\colorlet{poppinkL}{poppink!12!white}
\colorlet{popgoldL}{popgold!14!white}

\pagecolor{popcream}

% ============================================================
% Typographie
% ============================================================
\defaultfontfeatures{Ligatures=TeX}
\setmainfont{PlusJakartaSans-Medium.ttf}[Path=../../../fonts/,BoldFont=PlusJakartaSans-Bold.ttf]
% Maths en sans-serif (Fira Math) avec chiffres et lettres droites de Plus
% Jakarta, pour que les formules se fondent dans le texte.
\usepackage[math-style=ISO,bold-style=ISO]{unicode-math}
\setmathfont{FiraMath-Regular.otf}[Path=../../../fonts/]
\setmathfont{PlusJakartaSans-Medium.ttf}[Path=../../../fonts/,range={up/num,up/{Latin,latin},\mathup}]
\usepackage{icomma} % virgule décimale sans espace en mode math
% Caractères mathématiques Unicode absents des polices de texte (Plus Jakarta,
% Archivo Black) : rendus par leur équivalent mathématique, en texte comme en
% formule — sinon ils s'affichent en carré vide (ex. « Arithmétique dans ℤ »).
\usepackage{newunicodechar}
\newunicodechar{ℝ}{\ensuremath{\mathbb{R}}}
\newunicodechar{ℤ}{\ensuremath{\mathbb{Z}}}
\newunicodechar{ℂ}{\ensuremath{\mathbb{C}}}
\newunicodechar{ℕ}{\ensuremath{\mathbb{N}}}
\newunicodechar{ℚ}{\ensuremath{\mathbb{Q}}}
\newunicodechar{𝔻}{\ensuremath{\mathbb{D}}}
\newunicodechar{α}{\ensuremath{\alpha}}
\newunicodechar{β}{\ensuremath{\beta}}
\newunicodechar{γ}{\ensuremath{\gamma}}
\newunicodechar{δ}{\ensuremath{\delta}}
\newunicodechar{ε}{\ensuremath{\varepsilon}}
\newunicodechar{θ}{\ensuremath{\theta}}
\newunicodechar{λ}{\ensuremath{\lambda}}
\newunicodechar{μ}{\ensuremath{\mu}}
\newunicodechar{σ}{\ensuremath{\sigma}}
\newunicodechar{τ}{\ensuremath{\tau}}
\newunicodechar{φ}{\ensuremath{\varphi}}
\newunicodechar{ψ}{\ensuremath{\psi}}
\newunicodechar{ω}{\ensuremath{\omega}}
\newunicodechar{Σ}{\ensuremath{\Sigma}}
% majuscules produites par \MakeUppercase dans les titres (α-aminés -> Α)
\newunicodechar{Α}{\ensuremath{\alpha}}
\newunicodechar{Β}{\ensuremath{\beta}}
\newunicodechar{Γ}{\ensuremath{\Gamma}}
\newunicodechar{Δ}{\ensuremath{\Delta}}
\newunicodechar{Ε}{\ensuremath{\varepsilon}}
\newunicodechar{Θ}{\ensuremath{\Theta}}
\newunicodechar{Λ}{\ensuremath{\Lambda}}
\newunicodechar{Μ}{\ensuremath{\mu}}
\newunicodechar{Τ}{\ensuremath{\tau}}
\newunicodechar{Φ}{\ensuremath{\Phi}}
\newunicodechar{Ψ}{\ensuremath{\Psi}}
\newunicodechar{Ω}{\ensuremath{\Omega}}
\newunicodechar{↦}{\ensuremath{\mapsto}}
\newunicodechar{∈}{\ensuremath{\in}}
\newunicodechar{∉}{\ensuremath{\notin}}
\newunicodechar{∧}{\ensuremath{\wedge}}
\newunicodechar{⇔}{\ensuremath{\Leftrightarrow}}
\newunicodechar{⟺}{\ensuremath{\Longleftrightarrow}}
\newunicodechar{⇒}{\ensuremath{\Rightarrow}}
\newunicodechar{⟹}{\ensuremath{\Longrightarrow}}
\newunicodechar{∘}{\ensuremath{\circ}}
\newunicodechar{∪}{\ensuremath{\cup}}
\newunicodechar{∩}{\ensuremath{\cap}}
\newunicodechar{≡}{\ensuremath{\equiv}}
\newunicodechar{⊂}{\ensuremath{\subset}}
\newunicodechar{⊥}{\ensuremath{\perp}}
\newunicodechar{∃}{\ensuremath{\exists}}
\newunicodechar{∀}{\ensuremath{\forall}}
\newunicodechar{∛}{\ensuremath{\sqrt[3]{\,}}}
\newunicodechar{∜}{\ensuremath{\sqrt[4]{\,}}}
\newunicodechar{⃗}{\ensuremath{{}^{\to}}}
\newunicodechar{ⁿ}{\ifmmode^{n}\else\textsuperscript{\ensuremath{n}}\fi}
\newunicodechar{ˣ}{\ifmmode^{x}\else\textsuperscript{\ensuremath{x}}\fi}
\newunicodechar{ᵇ}{\ifmmode^{b}\else\textsuperscript{\ensuremath{b}}\fi}
\newunicodechar{ʳ}{\ifmmode^{r}\else\textsuperscript{\ensuremath{r}}\fi}
\newunicodechar{ᵘ}{\ifmmode^{u}\else\textsuperscript{\ensuremath{u}}\fi}
\newunicodechar{ʸ}{\ifmmode^{y}\else\textsuperscript{\ensuremath{y}}\fi}
\newunicodechar{ᵃ}{\ifmmode^{a}\else\textsuperscript{\ensuremath{a}}\fi}
\newunicodechar{ᵉ}{\ifmmode^{e}\else\textsuperscript{\ensuremath{e}}\fi}
\newunicodechar{ᵏ}{\ifmmode^{k}\else\textsuperscript{\ensuremath{k}}\fi}
\newunicodechar{ᵖ}{\ifmmode^{p}\else\textsuperscript{\ensuremath{p}}\fi}
\newunicodechar{ᴺ}{\ifmmode^{N}\else\textsuperscript{\ensuremath{N}}\fi}
\newunicodechar{ᶜ}{\ifmmode^{c}\else\textsuperscript{\ensuremath{c}}\fi}
\newunicodechar{ʰ}{\ifmmode^{h}\else\textsuperscript{\ensuremath{h}}\fi}
\newunicodechar{ᵠ}{\ifmmode^{q}\else\textsuperscript{\ensuremath{q}}\fi}
\newunicodechar{ₙ}{\ifmmode_{n}\else\textsubscript{\ensuremath{n}}\fi}
\newunicodechar{ₐ}{\ifmmode_{a}\else\textsubscript{\ensuremath{a}}\fi}
\newunicodechar{ᵢ}{\ifmmode_{i}\else\textsubscript{\ensuremath{i}}\fi}
\newunicodechar{ᵣ}{\ifmmode_{r}\else\textsubscript{\ensuremath{r}}\fi}
\newunicodechar{ₖ}{\ifmmode_{k}\else\textsubscript{\ensuremath{k}}\fi}
\newunicodechar{ᵦ}{\ifmmode_{\beta}\else\textsubscript{\ensuremath{\beta}}\fi}
\newunicodechar{ₓ}{\ifmmode_{x}\else\textsubscript{\ensuremath{x}}\fi}
\newfontfamily\popdisplay{ArchivoBlack-Regular.ttf}[Path=../../../fonts/]
\newfontfamily\poplogo{SpaceGrotesk-Bold.ttf}[Path=../../../fonts/]
\newfontfamily\popsemi{PlusJakartaSans-SemiBold.ttf}[Path=../../../fonts/]
\newfontfamily\popxbold{PlusJakartaSans-ExtraBold.ttf}[Path=../../../fonts/]
\newfontfamily\popxboldls{PlusJakartaSans-ExtraBold.ttf}[Path=../../../fonts/,LetterSpace=3.0]

\setlist[enumerate]{leftmargin=1.7em,itemsep=5pt,topsep=4pt}
\setlist[itemize]{leftmargin=1.7em,itemsep=4pt,topsep=4pt,label={\color{poporange}\textbullet}}
\setlength{\parindent}{0pt}
\setlength{\parskip}{5pt}
\renewcommand{\arraystretch}{1.25}

% pgfplots : style Pop par défaut
\pgfplotsset{
  every axis/.append style={axis line style={popink,line width=1pt},tick style={popink},
    grid style={popline,line width=0.5pt},label style={font=\small},tick label style={font=\footnotesize}},
  popcurve/.style={poporange,line width=1.8pt,no markers,smooth},
}
% Même style disponible dans un \draw TikZ simple (hors pgfplots)
\tikzset{popcurve/.style={poporange,line width=1.8pt}}
\tikzset{popgrid/.style={popline,very thin}}
\colorlet{popcurve}{poporange} % le nom du style est parfois utilisé comme couleur

% ============================================================
% Pill / squiggle
% ============================================================
\newcommand{\poppill}[3]{%
  \tikz[baseline=-0.55ex]\node[draw=popink,line width=1.2pt,rounded corners=2.6pt,%
    fill=#2,inner xsep=6.5pt,inner ysep=3pt]{%
    \popxboldls\fontsize{7}{7}\selectfont\color{#3}\MakeUppercase{#1}};%
}
\newcommand{\popsquiggle}[1]{%
  \tikz[baseline=-0.4ex]{
    \draw[#1,line width=2.6pt,line cap=round]
      plot [smooth] coordinates {(0,0.09) (0.34,0.01) (0.68,0.21) (1.02,0.01) (1.36,0.10)};
  }%
}
% Mot-clé surligné (marqueur orange sous le texte)
\newcommand{\cle}[1]{{\popxbold\color{popdark}#1}}

% ============================================================
% En-tête / pied
% ============================================================
\pagestyle{fancy}
\fancyhf{}
\renewcommand{\headrulewidth}{0pt}
\renewcommand{\footrulewidth}{0pt}
\lhead{\raisebox{-0.15ex}{\poplogo\fontsize{13}{13}\selectfont\color{popink}OnBuch\color{poporange}+}}
\rhead{\poppill{\DOCMATIERE\ \textbullet\ \DOCNIVEAU\ \textbullet\ Leçon \DOCLECON}{white}{popink}}
\lfoot{\popxbold\fontsize{7.6}{7.6}\selectfont\color{popmuted}\MakeUppercase{Cours signé OnBuch+}\ \textbullet\ {\popsemi\DOCTITRE}}
\rfoot{\tikz[baseline=-0.6ex]\node[rounded corners=8.5pt,fill=popink,minimum height=17pt,minimum width=17pt,inner xsep=5pt,inner sep=0]{\popxbold\fontsize{8}{8}\selectfont\color{popcream}\thepage\,/\,\pageref*{LastPage}};}

% ============================================================
% Titres
% ============================================================
\newcommand{\chaptitre}[2]{%
  \begin{tcolorbox}[enhanced,colback=poporange,colframe=popink,boxrule=2.6pt,arc=11pt,
      drop shadow=popink,left=15pt,right=15pt,top=12pt,bottom=11pt,before skip=3pt,after skip=18pt]
    \poppill{#2}{popink}{popcream}\par\vspace{9pt}
    {\raggedright\hyphenpenalty=10000\exhyphenpenalty=10000\popdisplay\fontsize{22}{24}\selectfont\color{popcream}\MakeUppercase{#1}\par}
  \end{tcolorbox}
}
% Section numérotée : pastille orange avec le numéro + titre Archivo
\newcounter{popsec}
\newcommand{\coursec}[1]{%
  \stepcounter{popsec}\setcounter{popsub}{0}%
  \par\vspace{16pt}\noindent
  \begin{minipage}{\linewidth}
  \tikz[baseline=-0.7ex]\node[draw=popink,line width=1.6pt,fill=poporange,rounded corners=4pt,minimum size=22pt,inner sep=0]{\popdisplay\fontsize{12}{12}\selectfont\color{popcream}\thepopsec};%
  \hspace{8pt}{\popdisplay\fontsize{15}{17}\selectfont\color{popink}\MakeUppercase{#1}}\par
  \vspace{4pt}\noindent{\color{poporange}\rule{36pt}{3.6pt}}
  \end{minipage}\par\nopagebreak\vspace{8pt}
}
\newcounter{popsub}
\newcommand{\courssub}[1]{%
  \stepcounter{popsub}%
  \par\vspace{9pt}\noindent{\popxbold\fontsize{11.5}{13}\selectfont\color{popink}{\color{poporange}\thepopsec.\thepopsub}\hspace{6pt}#1}\par\nopagebreak\vspace{4pt}
}

% ============================================================
% Boîtes pédagogiques — même recette : fond blanc, bordure épaisse colorée,
% ombre dure encre, badge plein. Titre optionnel [#1].
% ============================================================
\tcbset{popbase/.style={breakable,enhanced,colback=white,boxrule=2.3pt,arc=9pt,
  shadow={3pt}{-3pt}{0pt}{popink},left=14pt,right=14pt,top=11pt,bottom=11pt,before skip=11pt,after skip=15pt,
  fontupper=\popsemi\fontsize{10.6}{14.5}\selectfont\color{popink},
  fontlower=\popsemi\fontsize{10.6}{14.5}\selectfont\color{popink}}}
% Coche dessinée (le glyphe ✓ n'existe pas dans les polices du document)
\newcommand{\popcheck}[1]{\tikz[baseline=-0.5ex]\draw[#1,line width=1.6pt,line cap=round,line join=round](-0.13,0.0)--(-0.04,-0.09)--(0.13,0.1);}
\providecommand{\coche}{\popcheck{popgreen}}
\providecommand{\resultat}[1]{\par\smallskip\noindent\fcolorbox{popgreen}{popgreenL}{\parbox{\dimexpr\linewidth-2\fboxsep-2\fboxrule\relax}{\textbf{Résultat.} #1}}\par\smallskip}
\newcommand{\popboxhead}[3]{\poppill{#1}{#2}{white}\ifx\relax#3\relax\else\hspace{6pt}{\popxbold\fontsize{10.6}{12}\selectfont\color{popink}#3}\fi\par\vspace{7pt}}

\newtcolorbox{aretenir}[1][]{popbase,colframe=popgreen,before upper={\popboxhead{À retenir}{popgreen}{#1}}}
% Texte en allemand (document de lecture, dialogue, extrait) : cadre
% d'étude. Un titre optionnel (source, auteur) s'affiche dans l'en-tête.
\newtcolorbox{textebox}[1][]{popbase,colframe=popblue,before upper={\popboxhead{Text}{popblue}{#1}\begin{otherlanguage}{ngerman}},after upper={\end{otherlanguage}}}
% Marques ✓ / ✗ (forme correcte / fautive) souvent utilisées par les modèles
\providecommand{\cross}{\textcolor{poppink}{\ensuremath{\times}}}
\providecommand{\xmark}{\cross}\providecommand{\crossmark}{\cross}\providecommand{\cmark}{\ensuremath{\checkmark}}
% Ligne de réponse / justification à compléter (inventée par les modèles)
\providecommand{\justification}[1]{\par\noindent\textit{Justification~:} \dotfill\par}
% \textipa (package tipa non chargé) : la police ne couvre pas l'API -> simple passage
\providecommand{\textipa}[1]{#1}
% Soulignement (ulem non chargé) : \uline, \ul
\providecommand{\uline}[1]{\underline{#1}}\providecommand{\ul}[1]{\underline{#1}}
% \glossaire{\item ...} inventée par les modèles : liste de glossaire
\providecommand{\glossaire}[1]{\begin{itemize}#1\end{itemize}}
% \alert (beamer) inventée par les modèles : texte mis en évidence
\providecommand{\alert}[1]{\textbf{\textcolor{poppink}{#1}}}
% variantes ASCII d'ordinaux écrites en macro
\providecommand{\iere}{\textsuperscript{re}}\providecommand{\ieme}{\textsuperscript{e}}\providecommand{\ere}{\textsuperscript{re}}\providecommand{\eme}{\textsuperscript{e}}\providecommand{\er}{\textsuperscript{er}}\providecommand{\ie}{\textsuperscript{e}}
% Ordinaux français écrits en macro par les modèles : 1\ière, 3\ième
\newcommand{\ière}{\textsuperscript{re}}\newcommand{\ième}{\textsuperscript{e}}
% \exemple (inventée par les modèles) : étiquette « Exemple : »
\providecommand{\exemple}{\textbf{Exemple~:}\ }
% \square utilisable aussi hors mode math (cases à cocher)
\AtBeginDocument{\let\squareMath\square\renewcommand{\square}{\ensuremath{\squareMath}}}
% Mot ou phrase en allemand à l'intérieur d'un paragraphe français
\newcommand{\all}[1]{\ifmmode\text{\textit{#1}}\else\begingroup\selectlanguage{ngerman}\itshape #1\endgroup\fi}
\let\esp\all % alias toléré
\newtcolorbox{exemplebox}[1][]{popbase,colframe=poporange,before upper={\popboxhead{Exemple}{poporange}{#1}}}
\newtcolorbox{definition}[1][]{popbase,colframe=popblue,before upper={\popboxhead{Définition}{popblue}{#1}}}
\newtcolorbox{propriete}[1][]{popbase,colframe=poppurple,before upper={\popboxhead{Propriété}{poppurple}{#1}}}
\newtcolorbox{methode}[1][]{popbase,colframe=popgold,before upper={\popboxhead{Méthode}{popgold}{#1}}}
\newtcolorbox{attention}[1][]{popbase,colframe=poppink,before upper={\popboxhead{Attention}{poppink}{#1}}}
\newtcolorbox{experience}[1][]{popbase,colframe=popblue,colback=popblueL,before upper={\popboxhead{Expérience}{popblue}{#1}}}
% « Le savais-tu ? » : carte violette pleine (contexte, culture, Cameroun)
\newtcolorbox{savaistu}[1][]{popbase,colback=poppurple,colframe=popink,
  fontupper=\popsemi\fontsize{10.4}{14.2}\selectfont\color{white},
  before upper={\poppill{Le savais-tu\,?}{white}{poppurple}\ifx\relax#1\relax\else\hspace{6pt}{\popxbold\color{white}#1}\fi\par\vspace{7pt}}}
% Exercice résolu : énoncé en haut, \tcblower puis la solution sur fond crème
\newcounter{popexo}\newcounter{popexr}
\newtcolorbox{exoresolu}[1][]{popbase,colframe=poporange,colbacklower=popcream,
  segmentation style={popink,line width=1.2pt,dashed},
  before upper={\stepcounter{popexr}\popboxhead{Exercice résolu \thepopexr}{poporange}{#1}},
  before lower={\poppill{Solution}{popink}{popcream}\par\vspace{6pt}}}
% Exercice (non résolu en place) — la correction est dans la partie Corrigés
\newtcolorbox{exercice}[1][]{popbase,colframe=popink,
  before upper={\stepcounter{popexo}\popboxhead{Exercice \thepopexo}{popink}{#1}}}
% Corrigé
\newtcolorbox{corrige}[1][]{popbase,colframe=popgreen,colback=popgreenL,
  before upper={\popboxhead{Corrigé}{popgreen}{#1}}}
% Objectifs / prérequis / bilan
\newtcolorbox{objectifs}{popbase,colframe=popink,colback=poporangeL,
  before upper={\popboxhead{Objectifs de la leçon}{popink}{}}}
\newtcolorbox{prerequis}{popbase,colframe=popink,colback=popblueL,
  before upper={\popboxhead{Prérequis}{popblue}{}}}
\newtcolorbox{bilan}{popbase,colframe=popink,colback=white,boxrule=2.8pt,
  before upper={\popboxhead{Fiche bilan}{poporange}{L'essentiel à connaître pour l'examen}}}
% Figure encadrée avec légende
\newtcolorbox{popfigure}[1][]{enhanced,breakable=false,colback=white,colframe=popline,boxrule=1.4pt,arc=8pt,
  left=10pt,right=10pt,top=10pt,bottom=8pt,before skip=10pt,after skip=12pt,halign=center,
  lower separated=false,halign lower=center,
  fontlower=\popsemi\fontsize{9}{11.5}\selectfont\color{popmuted},
  #1}
\newcommand{\legende}[1]{\par\vspace{6pt}{\popsemi\fontsize{9}{11.5}\selectfont\color{popmuted}#1}}

% ============================================================
% Signature OnBuch+
% ============================================================
\newcommand{\poplogoinline}[1]{{\poplogo\fontsize{#1}{#1}\selectfont\color{popink}OnBuch\color{poporange}+}}
\newcommand{\signatureonbuch}{%
  \begin{tcolorbox}[enhanced,colback=popink,colframe=popink,boxrule=2.2pt,arc=10pt,
      drop shadow=poporange,left=14pt,right=14pt,top=10pt,bottom=10pt,before skip=0pt,after skip=0pt]
    \tikz[baseline=-0.7ex]\node[circle,fill=popgreen,draw=popcream,line width=1.3pt,minimum size=22pt,inner sep=0]{\popcheck{white}};%
    \hspace{9pt}\begin{minipage}[c]{0.62\linewidth}
      {\popxbold\fontsize{12}{13}\selectfont\color{popcream}Signé\ \textbullet\ {\poplogo OnBuch\color{poporange}+}}\par\vspace{2pt}
      {\raggedright\popsemi\fontsize{8}{10.5}\selectfont\color{popline}\MakeUppercase{Équipe pédagogique OnBuch+}\\\MakeUppercase{Conforme au programme officiel MINESEC}\par}
    \end{minipage}\hfill
    {\poplogo\fontsize{22}{22}\selectfont\color{popcream}OnBuch\color{poporange}+}
  \end{tcolorbox}%
}

% ============================================================
% Page de garde
% #1 titre · #2 matière · #3 niveau · #4 module · #5 n° leçon · #6 durée ·
% #7 description / objectifs (texte ou itemize)
% ============================================================
\newcommand{\pagegarde}[7]{%
  \thispagestyle{empty}
  \newgeometry{a4paper,margin=1.7cm,top=1.6cm,bottom=1.4cm}%
  \begin{tikzpicture}[remember picture,overlay]
    \fill[poporange!18] ([xshift=-3.2cm,yshift=-3.6cm]current page.north east) circle (4.6cm);
    \fill[poppurple!14] ([xshift=2.4cm,yshift=7.6cm]current page.south west) circle (3.8cm);
  \end{tikzpicture}%
  {\poplogo\fontsize{30}{30}\selectfont\color{popink}OnBuch\color{poporange}+}\hfill
  \raisebox{0.6ex}{\poppill{Cours signé \textbullet\ Premium}{popink}{popcream}}\par
  \vspace{1pt}\popsquiggle{poppurple}\par
  \vspace{8pt}
  \begin{tcolorbox}[enhanced,colback=poporange,colframe=popink,boxrule=2.8pt,arc=13pt,
      drop shadow=popink,left=18pt,right=18pt,top=12pt,bottom=13pt,after skip=10pt]
    \poppill{#2 \textbullet\ #3}{popink}{popcream}\hspace{5pt}\poppill{#4}{white}{popink}\par\vspace{8pt}
    {\popdisplay\fontsize{14}{14}\selectfont\color{popink}LEÇON #5}\par\vspace{4pt}
    {\raggedright\hyphenpenalty=10000\exhyphenpenalty=10000\popdisplay\fontsize{25}{27}\selectfont\color{popcream}\MakeUppercase{#1}\par}
    \vspace{8pt}
    {\popxbold\fontsize{9.5}{9.5}\selectfont\color{popink}\MakeUppercase{Durée conseillée : #6}}
  \end{tcolorbox}
  \begin{tcolorbox}[enhanced,colback=white,colframe=popink,boxrule=2.2pt,arc=10pt,
      drop shadow=popink,left=16pt,right=16pt,top=10pt,bottom=10pt,after skip=8pt]
    \poppill{Dans cette leçon}{poppurple}{white}\par\vspace{5pt}
    \popsemi\fontsize{9.3}{12.6}\selectfont\color{popink}#7
  \end{tcolorbox}
  \noindent\begin{minipage}[t]{0.487\linewidth}
    \begin{tcolorbox}[enhanced,colback=popgreen,colframe=popink,boxrule=2.2pt,arc=10pt,
        drop shadow=popink,left=11pt,right=11pt,top=10pt,bottom=10pt,height=3.1cm,valign=center]
      \tcbox[colback=white,colframe=popink,boxrule=1.3pt,arc=5pt,boxsep=0pt,nobeforeafter,
        left=3pt,right=3pt,top=3pt,bottom=3pt,valign=center]{\qrcode[height=1.9cm]{https://play.google.com/store/apps/details?id=com.luvvix.onbuch&hl=fr}}%
      \hspace{8pt}\begin{minipage}[c]{0.5\linewidth}
        {\popdisplay\fontsize{11}{12.5}\selectfont\color{white}\MakeUppercase{Télécharge OnBuch}}\par\vspace{4pt}
        {\raggedright\popsemi\fontsize{8.4}{11}\selectfont\color{white}Gratuit \textbullet\ Google Play\\Tuteur IA, annales, concours, résultats\par}
      \end{minipage}
    \end{tcolorbox}
  \end{minipage}\hfill
  \begin{minipage}[t]{0.487\linewidth}
    \begin{tcolorbox}[enhanced,colback=poppurple,colframe=popink,boxrule=2.2pt,arc=10pt,
        drop shadow=popink,left=11pt,right=11pt,top=10pt,bottom=10pt,height=3.1cm,valign=center]
      \tikz[baseline=-0.7ex]\node[circle,fill=white,draw=popink,line width=1.3pt,minimum size=1.6cm,inner sep=0]{\popdisplay\fontsize{24}{24}\selectfont\color{poppurple}+};%
      \hspace{8pt}\begin{minipage}[c]{0.6\linewidth}
        {\popdisplay\fontsize{11}{12.5}\selectfont\color{white}\MakeUppercase{Passe à OnBuch+}}\par\vspace{4pt}
        {\raggedright\popsemi\fontsize{8.4}{11}\selectfont\color{white}Tous les cours signés, Tuteur IA illimité, fiches PDF — onglet OnBuch+ de l'app\par}
      \end{minipage}
    \end{tcolorbox}
  \end{minipage}
  \vspace{8pt}
  \popcontenu
  \vfill
  \signatureonbuch
  \restoregeometry
  \newpage
}

% Rangée « Ce cours contient » (page de garde)
\newcommand{\poptile}[3]{% #1 couleur #2 gros texte #3 légende
  \begin{tcolorbox}[enhanced,colback=#1,colframe=popink,boxrule=2pt,arc=9pt,shadow={2.5pt}{-2.5pt}{0pt}{popink},
      left=6pt,right=6pt,top=9pt,bottom=9pt,halign=center,valign=center,height=1.75cm,width=\linewidth,nobeforeafter]
    {\popdisplay\fontsize{12.5}{13}\selectfont\color{white}\MakeUppercase{#2}}\par\vspace{3pt}
    {\centering\popsemi\fontsize{7.8}{9.5}\selectfont\color{white}#3\par}
  \end{tcolorbox}}
\newcommand{\popcontenu}{%
  \par\noindent{\popxboldls\fontsize{7.5}{7.5}\selectfont\color{popmuted}CE COURS SIGNÉ CONTIENT}\par\vspace{7pt}
  \noindent\begin{minipage}[t]{0.235\linewidth}\poptile{popblue}{Cours}{complet, illustré, pas à pas}\end{minipage}\hfill
  \begin{minipage}[t]{0.235\linewidth}\poptile{poporange}{Méthodes}{et exercices résolus}\end{minipage}\hfill
  \begin{minipage}[t]{0.235\linewidth}\poptile{poppink}{Exercices}{type examen + corrigés}\end{minipage}\hfill
  \begin{minipage}[t]{0.235\linewidth}\poptile{popgreen}{Fiche bilan}{l'essentiel à retenir}\end{minipage}\par}

% Sommaire
\newtcolorbox{sommaire}{enhanced,colback=white,colframe=popink,boxrule=2.2pt,arc=10pt,
  drop shadow=popink,left=18pt,right=18pt,top=13pt,bottom=13pt,before skip=6pt,after skip=20pt,
  before upper={\poppill{Sommaire}{poppurple}{white}\par\vspace{9pt}\popsemi\fontsize{10.6}{14.5}\selectfont\color{popink}}}

% Signature de fin de cours — poussée tout en bas de la dernière page
\newcommand{\findecours}{%
  \par\vfill\nopagebreak
  \begin{center}\popsquiggle{poporange}\end{center}\vspace{4pt}
  \signatureonbuch
}
\AtBeginDocument{\def\llbracket{\mathopen{[\mkern-3.5mu[}}\def\rrbracket{\mathclose{]\mkern-3.5mu]}}} % intervalles d'entiers (glyphe ⟦ absent de la police)
% Glyphes absents de Fira Math : redéfinis à partir de glyphes disponibles
\AtBeginDocument{%
  \def\setminus{\mathbin{\backslash}}\let\smallsetminus\setminus
  \def\vdots{\mathord{\vbox{\baselineskip3.2pt\lineskiplimit0pt\kern4pt\hbox{.}\hbox{.}\hbox{.}}}}%
  \def\ddots{\mathinner{\mkern1mu\raise6.5pt\hbox{.}\mkern2mu\raise3.6pt\hbox{.}\mkern2mu\raise0.7pt\hbox{.}\mkern1mu}}%
  \def\triangle{\mathord{\scalebox{0.9}{$\symup{\Delta}$}}}%
  \def\nabla{\mathord{\raisebox{\depth}{\scalebox{1}[-1]{$\symup{\Delta}$}}}}%
  \def\checkmark{\popcheck{popdark}}%
  \def\star{\mathbin{\ast}}\def\bigstar{\mathord{\ast}}%
  \def\diamond{\mathbin{\scalebox{0.6}{$\Diamond$}}}%
  \def\bigcup{\mathop{\mathchoice{\raisebox{-0.3em}{\scalebox{1.7}{$\cup$}}}{\scalebox{1.25}{$\cup$}}{\cup}{\cup}}\displaylimits}%
  \def\bigcap{\mathop{\mathchoice{\raisebox{-0.3em}{\scalebox{1.7}{$\cap$}}}{\scalebox{1.25}{$\cap$}}{\cap}{\cap}}\displaylimits}%
  \def\lesseqqgtr{\mathrel{\lessgtr}}\def\bigoplus{\mathop{\oplus}\displaylimits}%
}
\providecommand{\ssi}{si et seulement si} % abréviation fréquente
\AtBeginDocument{\providecommand{\wideparen}{\overparen}} % arc de cercle

\begin{document}

\pagegarde{Les cours à distance : les problèmes liés à l'utilisation du numérique}{Allemand}{Tle A}{Chapitre 4 : Média et communication}{AL20}{2 h}{Cette leçon propose l'étude d'un texte sur les cours à distance (Fernunterricht) et les problèmes liés au numérique. Elle permet d'acquérir le vocabulaire thématique, de maîtriser l'expression de l'opposition (einerseits … andererseits) ainsi que les propositions temporelles au plus-que-parfait, et de produire un commentaire et un court texte argumenté.
\vspace{4pt}
\begin{multicols}{2}\small
\begin{itemize}
\item À la fin de cette leçon, je sais comprendre globalement puis en détail un texte allemand sur les cours à distance et les problèmes du numérique
\item À la fin de cette leçon, je sais employer correctement le vocabulaire thématique (der Fernunterricht, die Verbindung, der Strom, das Netz, der Vorteil, der Nachteil) avec articles et pluriels
\item À la fin de cette leçon, je sais exprimer l'opposition avec einerseits … andererseits dans des phrases correctes
\item À la fin de cette leçon, je sais former et employer le plus-que-parfait (Plusquamperfekt) avec haben ou sein
\item À la fin de cette leçon, je sais construire des propositions temporelles avec nachdem, als, bevor, während au bon temps
\item À la fin de cette leçon, je sais répondre à des questions de compréhension en phrases complètes
\end{itemize}\end{multicols}}

\begin{sommaire}
\begin{multicols}{2}
\begin{enumerate}
\item Texte support : Fernunterricht – Chancen und Probleme
\item Compréhension du texte
\item Lexique thématique : der Fernunterricht und das Digitale
\item Grammaire 1 : l'opposition avec einerseits … andererseits
\item Grammaire 2 : le plus-que-parfait et les propositions temporelles
\item Méthode du commentaire et production guidée
\item Activité d'intégration
\item Méthodes et exercices résolus
\item Exercices
\item Corrigés des exercices
\item Fiche bilan
\end{enumerate}
\end{multicols}
\end{sommaire}

\begin{objectifs}
\begin{itemize}
\item À la fin de cette leçon, je sais comprendre globalement puis en détail un texte allemand sur les cours à distance et les problèmes du numérique
\item À la fin de cette leçon, je sais employer correctement le vocabulaire thématique (der Fernunterricht, die Verbindung, der Strom, das Netz, der Vorteil, der Nachteil) avec articles et pluriels
\item À la fin de cette leçon, je sais exprimer l'opposition avec einerseits … andererseits dans des phrases correctes
\item À la fin de cette leçon, je sais former et employer le plus-que-parfait (Plusquamperfekt) avec haben ou sein
\item À la fin de cette leçon, je sais construire des propositions temporelles avec nachdem, als, bevor, während au bon temps
\item À la fin de cette leçon, je sais répondre à des questions de compréhension en phrases complètes
\item À la fin de cette leçon, je sais rédiger un résumé et un court paragraphe argumenté sur les avantages et les inconvénients du numérique
\item À la fin de cette leçon, je sais commenter un texte selon la méthode du commentaire de texte au Bac
\end{itemize}
\end{objectifs}

\begin{prerequis}
\begin{itemize}
\item Le parfait (Perfekt) avec haben et sein (classe de Première)
\item Les propositions subordonnées et la place du verbe conjugué en fin de proposition
\item Les conjonctions temporelles de base (als, wenn, während)
\item Le vocabulaire général de l'école et des médias (die Schule, der Computer, das Internet)
\item L'expression de l'opinion (ich denke, meiner Meinung nach)
\end{itemize}
\end{prerequis}

\newpage

\coursec{Texte support : Fernunterricht – Chancen und Probleme}

\courssub{Situation d'accroche et problématique}

Pendant la pandémie de Covid-19, des millions d'élèves au Cameroun, en Allemagne, en Autriche et en Suisse ont dû suivre leurs cours à distance, devant un écran. À Yaoundé, on se connectait depuis le salon familial ou depuis un cybercafé du quartier ; à Berlin ou à Vienne, on s'installait devant un ordinateur portable dans sa chambre. Mais à Yaoundé comme à Berlin, les problèmes étaient parfois les mêmes : coupures d'électricité, connexion internet instable, isolement, difficulté à rester concentré. \all{Einerseits bietet der Fernunterricht viele Vorteile, andererseits gibt es auch große Probleme.} « D'un côté, l'enseignement à distance offre de nombreux avantages ; de l'autre, il pose aussi de grands problèmes. »

\textbf{Problématique :} comment comprendre un texte allemand qui raconte une expérience des cours à distance, et comment parler soi-même de ces avantages et de ces inconvénients en allemand ? C'est l'objet de cette leçon. Nous allons d'abord lire le texte globalement, puis le comprendre en détail, relever son vocabulaire et observer deux faits de grammaire que nous étudierons ensuite : l'opposition avec \all{einerseits \ldots{} andererseits} et le plus-que-parfait dans les propositions temporelles.

\courssub{Le texte : lecture globale}

\begin{attention}[Méthode de lecture : ne pas traduire mot à mot]
Face à un texte allemand, le réflexe du francophone est de traduire chaque mot dans l'ordre. C'est une erreur : l'ordre des mots allemand (verbe conjugué en deuxième position dans la phrase principale, verbe conjugué à la fin dans les subordonnées) rend la traduction mot à mot impossible. La bonne démarche :
\begin{enumerate}
\item Lire le texte entier une première fois \emph{sans dictionnaire} pour saisir le sens global (qui parle ? de quoi ? quel est le bilan ?).
\item Repérer les mots déjà connus et les mots transparents (\all{Internet, Computer, Transport, Probleme}).
\item Relire en vérifiant les détails avec le glossaire.
\item Seulement ensuite, répondre aux questions et, si nécessaire, traduire phrase par phrase en respectant la structure allemande.
\end{enumerate}
\end{attention}

\begin{textebox}[Fernunterricht – eine Erfahrung]
Während der Corona-Pandemie mussten viele Schüler in Kamerun und in Deutschland zu Hause bleiben. Nachdem die Schulen geschlossen hatten, begann der Unterricht per Internet, per Radio oder per Fernsehen.

Ich erinnere mich noch gut an die erste Woche. Mein Vater hatte mir ein altes Smartphone gegeben, und ich hatte eine Internetverbindung gekauft. Einerseits war das praktisch: Man konnte zu Hause lernen, und man sparte Zeit und Geld für den Transport. Ich musste nicht mehr um sechs Uhr aufstehen und mit dem Taxi oder dem Motorrad zur Schule fahren. Auch die Lehrer schickten uns Übungen per WhatsApp, und wir konnten die Lektionen wiederholen, so oft wir wollten.

Andererseits gab es viele Probleme. Die Verbindung war oft schlecht, manchmal brach sie mitten in der Stunde ab. Der Strom fiel aus, oft mehrere Stunden am Tag, und dann war das Telefon leer. Viele meiner Klassenkameraden hatten keinen Computer und kein Smartphone; sie konnten den Kursen gar nicht folgen. Nachdem ich drei Monate online gelernt hatte, merkte ich, dass mir der Kontakt zu meinen Freunden und Lehrern fehlte. Niemand konnte mir sofort helfen, wenn ich eine Frage hatte, und ich verstand manche Erklärungen nicht.

Am Ende war ich froh, als die Schulen wieder öffneten. Der Fernunterricht ist also eine Chance, besonders in schwierigen Zeiten, aber er kann die Schule und den Lehrer nicht ersetzen.
\end{textebox}

\textbf{Glossaire (allemand -- français)}

\begin{center}
\begin{tabularx}{\linewidth}{|X|X|}
\hline
\rowcolor{popblueL} \textbf{Deutsch} & \textbf{Français} \\
\hline
die Pandemie, -n & la pandémie \\
\hline
zu Hause bleiben (ist geblieben) & rester à la maison \\
\hline
schließen (hat geschlossen) & fermer \\
\hline
per Internet / per Radio & par internet / à la radio \\
\hline
sich erinnern an + Akk. & se souvenir de \\
\hline
die Internetverbindung, -en & la connexion internet \\
\hline
praktisch & pratique, commode \\
\hline
sparen & économiser \\
\hline
aufstehen (ist aufgestanden) & se lever \\
\hline
die Übung, -en & l'exercice \\
\hline
schicken & envoyer \\
\hline
wiederholen & répéter, réviser \\
\hline
abbrechen (bricht ab, brach ab, hat abgebrochen) & s'interrompre, couper \\
\hline
der Strom (Sg.) & l'électricité, le courant \\
\hline
ausfallen (fällt aus, fiel aus, ist ausgefallen) & tomber en panne, être coupé \\
\hline
leer (hier : das Telefon ist leer) & vide (ici : déchargé) \\
\hline
folgen + Dat. (ist gefolgt) & suivre (un cours) \\
\hline
merken & remarquer, se rendre compte \\
\hline
fehlen + Dat. & manquer à (\all{mir fehlt der Kontakt}) \\
\hline
die Erklärung, -en & l'explication \\
\hline
wieder öffnen & rouvrir \\
\hline
ersetzen & remplacer \\
\hline
\end{tabularx}
\end{center}

\begin{definition}[der Fernunterricht]
\all{der Fernunterricht} (masculin, en général sans pluriel) désigne l'enseignement à distance : des cours suivis par internet, par la radio ou par la télévision, sans présence physique en classe. Le mot est composé de \all{fern} (« loin, à distance », comme dans \all{der Fernseher}, le téléviseur, littéralement « le regarde-loin ») et de \all{der Unterricht} (« l'enseignement, le cours »). Prononciation : « fèr-oun-tèr-richte ».
\end{definition}

\begin{savaistu}[Homeschooling, Distanzunterricht, Distance Learning]
En Allemagne, on utilise aussi le mot anglais \all{das Homeschooling} et le terme \all{der Distanzunterricht}. En Autriche, le terme officiel pendant la pandémie était \all{Distance Learning}. En Suisse, on parlait souvent de \all{Fernlernen}. Ces mots montrent que la langue allemande emprunte volontiers à l'anglais dans le domaine du numérique : \all{das Internet, der Computer, das Smartphone, online, der Download}. Attention cependant : ces mots anglais prennent un article allemand et une majuscule !
\end{savaistu}

\courssub{Repérage du thème, du narrateur, du ton et de la structure}

Après une première lecture, répondons aux questions fondamentales de la compréhension globale.

\begin{center}
\begin{tabularx}{\linewidth}{|X|X|}
\hline
\rowcolor{popblueL} \textbf{Question} & \textbf{Réponse} \\
\hline
De quoi parle le texte ? (thème) & Des cours à distance (\all{der Fernunterricht}) pendant la pandémie de Covid-19, avec leurs avantages et leurs problèmes. \\
\hline
Qui parle ? (narrateur) & Un élève camerounais qui raconte sa propre expérience : narration à la première personne (\all{ich}). \\
\hline
Quel est le ton ? & Personnel et témoigné, mais réfléchi : le narrateur ne se plaint pas seulement, il pèse le pour et le contre. \\
\hline
Où et quand ? & Au Cameroun (et parallèlement en Allemagne), pendant la pandémie ; le récit est au passé. \\
\hline
Quelle est la thèse finale ? & \all{Der Fernunterricht ist eine Chance, aber er kann die Schule nicht ersetzen.} Le bilan est nuancé : une chance, mais pas un remplacement de l'école. \\
\hline
\end{tabularx}
\end{center}

\textbf{La structure du texte} suit un plan en quatre mouvements, typique du texte argumentatif à base de témoignage :

\begin{enumerate}
\item \textbf{Situation} (paragraphe 1) : la pandémie oblige les élèves à rester à la maison ; les cours commencent par internet, la radio ou la télévision.
\item \textbf{Avantages} (paragraphe 2) : apprendre à la maison, économiser le temps et l'argent du transport, répéter les leçons.
\item \textbf{Problèmes} (paragraphe 3) : connexion instable, coupures d'électricité, manque de matériel, isolement.
\item \textbf{Bilan} (paragraphe 4) : le Fernunterricht est une chance, mais il ne peut pas remplacer l'école.
\end{enumerate}

Ce plan \emph{situation $\rightarrow$ avantages $\rightarrow$ problèmes $\rightarrow$ bilan} est signalé dans le texte par le connecteur d'opposition \all{einerseits \ldots{} andererseits} : c'est lui qui organise l'argumentation.

\courssub{Lecture détaillée, paragraphe par paragraphe}

\textbf{Paragraphe 1 : la situation.} \all{Während der Corona-Pandemie mussten viele Schüler \ldots{} zu Hause bleiben.} « Pendant la pandémie de Corona, de nombreux élèves ont dû rester à la maison. » Notez le modal \all{mussten} (prétérit de \all{müssen}, « devoir ») qui exprime la contrainte. La subordonnée \all{Nachdem die Schulen geschlossen hatten} (« après que les écoles avaient fermé ») place le verbe conjugué \all{hatten} à la fin : c'est un plus-que-parfait que nous observerons plus bas. Reformulation : la pandémie a fermé les écoles, et l'école est entrée dans les maisons par les écrans et les ondes.

\textbf{Paragraphe 2 : les avantages.} Le narrateur donne d'abord un détail concret et touchant : \all{Mein Vater hatte mir ein altes Smartphone gegeben} (« mon père m'avait donné un vieux smartphone »). Puis vient la liste des avantages, introduite par \all{Einerseits} : apprendre à la maison, économiser le temps et l'argent du transport, ne plus se lever à six heures, recevoir les exercices par WhatsApp, répéter les leçons autant de fois qu'on veut. Notez la tournure \all{so oft wir wollten} (« autant de fois que nous le voulions »), subordonnée avec le verbe à la fin. Reformulation : le cours à distance rend l'élève plus autonome et allège les contraintes matérielles du quotidien.

\textbf{Paragraphe 3 : les problèmes.} \all{Andererseits gab es viele Probleme.} « D'un autre côté, il y avait beaucoup de problèmes. » Trois types de difficultés : (a) techniques -- \all{die Verbindung brach ab}, la connexion coupait en plein cours ; \all{der Strom fiel aus}, l'électricité était coupée ; (b) matérielles -- beaucoup de camarades n'avaient ni ordinateur ni smartphone ; (c) humaines -- \all{mir fehlte der Kontakt zu meinen Freunden und Lehrern}, le contact avec les amis et les professeurs manquait au narrateur. Notez la construction impersonnelle \all{es gab} (« il y avait », prétérit de \all{es gibt}) et le datif dans \all{mir fehlte} (« il me manquait »). Reformulation : la technique est fragile et l'écran ne remplace pas la présence humaine.

\textbf{Paragraphe 4 : le bilan.} \all{Am Ende war ich froh, als die Schulen wieder öffneten.} « À la fin, j'étais content quand les écoles ont rouvert. » La phrase finale formule la thèse : \all{Der Fernunterricht ist also eine Chance, aber er kann die Schule und den Lehrer nicht ersetzen.} Le connecteur \all{also} marque la conclusion ; \all{aber} introduit la nuance. Reformulation : l'enseignement à distance est un recours précieux en temps de crise, mais l'école reste irremplaçable.

\courssub{Observation guidée : les faits de grammaire dans le texte}

Avant d'étudier les règles dans les sections de grammaire, observons ce que le texte nous montre.

\textbf{1. L'opposition avec \all{einerseits \ldots{} andererseits}.} Relevez dans le texte :

\all{Einerseits war das praktisch: Man konnte zu Hause lernen \ldots{}} « D'un côté, c'était pratique \ldots{} »

\all{Andererseits gab es viele Probleme.} « D'un autre côté, il y avait beaucoup de problèmes. »

Observez : les deux éléments encadrent les deux faces d'un même sujet ; \all{einerseits} et \all{andererseits} occupent la première position de la phrase, et le verbe conjugué vient juste après (\all{war}, \all{gab}). C'est la règle de la place du verbe en deuxième position. Nous verrons dans la section de grammaire que ces connecteurs peuvent aussi se placer ailleurs dans la phrase.

\textbf{2. Le plus-que-parfait.} Relevez :

\all{Nachdem die Schulen geschlossen hatten, begann der Unterricht per Internet.} « Après que les écoles avaient fermé, les cours ont commencé par internet. »

\all{Mein Vater hatte mir ein altes Smartphone gegeben.} « Mon père m'avait donné un vieux smartphone. »

\all{Nachdem ich drei Monate online gelernt hatte, merkte ich \ldots{}} « Après avoir appris en ligne pendant trois mois, j'ai remarqué \ldots{} »

Observez : le plus-que-parfait allemand se forme avec l'auxiliaire \all{haben} (ou \all{sein}) au \emph{prétérit} + le participe passé : \all{hatte geschlossen, hatte gegeben, hatte gelernt}. Il exprime une action \emph{antérieure} à une autre action passée, exactement comme le plus-que-parfait français, et il est très souvent introduit par la conjonction temporelle \all{nachdem} (« après que »). Nous y reviendrons en détail.

\begin{exoresolu}[Compréhension globale et détaillée]
Répondez en allemand, par des phrases complètes, aux questions suivantes :

1. Warum mussten die Schüler zu Hause bleiben?

2. Nennen Sie zwei Vorteile des Fernunterrichts.

3. Welche technischen Probleme gab es?

4. Was fehlte dem Erzähler nach drei Monaten?

5. Was ist die Meinung des Erzählers am Ende?

\tcblower

\textbf{Solutions détaillées :}

1. \all{Sie mussten zu Hause bleiben, weil die Schulen während der Corona-Pandemie geschlossen waren.} « Ils ont dû rester à la maison parce que les écoles étaient fermées pendant la pandémie. » (Réponse avec \all{weil} + verbe conjugué à la fin.)

2. \all{Man konnte zu Hause lernen und Zeit und Geld für den Transport sparen. Außerdem konnte man die Lektionen so oft wiederholen, wie man wollte.} « On pouvait apprendre à la maison et économiser le temps et l'argent du transport. De plus, on pouvait répéter les leçons autant de fois qu'on voulait. »

3. \all{Die Verbindung war oft schlecht und brach manchmal ab, und der Strom fiel oft aus.} « La connexion était souvent mauvaise et coupait parfois, et l'électricité tombait souvent en panne. »

4. \all{Ihm fehlte der Kontakt zu seinen Freunden und Lehrern.} « Le contact avec ses amis et ses professeurs lui manquait. » (Attention au datif : \all{ihm fehlte}, littéralement « à lui manquait ».)

5. \all{Er meint, dass der Fernunterricht eine Chance ist, aber dass er die Schule und den Lehrer nicht ersetzen kann.} « Il pense que l'enseignement à distance est une chance, mais qu'il ne peut pas remplacer l'école et le professeur. »
\end{exoresolu}

\courssub{Carte mentale du texte}

Pour mémoriser la structure argumentative du texte et le vocabulaire du thème, voici une carte mentale que vous pouvez recopier et compléter dans votre cahier.

\begin{popfigure}
\begin{center}
\begin{tikzpicture}[mindmap, grow cyclic,
  every node/.style={concept, align=center, text=white},
  root concept/.append style={concept color=popblue, font=\large\bfseries},
  level 1/.append style={level distance=3.6cm, sibling angle=90},
  level 2/.append style={level distance=2.6cm, sibling angle=45, font=\small}]
\node[root concept]{Fern\-unterricht}
  child[concept color=popgreen] { node {Vorteile}
    child { node[concept color=popgreen!80]{zu Hause lernen} }
    child { node[concept color=popgreen!80]{Zeit und Geld sparen} }
    child { node[concept color=popgreen!80]{Lektionen wiederholen} }
  }
  child[concept color=poporange] { node {technische Probleme}
    child { node[concept color=poporange!80]{schlechte Verbindung} }
    child { node[concept color=poporange!80]{Strom fällt aus} }
    child { node[concept color=poporange!80]{kein Computer} }
  }
  child[concept color=poppurple] { node {soziale Probleme}
    child { node[concept color=poppurple!80]{Freunde fehlen} }
    child { node[concept color=poppurple!80]{Lehrer fehlen} }
    child { node[concept color=poppurple!80]{Isolation} }
  }
  child[concept color=poppink] { node {Bilanz}
    child { node[concept color=poppink!80]{eine Chance} }
    child { node[concept color=poppink!80]{ersetzt die Schule nicht} }
  };
\end{tikzpicture}
\end{center}
\legende{Carte mentale : le Fernunterricht, ses avantages, ses problèmes et le bilan}
\end{popfigure}

\courssub{Premier bilan lexical : les mots-clés du texte}

Avant la section de lexique détaillée, fixons les six mots-clés du programme, tels qu'ils apparaissent dans le texte. Apprenez chaque nom \emph{avec son article et son pluriel}.

\begin{center}
\begin{tabularx}{\linewidth}{|X|X|X|}
\hline
\rowcolor{popblueL} \textbf{Deutsch} & \textbf{Français} & \textbf{Exemple} \\
\hline
der Fernunterricht (Sg.) & l'enseignement à distance & \all{Der Fernunterricht ist eine Chance.} \\
\hline
die Verbindung, -en & la connexion, la liaison & \all{Die Verbindung war oft schlecht.} \\
\hline
der Strom (Sg.) & l'électricité, le courant & \all{Der Strom fiel aus.} \\
\hline
das Netz, -e & le réseau, le filet & \all{Das Netz funktioniert nicht.} \\
\hline
der Vorteil, -e & l'avantage & \all{Ein Vorteil ist: Man spart Zeit.} \\
\hline
der Nachteil, -e & l'inconvénient & \all{Ein Nachteil ist die Isolation.} \\
\hline
\end{tabularx}
\end{center}

\begin{attention}[Pièges fréquents des francophones]
\begin{itemize}
\item \all{der Vorteil} / \all{der Nachteil} se construisent comme en français avec « de » : \all{der Vorteil des Fernunterrichts} « l'avantage de l'enseignement à distance ». Dans un texte soigné, préférez le génitif à la tournure avec \all{von}.
\item \all{der Strom} ne signifie pas ici « le courant d'eau » mais « l'électricité » : \all{Der Strom fiel aus.} « Il y a eu une coupure d'électricité. »
\item \all{das Netz} est neutre (\all{das}), alors que « le réseau » est masculin en français. Méfiez-vous des genres : ils ne coïncident presque jamais entre les deux langues.
\item N'oubliez pas la majuscule à tous les noms : \emph{die verbindung} est fautif, il faut \all{die Verbindung}.
\item \all{die Verbindung} vient du verbe \all{verbinden} (« relier ») ; la famille de mots vous aide à mémoriser : \all{verbinden -- die Verbindung -- verbunden}.
\end{itemize}
\end{attention}

\begin{experience}[À vous de parler : votre expérience]
Par groupes de trois, répondez oralement à la question : \all{Wie war Ihre Erfahrung mit dem Fernunterricht?} « Quelle a été votre expérience des cours à distance ? » Utilisez obligatoirement dans vos réponses :
\begin{enumerate}
\item le connecteur \all{einerseits \ldots{} andererseits} ;
\item une phrase au plus-que-parfait avec \all{nachdem} (modèle : \all{Nachdem ich eine Woche online gelernt hatte, \ldots{}}) ;
\item au moins trois mots-clés du tableau ci-dessus.
\end{enumerate}
Un rapporteur de chaque groupe présente ensuite deux phrases au tableau ; la classe corrige collectivement la place du verbe et les articles.
\end{experience}

Nous avons maintenant compris le texte globalement et en détail, repéré son vocabulaire et observé les deux faits de grammaire qu'il contient. La section suivante approfondira la compréhension par des questions ciblées, avant que nous n'étudiions le lexique du numérique, puis l'opposition avec \all{einerseits \ldots{} andererseits} et le plus-que-parfait dans les propositions temporelles.

\coursec{Compréhension du texte}

Dans la section précédente, nous avons lu et découvert le texte \all{Fernunterricht – Chancen und Probleme}, ce récit dans lequel un élève camerounais raconte son expérience des cours à distance. Nous avons repéré le thème général et le vocabulaire nouveau. Il s'agit maintenant de passer à l'étape décisive de toute leçon de texte : la \cle{compréhension}, c'est-à-dire la lecture active, globale puis détaillée, qui permet de répondre à des questions, de vérifier des affirmations et de résumer le texte. C'est exactement la compétence exigée à l'épreuve d'allemand du Baccalauréat A.

Pour travailler cette section, relisons d'abord le texte support dans son intégralité.

\begin{textebox}[Fernunterricht – Chancen und Probleme]
Während der Coronazeit hatten viele Schüler in Kamerun Fernunterricht. Auch ich musste zu Hause bleiben und am Computer lernen. Einerseits war das praktisch : Ich musste nicht früh aufstehen und kein Geld für den Transport ausgeben. Andererseits gab es viele Probleme.

Das größte Problem war die Verbindung. Der Strom fiel oft aus, und das Netz war in unserem Viertel sehr schlecht. Manchmal hatte ich die Seite schon geöffnet, als der Strom wieder ausfiel. Dann musste ich warten, bis der Unterricht vorbei war.

Ein anderes Problem war sozial : Ich hatte meine Freunde lange nicht gesehen, bevor der Fernunterricht begann, und ich vermisste sie sehr. Nachdem ich drei Wochen allein gelernt hatte, fühlte ich mich einsam. Der Computer kann die Schule nicht ersetzen.

Trotzdem habe ich viel gelernt. Bevor der Kurs endete, hatte ich alle Aufgaben geschickt. Meine Lehrerin war zufrieden. Fernunterricht ist nützlich, aber er ersetzt nicht den Kontakt mit den Lehrern und den Freunden.
\end{textebox}

\textbf{Glossaire :} \all{der Fernunterricht} (les cours à distance, sans pluriel) ; \all{die Verbindung, -en} (la connexion) ; \all{der Strom} (le courant électrique, sans pluriel usuel) ; \all{ausfallen (fällt aus, fiel aus, ist ausgefallen)} (tomber en panne, être coupé) ; \all{das Netz, -e} (le réseau) ; \all{das Viertel, -} (le quartier) ; \all{vermissen} (regretter l'absence de quelqu'un) ; \all{einsam} (seul, isolé) ; \all{ersetzen} (remplacer) ; \all{die Aufgabe, -n} (le devoir, la tâche) ; \all{schicken} (envoyer) ; \all{nützlich} (utile).

\courssub{La méthode pour répondre à une question de compréhension}

Avant toute question, il faut adopter une démarche rigoureuse. Beaucoup d'élèves perdent des points non pas parce qu'ils n'ont pas compris le texte, mais parce qu'ils répondent mal : un mot isolé, une réponse en français, une phrase sans verbe conjugué. Voici la méthode en quatre étapes.

\begin{methode}[Répondre à une question de compréhension en quatre étapes]
\begin{enumerate}
\item \textbf{Lire la question attentivement} et souligner le mot interrogatif : \all{warum} (pourquoi), \all{was} (quoi), \all{wer} (qui), \all{wann} (quand), \all{wie} (comment), \all{wo} (où).
\item \textbf{Repérer le passage du texte} qui contient la réponse : chercher les mots-clés de la question dans le texte (souvent des synonymes ou des mots de la même famille).
\item \textbf{Reformuler ou citer} : recopier la phrase du texte ou la reformuler avec ses propres mots, en reprenant les termes de la question.
\item \textbf{Rédiger une phrase complète en allemand} : sujet + verbe conjugué à la deuxième place + compléments. Jamais un mot seul, jamais du français.
\end{enumerate}
\end{methode}

\begin{popfigure}
\begin{tikzpicture}[node distance=0.55cm, every node/.style={font=\small}]
\node[draw=popblue, fill=popblueL, rounded corners, align=center, text width=4.2cm] (a) {1. Lire la question\\ \emph{mot interrogatif}};
\node[draw=popgreen, fill=popgreenL, rounded corners, align=center, text width=4.2cm, below=of a] (b) {2. Repérer le passage\\ \emph{mots-clés dans le texte}};
\node[draw=poporange, fill=poporangeL, rounded corners, align=center, text width=4.2cm, below=of b] (c) {3. Reformuler ou citer\\ \emph{reprendre les termes}};
\node[draw=poppurple, fill=poppurpleL, rounded corners, align=center, text width=4.2cm, below=of c] (d) {4. Phrase complète\\ \emph{sujet + verbe + compléments}};
\draw[-{Stealth}, thick, popblue] (a) -- (b);
\draw[-{Stealth}, thick, popgreen] (b) -- (c);
\draw[-{Stealth}, thick, poporange] (c) -- (d);
\end{tikzpicture}
\legende{Organigramme de la méthode de réponse à une question de compréhension}
\end{popfigure}

\begin{exemplebox}[Un exemple de réponse modèle]
Question : \all{Warum war die Verbindung ein Problem ?} (Pourquoi la connexion était-elle un problème ?)

Réponse correcte : \all{Die Verbindung war ein Problem, weil der Strom oft ausfiel und das Netz schlecht war.}

« La connexion était un problème parce que le courant était souvent coupé et que le réseau était mauvais. »

Observez : la réponse reprend les termes de la question (\all{die Verbindung war ein Problem}), puis donne la cause avec \all{weil} et le verbe conjugué à la fin de la subordonnée.
\end{exemplebox}

\begin{attention}[Réponse juste mais mal formulée = points perdus]
Au Baccalauréat, une réponse juste mais rédigée \textbf{en français} ou en \textbf{mot isolé} perd des points, parfois la totalité des points de la question. Écrire \all{Strom} seul, ou « parce que le courant était coupé », ne suffit pas. Il faut toujours une \cle{phrase complète en allemand} avec un verbe conjugué. Autre piège : ne recopiez pas tout le paragraphe ; sélectionnez uniquement le passage qui répond à la question.
\end{attention}

\courssub{Compréhension globale du texte}

La compréhension globale vise à saisir le texte dans son ensemble : le thème, le locuteur, le point de vue de l'auteur. Ce sont souvent les premières questions du sujet.

\begin{exoresolu}[Questions de compréhension globale]
\textbf{1. Worum geht es in diesem Text ?} (De quoi parle ce texte ?)

\textbf{2. Wer spricht im Text ?} (Qui parle dans le texte ?)

\textbf{3. Was ist die Meinung des Schreibers ?} (Quelle est la conclusion de l'auteur ?)
\tcblower
\textbf{1.} \all{Der Text handelt vom Fernunterricht in Kamerun und von seinen Chancen und Problemen.} « Le texte parle des cours à distance au Cameroun, de leurs avantages et de leurs problèmes. » — On reprend le titre du texte, qui annonce toujours le thème. Remarquez la construction : \all{von etwas handeln} + datif (\all{vom Fernunterricht}, \all{von seinen Chancen}).

\textbf{2.} \all{Ein Schüler spricht im Text. Er erzählt seine Erfahrung mit dem Fernunterricht.} « C'est un élève qui parle. Il raconte son expérience des cours à distance. » — Indices : \all{ich}, \all{meine Lehrerin}, \all{ich musste zu Hause bleiben}.

\textbf{3.} \all{Der Schreiber findet : Fernunterricht ist nützlich, aber er ersetzt nicht den Kontakt mit den Lehrern und den Freunden.} « L'auteur conclut que les cours à distance sont utiles, mais qu'ils ne remplacent pas le contact avec les professeurs et les amis. » — La conclusion se trouve presque toujours dans la dernière phrase du texte.
\end{exoresolu}

Remarquez la stratégie : pour la question 1, on exploite le \textbf{titre} ; pour la question 2, on relève les \textbf{pronoms personnels} ; pour la question 3, on lit le \textbf{dernier paragraphe}. Ce sont trois réflexes de lecteur efficace.

\courssub{Compréhension détaillée du texte}

La compréhension détaillée exige de localiser des informations précises : les avantages cités, les problèmes techniques, le problème social. Chaque réponse doit être appuyée sur un passage précis du texte.

\begin{exoresolu}[Questions de compréhension détaillée]
\textbf{1. Welche Vorteile hatte der Fernunterricht für den Schüler ?} (Quels avantages les cours à distance avaient-ils pour l'élève ?)

\textbf{2. Welche technischen Probleme gab es ?} (Quels problèmes techniques y avait-il ?)

\textbf{3. Was fehlte dem Schüler sozial ?} (Que manquait-il à l'élève sur le plan social ?)
\tcblower
\textbf{1.} \all{Der Fernunterricht war praktisch : Der Schüler musste nicht früh aufstehen und kein Geld für den Transport ausgeben.} « Les cours à distance étaient pratiques : l'élève ne devait pas se lever tôt ni dépenser d'argent pour le transport. » — Passage : premier paragraphe, après \all{einerseits}.

\textbf{2.} \all{Der Strom fiel oft aus, und das Netz war sehr schlecht.} « Le courant tombait souvent en panne et le réseau était très mauvais. » — Passage : deuxième paragraphe.

\textbf{3.} \all{Dem Schüler fehlten seine Freunde. Er hatte sie lange nicht gesehen und vermisste sie sehr.} « Ses amis manquaient à l'élève. Il ne les avait pas vus depuis longtemps et ils lui manquaient beaucoup. » — Passage : troisième paragraphe. Attention : \all{jemandem fehlen} signifie « manquer à quelqu'un » (la personne qui éprouve le manque est au \textbf{datif} : \all{dem Schüler}) ; le verbe \all{vermissen} exprime le même sentiment, mais avec un sujet nominatif et un complément à l'accusatif : \all{Ich vermisse meine Freunde.}
\end{exoresolu}

\begin{exemplebox}[Deux questions types avec réponses modèles]
\all{Warum war der Fernunterricht praktisch ?} → \all{Er war praktisch, weil der Schüler nicht früh aufstehen musste und kein Geld für den Transport ausgeben musste.} « Il était pratique parce que l'élève ne devait pas se lever tôt ni payer le transport. »

\all{Was fehlte dem Schüler ?} → \all{Dem Schüler fehlte der Kontakt mit seinen Freunden.} « Le contact avec ses amis manquait à l'élève. »
\end{exemplebox}

\courssub{Richtig oder falsch ? Justifier par le texte}

L'exercice \all{Richtig oder falsch ?} (vrai ou faux) est classique au Baccalauréat. La consigne complète est souvent : \all{Begründen Sie mit dem Text} (justifiez à l'aide du texte). Il ne suffit donc pas d'écrire \all{richtig} ou \all{falsch} : il faut \textbf{citer la phrase du texte} qui prouve votre jugement.

\begin{methode}[Réussir l'exercice « vrai ou faux »]
\begin{enumerate}
\item Lire l'affirmation et repérer le passage correspondant dans le texte.
\item Comparer mot à mot : le texte dit-il exactement la même chose, ou le contraire, ou autre chose ?
\item Écrire \all{Richtig} ou \all{Falsch}, puis citer la phrase du texte entre guillemets allemands „ … “.
\item Si l'affirmation est fausse, on peut ajouter la correction en phrase complète.
\end{enumerate}
\end{methode}

\begin{exoresolu}[Richtig oder falsch ? Begründen Sie mit dem Text]
\textbf{1.} Der Schüler musste jeden Morgen früh aufstehen.

\textbf{2.} Das Netz war in seinem Viertel sehr gut.

\textbf{3.} Der Schüler hatte alle Aufgaben geschickt, bevor der Kurs endete.

\textbf{4.} Der Computer kann die Schule ersetzen.
\tcblower
\textbf{1. Falsch.} Le texte dit : \all{„Ich musste nicht früh aufstehen.“} L'affirmation dit le contraire du texte.

\textbf{2. Falsch.} Le texte dit : \all{„Das Netz war in unserem Viertel sehr schlecht.“} L'adjectif est inversé : \all{schlecht} (mauvais) et non \all{gut}.

\textbf{3. Richtig.} Le texte dit : \all{„Bevor der Kurs endete, hatte ich alle Aufgaben geschickt.“} L'affirmation reprend exactement cette information.

\textbf{4. Falsch.} Le texte dit : \all{„Der Computer kann die Schule nicht ersetzen.“} La négation \all{nicht} a été supprimée dans l'affirmation : piège classique !
\end{exoresolu}

\begin{attention}[Les pièges du vrai/faux]
Les affirmations fausses sont fabriquées de trois façons : \textbf{inversion} (\all{schlecht} devient \all{gut}), \textbf{suppression de la négation} (\all{kann nicht ersetzen} devient \all{kann ersetzen}), \textbf{exagération} (\all{oft} « souvent » devient \all{immer} « toujours »). Vérifiez donc systématiquement les négations (\all{nicht}, \all{kein}), les adverbes de fréquence et les adjectifs. Et rappelez-vous : une réponse sans citation du texte est une réponse incomplète.
\end{attention}

\courssub{Recherche des familles de mots dans le texte}

Un autre exercice fréquent consiste à retrouver dans le texte des mots de la même famille qu'un mot donné. Cela entraîne le regard à reconnaître les \textbf{radicaux} et les \textbf{suffixes nominaux} (\all{-ung}, infinitif nominalisé, \all{Er-}).

\begin{center}
\begin{tabularx}{\linewidth}{|l|l|l|X|}
\hline
\rowcolor{popblueL} Verbe (infinitif) & Nom & Traduction & Remarque \\
\hline
verbinden & die Verbindung & la connexion & suffixe \all{-ung} → nom féminin \\
\hline
lernen & das Lernen & l'apprentissage & infinitif nominalisé → neutre \\
\hline
ersetzen & der Ersatz & le remplacement & radical + préfixe \all{Er-} → nom masculin \\
\hline
aufstehen & das Aufstehen & le lever & infinitif nominalisé → neutre \\
\hline
aufgeben & die Aufgabe & le devoir, la tâche & radical + préfixe \all{Auf-} → nom féminin \\
\hline
\end{tabularx}
\end{center}

\begin{propriete}[Former des noms à partir de verbes]
Trois procédés très productifs en allemand :
\begin{itemize}
\item \textbf{radical + -ung} : \all{verbinden → die Verbindung}, \all{üben → die Übung}. Ces noms sont toujours \textbf{féminins}.
\item \textbf{infinitif nominalisé} (majuscule + article \all{das}) : \all{lernen → das Lernen}, \all{lesen → das Lesen}. Ces noms sont toujours \textbf{neutres}.
\item \textbf{nom déverbal sans suffixe} : \all{ersetzen → der Ersatz}, \all{aufgeben → die Aufgabe}. Le genre se mémorise avec le mot.
\end{itemize}
Ce procédé est précieux pour le résumé : un nom remplace toute une phrase. « Le courant tombait en panne » peut devenir \all{der Stromausfall} (mot composé : \all{der Strom} + \all{der Ausfall}).
\end{propriete}

\courssub{Le résumé guidé : die Zusammenfassung}

Dernière étape de la compréhension : produire un résumé du texte en trois ou quatre phrases, avec des connecteurs imposés. Le résumé prouve que l'élève a distingué l'essentiel de l'accessoire.

\begin{propriete}[Les règles du résumé en allemand]
\begin{itemize}
\item \textbf{Temps} : on utilise le \textbf{présent} ou le \textbf{parfait} (Perfekt) ; le prétérit est réservé au récit écrit.
\item \textbf{Sélection} : on supprime les exemples, les détails, les anecdotes ; on garde l'idée principale de chaque paragraphe.
\item \textbf{Connecteurs} : on relie les idées avec les connecteurs imposés, par exemple \all{einerseits … andererseits} (d'une part … d'autre part), \all{aber} (mais), \all{deshalb} (c'est pourquoi), \all{trotzdem} (pourtant).
\item \textbf{Longueur} : trois à quatre phrases maximum ; le résumé ne doit pas dépasser un tiers du texte.
\end{itemize}
\end{propriete}

\begin{exoresolu}[Rédiger une Zusammenfassung en 3-4 phrases]
Consigne : \all{Fassen Sie den Text in drei bis vier Sätzen zusammen. Benutzen Sie : einerseits … andererseits, aber, trotzdem.}
\tcblower
\all{Der Text handelt vom Fernunterricht in Kamerun. Einerseits war er praktisch, weil die Schüler zu Hause lernen konnten, andererseits gab es technische Probleme wie Stromausfälle und ein schlechtes Netz. Der Schüler vermisste auch seine Freunde, aber er hat trotzdem viel gelernt. Der Text zeigt : Fernunterricht ist nützlich, ersetzt aber nicht die Schule.}

« Le texte traite des cours à distance au Cameroun. D'une part ils étaient pratiques parce que les élèves pouvaient apprendre à la maison, d'autre part il y avait des problèmes techniques comme les coupures de courant et un mauvais réseau. L'élève regrettait aussi l'absence de ses amis, mais il a pourtant beaucoup appris. Le texte montre que les cours à distance sont utiles mais ne remplacent pas l'école. »

Vérification : quatre phrases, présent et parfait, aucun détail superflu, les trois connecteurs imposés sont utilisés, la conclusion du texte est conservée.
\end{exoresolu}

\begin{exercice}[À vous de jouer]
\textbf{1.} Répondez en phrases complètes : \all{Warum fühlte sich der Schüler einsam ? Wann hatte er die Seite schon geöffnet ?}

\textbf{2.} \all{Richtig oder falsch ? Begründen Sie mit dem Text.} a) Die Lehrerin war unzufrieden. b) Der Schüler lernte drei Wochen allein.

\textbf{3.} Trouvez dans le texte le nom formé sur le verbe \all{aufstehen} (indice : un infinitif nominalisé) et le nom de la même famille que le verbe \all{aufgeben} (indice : un nom féminin qui désigne un devoir scolaire).

\textbf{4.} Rédigez votre propre résumé du texte en trois phrases avec \all{zuerst}, \all{dann} et \all{am Ende}.
\end{exercice}

\begin{corrige}[Exercice « À vous de jouer »]
\textbf{1.} \all{Der Schüler fühlte sich einsam, weil er drei Wochen allein gelernt hatte und seine Freunde vermisste.} « L'élève se sentait seul parce qu'il avait appris seul pendant trois semaines et que ses amis lui manquaient. » — \all{Er hatte die Seite schon geöffnet, als der Strom wieder ausfiel.} « Il avait déjà ouvert la page quand le courant est retombé en panne. »

\textbf{2.} a) \textbf{Falsch.} Le texte dit : \all{„Meine Lehrerin war zufrieden.“} b) \textbf{Richtig.} Le texte dit : \all{„Nachdem ich drei Wochen allein gelernt hatte, fühlte ich mich einsam.“}

\textbf{3.} \all{aufstehen → das Aufstehen} (le lever) ; \all{aufgeben → die Aufgabe} (le devoir, la tâche).

\textbf{4.} Modèle : \all{Zuerst erzählt der Text, dass der Fernunterricht praktisch war. Dann beschreibt er die technischen und sozialen Probleme. Am Ende sagt der Schreiber, dass Fernunterricht nützlich ist, aber die Schule nicht ersetzt.}
\end{corrige}

\begin{aretenir}[L'essentiel de la section]
Pour exploiter un texte : répondre toujours en \textbf{phrase complète en allemand}, en reprenant les termes de la question ; justifier le vrai/faux par une \textbf{citation exacte} du texte entre guillemets „ … “ ; reconnaître les familles de mots (\all{-ung} → féminin, infinitif nominalisé → neutre) ; résumer au \textbf{présent ou au parfait}, en supprimant les détails et en gardant l'idée principale de chaque paragraphe, avec les connecteurs imposés. Ces quatre savoir-faire seront réinvestis dans la section 6 consacrée à la production guidée.
\end{aretenir}

\coursec{Lexique thématique : der Fernunterricht und das Digitale}

Dans la section précédente, nous avons lu et compris le texte \all{Fernunterricht – Chancen und Probleme} : nous savons maintenant de quoi il parle et quelles idées il défend. Mais comprendre globalement ne suffit pas au baccalauréat : il faut pouvoir \emph{réemployer} les mots du texte dans vos réponses, vos résumés et vos productions. Cette section est donc consacrée au \cle{lexique thématique} des cours à distance et du numérique. Nous allons observer les mots dans leur contexte, les classer par champs lexicaux, mémoriser systématiquement leur \cle{article} et leur \cle{pluriel}, construire des familles de mots et apprendre des expressions prêtes à l'emploi.

\courssub{Observer d'abord : les mots dans un contexte}

Avant toute liste, lisons ce court texte original qui réutilise presque tout le vocabulaire de la leçon. Soulignez mentalement les noms et leurs articles.

\begin{textebox}[Ein Schüler aus Yaoundé erzählt]
Ich heiße Arnaud und ich besuche die Terminale an einem Gymnasium in Yaoundé. Letztes Jahr hatten wir oft Fernunterricht, weil die Schule manchmal geschlossen war. Unser Lehrer hat eine Plattform benutzt, und wir haben die Lektionen per Videokonferenz verfolgt. Das Lernen zu Hause hat Vorteile: Ich musste nicht mit dem Bus fahren, und ich konnte die Lektionen wiederholen. Aber es gab auch viele Probleme. Die Verbindung war oft schlecht, das Netz war langsam, und zweimal ist der Strom ausgefallen – mitten in der Videokonferenz! Mein Handy war mein einziger Computer, und der Bildschirm war sehr klein. Einmal habe ich mein Passwort vergessen und konnte mich nicht anmelden. Einerseits ist der Fernunterricht praktisch und nützlich, andererseits kann er den Kontakt zum Lehrer und zu den Freunden nicht ersetzen. Ich denke: Der Fernunterricht ist eine Chance, aber auch eine Gefahr, wenn die Technik nicht funktioniert.
\end{textebox}

\textbf{Glossaire (ordre d'apparition) :} \all{der Fernunterricht} : les cours à distance ; \all{die Plattform} : la plateforme ; \all{die Videokonferenz} : la visioconférence ; \all{verfolgen} : suivre (un cours) ; \all{die Verbindung} : la connexion ; \all{das Netz} : le réseau ; \all{der Strom} : le courant électrique ; \all{ausfallen} : tomber en panne, être coupé ; \all{der Bildschirm} : l'écran ; \all{das Passwort} : le mot de passe ; \all{sich anmelden} : se connecter (à un compte) ; \all{ersetzen} : remplacer ; \all{die Gefahr} : le danger.

\textbf{Questions rapides :} 1) Welche Geräte benutzt Arnaud? 2) Welche Probleme hat er? 3) Findet er den Fernunterricht positiv oder negativ?

\begin{corrige}[Questions rapides]
1) Er benutzt sein Handy (als Computer). 2) Die Verbindung ist schlecht, das Netz ist langsam, der Strom fällt aus, der Bildschirm ist klein, und er vergisst sein Passwort. 3) Er findet ihn einerseits praktisch und nützlich, andererseits problematisch: Er ist eine Chance, aber auch eine Gefahr.
\end{corrige}

\courssub{Champ lexical 1 : les cours à distance}

\begin{definition}[der Fernunterricht]
\all{Der Fernunterricht} (meist Singular) désigne l'enseignement à distance : \all{fern} = lointain, à distance ; \all{der Unterricht} = le cours, l'enseignement. On dit : \all{Fernunterricht haben} (avoir des cours à distance), \all{am Fernunterricht teilnehmen} (participer aux cours à distance).
\end{definition}

\begin{center}
\begin{tabularx}{\linewidth}{|X|X|X|}
\hline
\rowcolor{popblueL} Deutsch & Français & Beispiel \\
\hline
der Fernunterricht (Sg.) & les cours à distance & \all{Der Fernunterricht beginnt um 8 Uhr.} \\
\hline
der Unterricht (-e) & le cours, l'enseignement & \all{Der Unterricht ist heute online.} \\
\hline
die Lektion (-en) & la leçon & \all{Die Lektion ist schwierig.} \\
\hline
die Videokonferenz (-en) & la visioconférence & \all{Die Videokonferenz dauert eine Stunde.} \\
\hline
die Plattform (-en) & la plateforme & \all{Die Plattform heißt OnBuch.} \\
\hline
das Lernen (Sg.) & l'apprentissage & \all{Das Lernen zu Hause ist ruhig.} \\
\hline
der Schüler (-) / die Schülerin (-nen) & l'élève (m./f.) & \all{Die Schülerin lernt online.} \\
\hline
der Kurs (-e) & le cours, la formation & \all{Ich besuche einen Deutschkurs.} \\
\hline
\end{tabularx}
\end{center}

\courssub{Champ lexical 2 : la technique et les appareils}

\begin{center}
\begin{tabularx}{\linewidth}{|X|X|X|}
\hline
\rowcolor{popblueL} Deutsch & Français & Beispiel \\
\hline
die Verbindung (-en) & la connexion, la liaison & \all{Die Verbindung ist schlecht.} \\
\hline
das Netz (-e) & le réseau & \all{Das Netz ist heute langsam.} \\
\hline
das Internet (-s) & l'Internet & \all{Das Internet funktioniert nicht.} \\
\hline
der Strom (Sg.) & le courant électrique & \all{Der Strom ist ausgefallen.} \\
\hline
der Computer (-) & l'ordinateur & \all{Der Computer ist alt.} \\
\hline
das Handy (-s) & le portable & \all{Mein Handy hat kein Netz.} \\
\hline
der Bildschirm (-e) & l'écran & \all{Der Bildschirm ist klein.} \\
\hline
die App (-s) & l'application & \all{Die App ist kostenlos.} \\
\hline
das Passwort (die Passwörter) & le mot de passe & \all{Ich habe das Passwort vergessen.} \\
\hline
der Laptop (-s) & l'ordinateur portable & \all{Der Laptop ist neu.} \\
\hline
\end{tabularx}
\end{center}

\begin{definition}[die Verbindung]
\all{Die Verbindung} (-en) signifie la connexion, la liaison. Famille de mots : le verbe \all{verbinden} (relier, connecter) → \all{die Verbindung} (la connexion) → \all{verbunden} (relié, connecté). Expressions : \all{eine gute/schlechte Verbindung haben} (avoir une bonne/mauvaise connexion).
\end{definition}

\begin{attention}[Faux amis : der Strom et das Netz]
\all{Der Strom} signifie \emph{le courant électrique} uniquement ; il ne traduit jamais « le courant » au sens figuré (« être au courant » = \all{informiert sein}). De même, \all{das Netz} est \emph{le réseau} (de téléphone, de données) : ne le confondez pas avec \all{das Internet}. On dit : \all{Ich habe kein Netz} (je n'ai pas de réseau), mais \all{Ich bin im Internet} (je suis sur Internet).
\end{attention}

\courssub{Champ lexical 3 : les problèmes}

\begin{center}
\begin{tabularx}{\linewidth}{|X|X|X|}
\hline
\rowcolor{popblueL} Deutsch & Français & Beispiel \\
\hline
der Stromausfall (die Stromausfälle) & la coupure de courant & \all{Der Stromausfall dauerte zwei Stunden.} \\
\hline
die Störung (-en) & la panne, le dérangement & \all{Es gibt eine Störung im Netz.} \\
\hline
langsam & lent & \all{Das Internet ist sehr langsam.} \\
\hline
ausfallen (fällt aus, fiel aus, ist ausgefallen) & tomber en panne, être annulé & \all{Der Unterricht fällt heute aus.} \\
\hline
fehlen & manquer, faire défaut & \all{Mir fehlt ein Computer.} \\
\hline
das Problem (-e) & le problème & \all{Das Problem ist das Netz.} \\
\hline
die Schwierigkeit (-en) & la difficulté & \all{Ich habe Schwierigkeiten mit der App.} \\
\hline
\end{tabularx}
\end{center}

\begin{propriete}[Le verbe fort ausfallen]
\all{Ausfallen} est un verbe fort à particule séparable : Präsens \all{er fällt aus}, Präteritum \all{er fiel aus}, Perfekt \all{er ist ausgefallen} (auxiliaire \all{sein}, car il exprime un changement d'état). Attention à la place de la particule : \all{Der Strom ist gestern ausgefallen.} (Le courant a été coupé hier.) Autre sens : \all{Der Unterricht fällt aus} = le cours est annulé.
\end{propriete}

\begin{exemplebox}[Exemples traduits]
\all{Die Verbindung ist schlecht.} « La connexion est mauvaise. »

\all{Der Strom ist ausgefallen.} « Le courant a été coupé. »

\all{Der Fernunterricht hat Vor- und Nachteile.} « Les cours à distance ont des avantages et des inconvénients. »

\all{Mir fehlt das Geld für einen Laptop.} « Il me manque l'argent pour un ordinateur portable. » (Notez : \all{fehlen} se construit avec le datif de la personne.)
\end{exemplebox}

\courssub{Champ lexical 4 : évaluer, donner son avis}

\begin{center}
\begin{tabularx}{\linewidth}{|X|X|X|}
\hline
\rowcolor{popblueL} Deutsch & Français & Beispiel \\
\hline
der Vorteil (-e) & l'avantage & \all{Ein Vorteil ist die Flexibilität.} \\
\hline
der Nachteil (-e) & l'inconvénient & \all{Ein Nachteil ist die schlechte Verbindung.} \\
\hline
die Chance (-n) & la chance, l'opportunité & \all{Der Fernunterricht ist eine Chance.} \\
\hline
die Gefahr (-en) & le danger & \all{Die Gefahr ist die Isolation.} \\
\hline
nützlich & utile & \all{Die App ist sehr nützlich.} \\
\hline
praktisch & pratique & \all{Online lernen ist praktisch.} \\
\hline
gefährlich & dangereux & \all{Zu viel Bildschirm ist gefährlich.} \\
\hline
ersetzen & remplacer & \all{Der Computer ersetzt den Lehrer nicht.} \\
\hline
\end{tabularx}
\end{center}

\courssub{Familles de mots : un verbe, plusieurs mots}

L'allemand construit des familles de mots très régulières autour d'un verbe. Les mémoriser ensemble multiplie votre vocabulaire sans effort supplémentaire.

\begin{center}
\begin{tabularx}{\linewidth}{|X|X|X|}
\hline
\rowcolor{popgreenL} Verbe & Noms dérivés & Adjectif / participe \\
\hline
verbinden (relier) & die Verbindung (-en) & verbunden (connecté) \\
\hline
lernen (apprendre) & das Lernen, der Lerner (-) & — \\
\hline
nützen (servir) & der Nutzen (l'utilité) & nützlich (utile) \\
\hline
teilnehmen (participer) & die Teilnahme (la participation) & — \\
\hline
unterrichten (enseigner) & der Unterricht (-e), der Fernunterricht & — \\
\hline
stören (déranger) & die Störung (-en) & gestört (perturbé) \\
\hline
\end{tabularx}
\end{center}

\begin{propriete}[Règle d'or : toujours apprendre article + pluriel]
En allemand, un nom ne s'apprend jamais seul. Mémorisez toujours la forme complète : \all{der Vorteil, -e} ; \all{die Verbindung, -en} ; \all{das Netz, -e} ; \all{das Passwort, die Passwörter}. Le genre détermine toute la grammaire de la phrase (articles, adjectifs, pronoms). Astuce : les noms en \all{-ung}, \all{-heit}, \all{-keit}, \all{-ion} sont presque toujours féminins (\all{die Verbindung, die Schwierigkeit, die Lektion}).
\end{propriete}

\begin{aretenir}[Expressions utiles à réemployer]
\begin{itemize}
\item \all{online lernen} : apprendre en ligne
\item \all{eine gute/schlechte Verbindung haben} : avoir une bonne/mauvaise connexion
\item \all{der Strom fällt aus} : le courant est coupé
\item \all{Kontakt halten} : garder le contact
\item \all{einen Kurs besuchen} : suivre un cours
\item \all{am Fernunterricht teilnehmen} : participer aux cours à distance
\end{itemize}
\end{aretenir}

\courssub{Carte mentale : das Digitale}

\begin{popfigure}
\begin{center}
\begin{tikzpicture}[mindmap, grow cyclic, every node/.style={concept, align=center, font=\small},
root concept/.append style={concept color=popblue, font=\bfseries},
level 1/.append style={level distance=3.2cm, sibling angle=90}]
\node[root concept]{das Digitale}
  child[concept color=poporange]{ node[align=center]{Geräte\\ \all{das Handy}\\ \all{der Computer}\\ \all{der Bildschirm}} }
  child[concept color=popgreen]{ node[align=center]{Internet\\ \all{das Netz}\\ \all{die Verbindung}\\ \all{die Plattform}} }
  child[concept color=poppink]{ node[align=center]{Probleme\\ \all{der Stromausfall}\\ \all{die Störung}\\ \all{langsam}} }
  child[concept color=popgold]{ node[align=center]{Meinung\\ \all{der Vorteil}\\ \all{der Nachteil}\\ \all{nützlich}} };
\end{tikzpicture}
\end{center}
\legende{Carte mentale du vocabulaire du numérique : quatre branches à apprendre par cœur.}
\end{popfigure}

\begin{savaistu}[Das Handy : un mot « anglais » inventé en Allemagne]
L'allemand du numérique emprunte beaucoup à l'anglais : \all{der Laptop}, \all{das WLAN}, \all{die App}, \all{das Passwort}. Mais attention au piège amusant : \all{das Handy} (le téléphone portable) \emph{n'existe pas en anglais} ! C'est un mot inventé en Allemagne qui a l'air anglais — un « faux anglicisme ». Un Anglais dira \emph{mobile phone}, un Américain \emph{cell phone}. Au Cameroun, où le portable est souvent le seul outil pour suivre les cours en ligne, \all{das Handy} est un mot à connaître absolument.
\end{savaistu}

\begin{exoresolu}[Compléter avec le bon mot et le bon article]
Énoncé : Complétez avec le mot qui convient (avec l'article) : \all{Verbindung, Strom, Passwort, Vorteil, Plattform, Netz}.

1) Ich kann mich nicht anmelden – ich habe mein \ldots\ vergessen. \quad 2) \ldots\ ist heute sehr langsam, die Seite lädt nicht. \quad 3) Ein großer \ldots\ des Fernunterrichts: Man muss nicht zur Schule fahren. \quad 4) \ldots\ ist ausgefallen, wir hatten kein Licht mehr. \quad 5) Unser Lehrer benutzt eine neue \ldots\ für die Lektionen. \quad 6) Die \ldots\ ist schlecht, ich höre den Lehrer nicht.
\tcblower
Solution détaillée : 1) \all{mein Passwort} (\all{das Passwort}, accusatif neutre après \all{vergessen}) : « j'ai oublié mon mot de passe ». 2) \all{Das Netz} (nominatif, sujet de \all{ist}) : le réseau est lent. 3) \all{Vorteil} (\all{ein großer Vorteil}, nominatif masculin ; notez le génitif \all{des Fernunterrichts}) : un grand avantage. 4) \all{Der Strom} (nominatif, sujet de \all{ist ausgefallen}) : le courant a été coupé. 5) \all{Plattform} (\all{eine neue Plattform}, accusatif féminin après \all{benutzen}). 6) \all{Verbindung} (\all{die Verbindung}, nominatif féminin) : la connexion est mauvaise.
\end{exoresolu}

\begin{attention}[Erreurs fréquentes des francophones]
\begin{itemize}
\item Ne dites pas \all{*das Verbindung*} : c'est \all{die Verbindung} (féminin, comme tous les noms en \all{-ung}).
\item « Suivre un cours » se dit \all{einen Kurs besuchen} ou \all{einer Lektion folgen} (avec le datif !), pas \all{*einen Kurs folgen*}.
\item \all{Der Unterricht fällt aus} signifie « le cours est annulé », pas « le cours tombe ».
\item N'oubliez pas la majuscule à tous ces noms : \all{das Internet, das Handy, das Netz}.
\item \all{teilnehmen} se construit avec \all{an} + datif : \all{Ich nehme am Fernunterricht teil} (et la particule \all{teil} va en fin de phrase).
\end{itemize}
\end{attention}

\begin{exercice}[À vous]
1) Traduisez : a) La connexion est mauvaise aujourd'hui. b) Le courant a été coupé pendant la visioconférence. c) Les cours à distance ont des avantages et des inconvénients. \quad 2) Construisez trois phrases avec : \all{das Netz}, \all{nützlich}, \all{Kontakt halten}.
\end{exercice}

\begin{corrige}[Exercice « À vous »]
1) a) \all{Die Verbindung ist heute schlecht.} b) \all{Der Strom ist während der Videokonferenz ausgefallen.} c) \all{Der Fernunterricht hat Vor- und Nachteile.} 2) Exemples : \all{Das Netz ist in meinem Quartier sehr langsam.} / \all{Die App ist für das Lernen sehr nützlich.} / \all{Mit dem Handy kann ich Kontakt zu meinen Freunden halten.}
\end{corrige}

Ce lexique est désormais votre boîte à outils : gardez-le à portée de main, car les sections suivantes — l'opposition avec \all{einerseits … andererseits} et le plus-que-parfait — vous demanderont de réemployer ces mots dans des phrases complètes et un court texte d'opinion.

\coursec{Grammaire 1 : l'opposition avec \all{einerseits} … \all{andererseits}}

Dans la section précédente, nous avons construit le lexique du \cle{Fernunterricht} : \all{der Vorteil}, \all{der Nachteil}, \all{die Verbindung}, \all{das Netz}. Mais posséder des mots ne suffit pas : au baccalauréat, on vous demandera souvent de \emph{présenter les avantages et les inconvénients} d'un phénomène, c'est-à-dire d'organiser une argumentation \emph{nuancée}. L'outil grammatical roi pour cela est la structure d'opposition \all{einerseits … andererseits}. C'est elle que nous allons observer, comprendre et pratiquer maintenant.

\courssub{Observation : repérer la structure dans le texte}

Relisons deux phrases clés du texte support \emph{Fernunterricht – Chancen und Probleme} :

\begin{exemplebox}[Deux phrases du texte support]
\all{Einerseits spart man beim Fernunterricht Zeit und Geld, andererseits fehlt der direkte Kontakt zu den Lehrern.}

« D'une part, on gagne du temps et de l'argent avec les cours à distance, d'autre part, le contact direct avec les professeurs manque. »

\all{Einerseits ist das Lernen zu Hause bequem, andererseits ist es oft anstrengend und einsam.}

« D'une part, apprendre à la maison est confortable, d'autre part, c'est souvent fatigant et solitaire. »
\end{exemplebox}

Observons attentivement. Que remarquez-vous ?

\begin{itemize}
\item Chaque phrase est composée de \textbf{deux parties} séparées par une virgule.
\item La première partie commence par \all{einerseits}, la deuxième par \all{andererseits}.
\item La première partie présente un \textbf{aspect positif} (\all{Zeit und Geld sparen}, \all{bequem}), la deuxième un \textbf{aspect négatif} (\all{der Kontakt fehlt}, \all{anstrengend}).
\item Juste après \all{einerseits} et \all{andererseits}, on trouve immédiatement le \textbf{verbe conjugué} (\all{spart}, \all{fehlt}, \all{ist}) : le sujet vient \emph{après} le verbe.
\end{itemize}

Nous venons de découvrir par l'observation les deux caractéristiques essentielles de cette structure : sa \textbf{symétrie} (deux aspects opposés) et sa \textbf{syntaxe} (l'inversion du sujet). Formalisons maintenant la règle.

\courssub{La règle : forme, place et emploi}

\begin{propriete}[Règle de la structure \all{einerseits} … \all{andererseits}]
\all{Einerseits … andererseits} signifie « d'une part … d'autre part ». La structure encadre une argumentation en deux temps : \all{einerseits} introduit le \textbf{premier aspect} (par exemple les avantages), \all{andererseits} introduit l'\textbf{aspect opposé} (par exemple les inconvénients).

\textbf{Syntaxe obligatoire :} chaque connecteur occupe la \textbf{1\textsuperscript{re} position} de sa proposition. Comme tout élément placé en tête de phrase en allemand, il provoque l'\textbf{inversion du sujet} : le verbe conjugué reste en \textbf{2\textsuperscript{e} position}, et le sujet se place \textbf{après le verbe}.

\begin{center}
\begin{tabular}{|l|l|l|l|}
\hline
\rowcolor{popblueL} Position 1 & Position 2 (verbe) & Sujet & Reste de la phrase \\
\hline
Einerseits & spart & man & Zeit und Geld, \\
\hline
andererseits & fehlt & der Kontakt & zu den Lehrern. \\
\hline
\end{tabular}
\end{center}

Les deux propositions sont séparées par une \textbf{virgule}. On peut aussi construire deux phrases séparées : \all{Einerseits spart man Zeit. Andererseits fehlt der Kontakt.}
\end{propriete}

\begin{aretenir}[Emplois de la structure]
On utilise \all{einerseits … andererseits} pour :
\begin{itemize}
\item présenter les \textbf{avantages et les inconvénients} d'une situation : \all{Einerseits ist das Netz in Yaoundé schnell, andererseits fällt der Strom oft aus.} « D'une part, le réseau est rapide à Yaoundé, d'autre part, l'électricité coupe souvent. »
\item opposer \textbf{deux points de vue} : \all{Einerseits lieben die Schüler das Handy, andererseits stört es den Unterricht.} « D'une part, les élèves adorent le téléphone portable, d'autre part, il perturbe le cours. »
\item établir un \textbf{bilan nuancé} dans un essai ou un commentaire de texte : c'est la structure idéale pour la conclusion d'un devoir d'expression écrite au Bac.
\end{itemize}
\end{aretenir}

\begin{popfigure}
\begin{tikzpicture}[scale=1.0, every node/.style={font=\small}]
% Poteau central
\draw[fill=popdark, popdark] (-0.15,0) rectangle (0.15,3.2);
\draw[fill=popdark, popdark] (-1.2,0) rectangle (1.2,0.25);
% Flèche centrale / pivot
\draw[fill=popgold, popgold] (0,3.2) circle (0.12);
% Fléau
\draw[line width=2pt, popblue] (-3.6,3.6) -- (3.6,3.6);
% Cordes gauches
\draw[popblue] (-3.6,3.6) -- (-4.4,2.4);
\draw[popblue] (-3.6,3.6) -- (-2.8,2.4);
% Cordes droites
\draw[popblue] (3.6,3.6) -- (2.8,2.4);
\draw[popblue] (3.6,3.6) -- (4.4,2.4);
% Plateaux
\draw[fill=popgreenL, draw=popgreen, line width=1pt] (-4.4,2.4) arc (180:360:0.8) -- cycle;
\draw[fill=poporangeL, draw=poporange, line width=1pt] (2.8,2.4) arc (180:360:0.8) -- cycle;
% Étiquettes
\node[align=center, fill=popgreenL, draw=popgreen, rounded corners, inner sep=4pt] at (-3.6,1.2) {\textbf{einerseits}\\ die Vorteile\\ \all{Zeit sparen, bequem}};
\node[align=center, fill=poporangeL, draw=poporange, rounded corners, inner sep=4pt] at (3.6,1.2) {\textbf{andererseits}\\ die Nachteile\\ \all{kein Kontakt, Strom weg}};
\node[align=center, fill=popblueL, draw=popblue, rounded corners, inner sep=4pt] at (0,4.3) {L'opposition équilibrée : inversion du sujet des deux côtés};
\end{tikzpicture}
\legende{La balance de l'argumentation : \all{einerseits} pèse les avantages, \all{andererseits} les inconvénients ; la phrase reste en équilibre grâce à la même syntaxe des deux côtés.}
\end{popfigure}

\courssub{Comparaison avec le français et alternatives pour le Bac}

En français, « d'une part, on gagne du temps » : le sujet « on » reste \emph{avant} le verbe. En allemand, c'est impossible : \all{einerseits} occupe la première position, donc le verbe doit suivre immédiatement. C'est la différence majeure entre les deux langues.

\begin{center}
\begin{tabularx}{\linewidth}{|X|X|X|}
\hline
\rowcolor{popblueL} Français & Deutsch & Remarque \\
\hline
D'une part, on gagne du temps. & Einerseits spart man Zeit. & Inversion obligatoire : le verbe \all{spart} précède le sujet \all{man}. \\
\hline
D'autre part, le contact manque. & Andererseits fehlt der Kontakt. & Même règle dans la deuxième proposition. \\
\hline
Certes c'est pratique, mais c'est fatigant. & Zwar ist es praktisch, aber es ist anstrengend. & \all{zwar … aber} : concession + opposition. \\
\hline
Par contre, la connexion est lente. & Im Gegensatz dazu ist die Verbindung langsam. & \all{im Gegensatz dazu} : « par contre », avec inversion aussi. \\
\hline
D'un côté … de l'autre côté … & Auf der einen Seite … auf der anderen Seite … & Variante plus longue, très appréciée à l'écrit. \\
\hline
\end{tabularx}
\end{center}

\begin{definition}[Les alternatives utiles au Bac]
Pour varier votre expression écrite, remplacez parfois \all{einerseits … andererseits} par :
\begin{itemize}
\item \all{auf der einen Seite … auf der anderen Seite} : « d'un côté … de l'autre côté » — \all{Auf der einen Seite ist Fernunterricht flexibel, auf der anderen Seite braucht man ein gutes Netz.} « D'un côté, le cours à distance est flexible, de l'autre, il faut un bon réseau. »
\item \all{zwar … aber} : « certes … mais » — \all{Das Handy ist zwar nützlich, aber es lenkt ab.} « Le portable est certes utile, mais il distrait. »
\item \all{im Gegensatz dazu} : « par contre » — \all{In der Stadt ist das Netz schnell. Im Gegensatz dazu ist es auf dem Land oft schlecht.} « En ville, le réseau est rapide. Par contre, il est souvent mauvais à la campagne. »
\end{itemize}
Attention : \all{auf der einen Seite} et \all{im Gegensatz dazu} occupent aussi la 1\textsuperscript{re} position et provoquent l'inversion ; \all{aber} est en position 0 et ne la provoque pas.
\end{definition}

\begin{exemplebox}[Exemples variés autour du thème de la leçon]
\all{Einerseits können die Schüler in Douala die Lektionen wiederholen, andererseits kosten die Daten viel Geld.} « D'une part, les élèves de Douala peuvent réécouter les leçons, d'autre part, les données coûtent cher. »

\all{Einerseits braucht man für den Fernunterricht nur ein Handy, andererseits ist ein Computer besser für die Hausaufgaben.} « D'une part, on n'a besoin que d'un téléphone pour les cours à distance, d'autre part, un ordinateur est meilleur pour les devoirs. »

\all{Einerseits lernen die Kinder selbstständig, andererseits brauchen sie die Hilfe der Eltern.} « D'une part, les enfants apprennent de manière autonome, d'autre part, ils ont besoin de l'aide des parents. »

\all{Einerseits war die Schule während der Krise geschlossen, andererseits ging der Unterricht online weiter.} « D'une part, l'école était fermée pendant la crise, d'autre part, les cours continuaient en ligne. »
\end{exemplebox}

\courssub{Texte d'application : un débat au lycée}

\begin{textebox}[Ein Streitgespräch im Lycée : „Fernunterricht – ja oder nein ?"]
Nach der Krise diskutieren die Schüler der Terminale in Yaoundé über den Fernunterricht. Die Lehrerin hat eine Frage an die Tafel geschrieben: „Ist das Lernen zu Hause besser als der Unterricht in der Schule ?"

Amina steht auf und sagt: „Einerseits ist der Fernunterricht sehr praktisch. Ich muss nicht mit dem Bendskin durch den Verkehr fahren, und ich spare Zeit und Geld. Andererseits sitze ich den ganzen Tag allein vor dem Bildschirm, und das ist traurig. Wenn das Netz schlecht ist, verstehe ich die Erklärungen der Lehrer nicht."

Ihr Freund Martial nickt und ergänzt: „Einerseits kann ich die Videos stoppen und die Lektion wiederholen, wenn ich etwas nicht verstehe. Andererseits fehlt der Kontakt zu meinen Freunden. In der Pause spielen wir normalerweise Fußball – zu Hause ist das unmöglich."

Ein dritter Schüler, Etienne, denkt an seine Familie auf dem Land: „Einerseits hat mein Cousin im Dorf ein Smartphone, andererseits gibt es dort oft keinen Strom. Wie soll er dann lernen ? Für ihn ist der Fernunterricht eine Illusion."

Die Lehrerin lächelt: „Sehr gut! Ihr benutzt die Struktur \emph{einerseits … andererseits} perfekt. Zum Schluss schreibt bitte einen kurzen Text mit eurer persönlichen Meinung. Beginnt mit den Vorteilen, nennt dann die Nachteile und endet mit \emph{ich denke, dass …}"

Am Ende der Stunde sind sich alle einig: Der Fernunterricht ist einerseits eine große Chance, andererseits darf er die Schule nicht ersetzen.

\textbf{Glossaire :} das Streitgespräch, -e (débat) ; die Tafel, -n (le tableau) ; der Bendskin, -s (moto-taxi) ; der Verkehr (la circulation) ; der Bildschirm, -e (l'écran) ; nicken (hocher la tête) ; ergänzen (ajouter, compléter) ; die Pause, -n (la récréation) ; das Dorf, „-er (le village) ; die Illusion, -en (l'illusion) ; ersetzen (remplacer) ; sich einig sein (être d'accord).

\textbf{Questions de compréhension :} 1. Welche Vorteile nennt Amina ? 2. Was fehlt Martial zu Hause ? 3. Warum ist der Fernunterricht für Etiennes Cousin schwierig ? 4. Was sollen die Schüler am Ende schreiben ?
\end{textebox}

\courssub{Entraînement guidé : transformer des phrases simples}

\begin{methode}[Construire une phrase d'opposition en quatre étapes]
\begin{enumerate}
\item \textbf{Identifier} les deux idées opposées (un avantage, un inconvénient).
\item \textbf{Placer} \all{einerseits} en tête de la première proposition, \all{andererseits} en tête de la seconde.
\item \textbf{Vérifier l'inversion} : verbe conjugué en 2\textsuperscript{e} position, sujet après le verbe, dans les deux propositions.
\item \textbf{Relier} les deux propositions par une virgule et relire la phrase entière.
\end{enumerate}
\end{methode}

\begin{exoresolu}[Transformation de phrases simples]
Énoncé : transformez les deux phrases simples suivantes en une seule phrase d'opposition avec \all{einerseits … andererseits} : \all{Man spart Zeit und Geld. Der Kontakt zu den Lehrern fehlt.}
\tcblower
Solution détaillée :

Étape 1 : la première phrase exprime un avantage, la deuxième un inconvénient — l'opposition est possible.

Étape 2 : on place les connecteurs en tête : \all{Einerseits …} / \all{andererseits …}.

Étape 3 : on applique l'inversion. Le verbe \all{spart} passe devant le sujet \all{man} ; le verbe \all{fehlt} passe devant le sujet \all{der Kontakt zu den Lehrern}.

Étape 4 : phrase finale : \all{Einerseits spart man Zeit und Geld, andererseits fehlt der Kontakt zu den Lehrern.} « D'une part, on gagne du temps et de l'argent, d'autre part, le contact avec les professeurs manque. »
\end{exoresolu}

\begin{exercice}[À vous de jouer]
Transformez les paires de phrases suivantes en phrases d'opposition :
\begin{enumerate}
\item \all{Das Lernen zu Hause ist bequem. Es ist anstrengend.}
\item \all{Die Schüler können die Videos wiederholen. Die Verbindung ist oft schlecht.}
\item \all{Der Fernunterricht ist flexibel. Viele Familien haben keinen Computer.}
\end{enumerate}
\end{exercice}

\begin{corrige}[Exercice « À vous de jouer »]
1. \all{Einerseits ist das Lernen zu Hause bequem, andererseits ist es anstrengend.} « D'une part, apprendre à la maison est confortable, d'autre part, c'est fatigant. »

2. \all{Einerseits können die Schüler die Videos wiederholen, andererseits ist die Verbindung oft schlecht.} « D'une part, les élèves peuvent réécouter les vidéos, d'autre part, la connexion est souvent mauvaise. » (Avec un verbe modal, c'est le modal conjugué qui prend la 2\textsuperscript{e} position.)

3. \all{Einerseits ist der Fernunterricht flexibel, andererseits haben viele Familien keinen Computer.} « D'une part, le cours à distance est flexible, d'autre part, beaucoup de familles n'ont pas d'ordinateur. »
\end{corrige}

\begin{attention}[Erreurs fréquentes des francophones]
\textbf{1. Oublier l'inversion.} Le français « d'une part, on gagne du temps » tente de calquer : \all{Einerseits man spart Zeit.} ✗ Faux ! Il faut : \all{Einerseits spart man Zeit.} ✓ Le verbe conjugué suit immédiatement le connecteur.

\textbf{2. Placer \all{andererseits} sans inversion dans la deuxième partie.} \all{…, andererseits der Kontakt fehlt.} ✗ → \all{…, andererseits fehlt der Kontakt.} ✓

\textbf{3. Traduire mot à mot « d'une part ».} On ne dit jamais \all{auf einer Hand} ✗ : la forme correcte est \all{einerseits} (un seul mot, adverbe) ou \all{auf der einen Seite}.

\textbf{4. Utiliser la structure pour deux idées qui ne s'opposent pas.} \all{Einerseits … andererseits} exige un contraste réel (avantage/inconvénient, point de vue A/point de vue B). Pour une simple addition, utilisez \all{und} ou \all{außerdem} (« de plus »).
\end{attention}

\begin{methode}[Rédiger un paragraphe d'opinion équilibré au Bac]
\begin{enumerate}
\item \textbf{Introduction} : une phrase qui pose le sujet, par exemple \all{Viele Schüler lernen heute zu Hause.} « Beaucoup d'élèves apprennent aujourd'hui à la maison. »
\item \textbf{Deux arguments \all{einerseits}} : \all{Einerseits spart man Zeit, und man kann die Lektionen wiederholen.}
\item \textbf{Deux arguments \all{andererseits}} : \all{Andererseits fehlt der Kontakt zu den Lehrern, und das Netz ist oft schlecht.}
\item \textbf{Conclusion personnelle} avec une subordonnée : \all{Ich denke, dass der Fernunterricht eine gute Ergänzung, aber kein Ersatz für die Schule ist.} « Je pense que le cours à distance est un bon complément, mais pas un remplacement de l'école. » (Attention : après \all{dass}, le verbe conjugué va à la fin.)
\end{enumerate}
\end{methode}

\begin{aretenir}[L'essentiel de la section]
\all{Einerseits … andererseits} = « d'une part … d'autre part » : deux aspects opposés, deux propositions symétriques, et une règle d'or : le connecteur en 1\textsuperscript{re} position, le verbe en 2\textsuperscript{e}, le sujet après le verbe. Variantes : \all{auf der einen Seite … auf der anderen Seite}, \all{zwar … aber}, \all{im Gegensatz dazu}. Cette structure sera votre meilleure alliée pour le commentaire de texte et l'expression écrite. Dans la section suivante, nous verrons comment raconter les expériences passées du cours à distance avec le plus-que-parfait et les propositions temporelles.
\end{aretenir}

\coursec{Grammaire 2 : le plus-que-parfait et les propositions temporelles}

Dans la section précédente, nous avons appris à opposer les avantages et les inconvénients du cours à distance grâce à \all{einerseits \dots\ andererseits}. Mais un commentaire de texte ne se contente pas d'opposer des idées : il raconte aussi une suite d'événements. Or le texte \emph{Fernunterricht – Chancen und Probleme} est construit sur une chronologie : d'abord les écoles ont fermé, \emph{ensuite} les cours en ligne ont commencé. Comment l'allemand exprime-t-il cet « avant » et cet « après » ? C'est toute la question du plus-que-parfait et des propositions temporelles, que nous allons maintenant étudier.

\courssub{Observation : deux actions, deux moments}

Relisons la phrase clé du texte support :

\begin{exemplebox}[Phrase d'observation]
\all{Nachdem die Schulen geschlossen hatten, begann der Unterricht per Internet.}

« Après que les écoles eurent fermé, les cours commencèrent par Internet. »
\end{exemplebox}

Posons-nous trois questions :

\begin{enumerate}
\item Combien y a-t-il d'actions ? \textbf{Deux} : \all{die Schulen schließen} (les écoles ferment) et \all{der Unterricht beginnt} (les cours commencent).
\item Laquelle arrive en premier ? La fermeture des écoles. Elle est exprimée par \all{hatten geschlossen}.
\item Laquelle arrive ensuite ? Le début des cours en ligne, exprimé par \all{begann} (imparfait, \emph{Präteritum}).
\end{enumerate}

L'allemand, comme le français, dispose donc d'un temps spécial pour dire qu'une action passée est \emph{antérieure} à une autre action passée : le \cle{Plusquamperfekt}, c'est-à-dire le \textbf{plus-que-parfait}.

\begin{popfigure}
\begin{tikzpicture}[scale=1]
\draw[-{Stealth},thick,popblue] (0,0) -- (11.5,0) node[below left=0.1cm and 0cm,popdark]{Zeit (le temps)};
\fill[poporange] (2,0) circle (0.12);
\node[above=0.25cm,align=center,text=poporange] at (2,0) {\textbf{1. Schulen schließen}\\ \footnotesize\all{hatten geschlossen}\\ \footnotesize Plusquamperfekt};
\fill[popgreen] (6,0) circle (0.12);
\node[above=0.25cm,align=center,text=popgreen] at (6,0) {\textbf{2. Unterricht beginnt}\\ \footnotesize\all{begann}\\ \footnotesize Präteritum};
\fill[popink] (10,0) circle (0.12);
\node[above=0.25cm,align=center,text=popink] at (10,0) {\textbf{maintenant}\\ \footnotesize (aujourd'hui)};
\draw[-{Stealth},thick,poporange] (2.6,-0.6) -- (5.4,-0.6) node[midway,below,popdark]{\footnotesize antériorité : l'action 1 est \emph{avant} l'action 2};
\end{tikzpicture}
\legende{Frise des temps : le Plusquamperfekt situe une action avant une autre action passée.}
\end{popfigure}

\courssub{Formation du Plusquamperfekt}

\begin{propriete}[Formation du Plusquamperfekt]
Le \cle{Plusquamperfekt} se forme avec :
\begin{center}
\textbf{l'imparfait (Präteritum) de \all{haben} ou \all{sein} + le participe passé du verbe}
\end{center}
\begin{itemize}
\item \all{ich hatte gelernt} — j'avais appris
\item \all{ich war gefahren} — j'étais allé(e) (en véhicule)
\end{itemize}
C'est exactement la même construction qu'en français : \emph{j'avais appris} = imparfait de « avoir » + participe passé.
\end{propriete}

Le choix de l'auxiliaire suit \textbf{les mêmes règles qu'au Perfekt}, que vous connaissez depuis la classe de Première :

\begin{aretenir}[Choix de l'auxiliaire : haben ou sein ?]
\begin{itemize}
\item \all{sein} (\all{ich war, du warst, er/sie/es war, wir waren, ihr wart, sie/Sie waren}) avec les verbes de \textbf{mouvement} (\all{gehen, fahren, kommen, fliegen, aufstehen}) et de \textbf{changement d'état} (\all{sterben, einschlafen, aufwachen, passieren}), ainsi que \all{sein} et \all{bleiben}.
\item \all{haben} (\all{ich hatte, du hattest, er/sie/es hatte, wir hatten, ihr hattet, sie/Sie hatten}) avec \textbf{tous les autres verbes} : transitifs, réfléchis, verbes d'activité ou d'état (\all{lernen, arbeiten, schließen, vergessen, beginnen}).
\end{itemize}
\end{aretenir}

\begin{center}
\begin{tabular}{|l|l|l|}
\hline
\rowcolor{popblueL} Personne & \all{lernen} (haben) & \all{gehen} (sein) \\
\hline
ich & hatte gelernt & war gegangen \\
\hline
du & hattest gelernt & warst gegangen \\
\hline
er/sie/es & hatte gelernt & war gegangen \\
\hline
wir & hatten gelernt & waren gegangen \\
\hline
ihr & hattet gelernt & wart gegangen \\
\hline
sie/Sie & hatten gelernt & waren gegangen \\
\hline
\end{tabular}
\end{center}

\begin{popfigure}
\begin{tikzpicture}[
  box/.style={draw=popblue,thick,rounded corners,fill=popblueL,align=center,minimum width=3.4cm,minimum height=0.9cm},
  aux/.style={draw=poporange,thick,rounded corners,fill=poporangeL,align=center,minimum width=2.6cm,minimum height=0.9cm}]
\node[aux, align=center] (hab) at (0,2){\all{hatte}\\ \footnotesize (haben)};
\node[aux, align=center] (war) at (0,0){\all{war}\\ \footnotesize (sein)};
\node[box, align=center] (part) at (5.5,1){\textbf{+ participe passé}\\ \all{gelernt / geschlossen}\\ \all{gegangen / gefahren}};
\node[box,fill=poppurpleL,draw=poppurple, align=center] (res) at (10.5,1){\textbf{Plusquamperfekt}\\ \all{ich hatte gelernt}\\ \all{ich war gegangen}};
\draw[-{Stealth},thick,popdark] (hab.east) -- (part.west);
\draw[-{Stealth},thick,popdark] (war.east) -- (part.west);
\draw[-{Stealth},thick,popdark] (part.east) -- (res.west);
\end{tikzpicture}
\legende{Schéma de formation : imparfait de l'auxiliaire + participe passé.}
\end{popfigure}

\courssub{Emploi : l'antériorité}

\begin{propriete}[Emploi du Plusquamperfekt]
Le Plusquamperfekt exprime une action \textbf{achevée AVANT} une autre action passée. L'action de référence est, elle, au \emph{Präteritum} (ou au Perfekt à l'oral). En français, on emploie de la même façon le plus-que-parfait : « Quand j'\emph{avais} fini mes devoirs, je me suis connecté. »
\end{propriete}

\begin{exemplebox}[Exemples traduits]
\begin{itemize}
\item \all{Nachdem ich drei Monate online gelernt hatte, merkte ich, dass mir der Kontakt fehlte.} — « Après avoir appris en ligne pendant trois mois, je remarquai que le contact me manquait. »
\item \all{Als der Strom ausgefallen war, konnten wir nicht mehr lernen.} — « Quand le courant eut été coupé, nous ne pûmes plus étudier. »
\item \all{Bevor der Kurs begann, hatte ich das Passwort vergessen.} — « Avant que le cours commence, j'avais oublié le mot de passe. »
\item \all{Die Lehrerin hatte die Aufgaben geschickt, bevor die Verbindung abbrach.} — « La professeure avait envoyé les devoirs avant que la connexion ne coupe. »
\end{itemize}
\end{exemplebox}

\courssub{Les propositions temporelles et la concordance des temps}

Le plus-que-parfait apparaît très souvent dans les \textbf{subordonnées temporelles}. Voici les quatre conjonctions essentielles du programme de Terminale :

\begin{center}
\begin{tabularx}{\linewidth}{|l|X|X|}
\hline
\rowcolor{popblueL} Conjonction & Sens et temps & Exemple \\
\hline
\all{nachdem} & après que ; antériorité : \textbf{+ Plusquamperfekt} si la principale est au passé & \all{Nachdem die Schulen geschlossen hatten, begann der Unterricht.} \\
\hline
\all{als} & quand, lorsque ; action \textbf{unique} au passé : + Präteritum & \all{Als der Strom ausfiel, konnten wir nicht lernen.} \\
\hline
\all{während} & pendant que ; \textbf{simultanéité} : mêmes temps des deux côtés & \all{Während ich lernte, schlief mein Bruder.} \\
\hline
\all{bevor / ehe} & avant que ; action \textbf{postérieure} & \all{Bevor der Kurs begann, prüfte ich das Netz.} \\
\hline
\end{tabularx}
\end{center}

\begin{propriete}[Règle de concordance avec nachdem]
\begin{itemize}
\item \all{nachdem} + \textbf{Plusquamperfekt} $\rightarrow$ principale au \textbf{Präteritum} : \all{Nachdem ich gegessen hatte, ging ich an den Computer.} (« Après avoir mangé, j'allai à l'ordinateur. »)
\item \all{nachdem} + \textbf{Perfekt} $\rightarrow$ principale au \textbf{présent} : \all{Nachdem ich gegessen habe, gehe ich an den Computer.} (« Après avoir mangé, je vais à l'ordinateur. »)
\end{itemize}
Autrement dit : la subordonnée en \all{nachdem} a toujours « un temps d'avance » sur la principale.
\end{propriete}

Comparons les deux langues :

\begin{center}
\begin{tabularx}{\linewidth}{|X|X|}
\hline
\rowcolor{popblueL} Français & Deutsch \\
\hline
Après que les écoles eurent fermé, les cours commencèrent. & \all{Nachdem die Schulen geschlossen hatten, begann der Unterricht.} \\
\hline
Quand j'avais fini mes devoirs, je me connectais. & \all{Nachdem ich meine Hausaufgaben gemacht hatte, loggte ich mich ein.} \\
\hline
Pendant que je travaillais, la connexion coupait souvent. & \all{Während ich arbeitete, brach die Verbindung oft ab.} \\
\hline
\end{tabularx}
\end{center}

\begin{attention}[Piège n° 1 : jamais le Präteritum après nachdem au passé]
Ne dites jamais \all{nachdem die Schulen schlossen} ✗. Avec \all{nachdem}, quand la principale est au passé, la subordonnée est \textbf{obligatoirement} au Plusquamperfekt : \all{Nachdem die Schulen geschlossen hatten, begann der Unterricht.} ✓ Le Präteritum simple après \all{nachdem} est une faute classique de francophone (et même d'Allemand pressé !).
\end{attention}

\begin{attention}[Piège n° 2 : la place des verbes]
Deux règles d'ordre des mots à ne jamais oublier :
\begin{enumerate}
\item Dans la subordonnée temporelle, le verbe conjugué va \textbf{à la fin} : \all{Nachdem ich drei Monate online gelernt hatte, \dots} (et non \all{Nachdem ich hatte gelernt} ✗). Le participe passé se place juste avant l'auxiliaire.
\item Quand la subordonnée ouvre la phrase, la principale commence par \textbf{son verbe conjugué} : \all{\dots, begann der Unterricht.} (et non \all{\dots, der Unterricht begann} ✗). La virgule joue le rôle de « pivot » : verbe + verbe.
\end{enumerate}
\end{attention}

\courssub{Texte d'entraînement}

\begin{textebox}[Ein schwieriger Morgen]
\all{Nachdem Amadou aufgestanden war, schaltete er sofort seinen Laptop ein. Er hatte am Abend vorher seine Hausaufgaben fertig gemacht, und er wollte sie der Lehrerin schicken. Aber als er sich einloggen wollte, merkte er, dass das Netz nicht funktionierte. Während er auf die Verbindung wartete, frühstückte er schnell. Bevor der Online-Kurs begann, hatte er schon dreimal das Passwort neu eingegeben. Nachdem der Strom ausgefallen war, konnte niemand in seiner Familie mehr arbeiten. Seine Schwester, die die Nacht vorher bis Mitternacht gelernt hatte, war sehr enttäuscht. Als der Strom endlich zurückkam, war der Kurs schon fast zu Ende. Amadou ärgerte sich, aber dann dachte er: „Morgen stehe ich früher auf und prüfe zuerst das Netz!“}

\emph{Glossaire :} der Laptop, -s (l'ordinateur portable) ; einloggen (sich) (se connecter) ; das Netz, -e (le réseau) ; die Verbindung, -en (la connexion) ; das Frühstück (le petit-déjeuner) : frühstücken (prendre le petit-déjeuner) ; der Strom (le courant électrique) ; ausfallen (fiel aus, ist ausgefallen) (tomber en panne, être coupé) ; enttäuscht (déçu) ; sich ärgern (s'énerver) ; prüfen (vérifier).
\end{textebox}

\textbf{Questions de compréhension :}
\begin{enumerate}
\item \all{Was hatte Amadou am Abend vorher gemacht?}
\item \all{Was funktionierte nicht, als er sich einloggen wollte?}
\item \all{Was machte er, während er wartete?}
\item \all{Warum war seine Schwester enttäuscht?}
\end{enumerate}

\courssub{Exercice résolu et entraînement}

\begin{exoresolu}[Mettez les verbes au temps correct]
Énoncé : conjuguez les verbes entre parenthèses (subordonnée : Plusquamperfekt ou Präteritum selon la conjonction ; principale : Präteritum).

\begin{enumerate}
\item \all{Nachdem er das Passwort (eingeben), (funktionieren) das Netz nicht.}
\item \all{Als der Strom (ausfallen), (sitzen) wir vor dem Laptop.}
\item \all{Während sie (lernen), (schlafen) ihr Bruder.}
\item \all{Bevor der Kurs (beginnen), (vergessen – er) das Passwort.}
\end{enumerate}

\tcblower

Solution détaillée :
\begin{enumerate}
\item \all{Nachdem er das Passwort \textbf{eingegeben hatte}, \textbf{funktionierte} das Netz nicht.} — \all{nachdem} + principale au passé $\rightarrow$ Plusquamperfekt dans la subordonnée (\all{haben} + \all{eingegeben}), verbe conjugué à la fin, la principale commence par le verbe.
\item \all{Als der Strom \textbf{ausfiel}, \textbf{saßen} wir vor dem Laptop.} — \all{als} + action unique passée $\rightarrow$ Präteritum des deux côtés ; \all{ausfallen} est un verbe fort : \all{fiel aus}.
\item \all{Während sie \textbf{lernte}, \textbf{schlief} ihr Bruder.} — \all{während} = simultanéité : Präteritum des deux côtés.
\item \all{Bevor der Kurs \textbf{begann}, \textbf{hatte er} das Passwort \textbf{vergessen}.} — l'oubli est antérieur au début du cours : Plusquamperfekt ; \all{vergessen} se conjugue avec \all{haben}.
\end{enumerate}
\end{exoresolu}

\begin{exercice}[À vous de jouer]
Traduisez en allemand :
\begin{enumerate}
\item Après que la connexion eut été coupée, nous ne pûmes plus suivre le cours.
\item Quand j'avais fini mes devoirs, j'éteignis l'ordinateur.
\item Avant que le cours commence, elle avait vérifié le réseau.
\end{enumerate}
\end{exercice}

\begin{corrige}[Exercice « À vous de jouer »]
\begin{enumerate}
\item \all{Nachdem die Verbindung abgebrochen war, konnten wir den Kurs nicht mehr verfolgen.}
\item \all{Nachdem ich meine Hausaufgaben gemacht hatte, schaltete ich den Computer aus.}
\item \all{Bevor der Kurs begann, hatte sie das Netz geprüft.}
\end{enumerate}
\end{corrige}

\courssub{Méthode du commentaire et production guidée}

\begin{methode}[Méthode : Commenter un texte et produire un paragraphe argumenté]
Pour réussir un commentaire de texte ou une production écrite sur le thème du numérique (Bac LV2), suivez ces étapes :
\begin{enumerate}
\item \textbf{Lecture globale :} Identifiez le thème (ici : \all{Fernunterricht}), la source, le ton (journalistique, témoignage) et la thèse principale.
\item \textbf{Analyse structurelle :} Repérez les deux grands axes : \textbf{Opposition} (\all{einerseits \dots\ andererseits}, \all{Vorteil}/\all{Nachteil}) et \textbf{Chronologie} (connecteurs temporels : \all{nachdem, als, während, bevor} + \textbf{Plusquamperfekt} / \textbf{Präteritum}).
\item \textbf{Lexique actif :} Listez le vocabulaire thématique avec articles et pluriels (\all{der Strom, die Verbindung, das Netz, der Vorteil, die Nachteile, die Herausforderung}).
\item \textbf{Grammaire en action :} Vérifiez l'ordre des mots (verbe en position finale en subordonnée, inversion verbe-sujet en principale après subordonnée initiale) et la concordance des temps (\all{nachdem} + Plusquamperfekt $\rightarrow$ Präteritum).
\item \textbf{Rédaction guidée :} Rédigez un paragraphe de 8–10 lignes structuré ainsi :
      \begin{itemize}
      \item Phrase d'accroche (contexte : \all{Seit der Pandemie \dots}).
      \item Argument 1 (avantages) avec \all{einerseits} + exemple personnel.
      \item Argument 2 (inconvénients) avec \all{andererseits} + cause technique (\all{Stromausfall, schlechtes Netz}).
      \item Synthèse avec \all{bevor / nachdem} pour situer une action antérieure (\all{Bevor ich online gehe, prüfe ich die Verbindung.} / \all{Nachdem der Strom ausgefallen war, \dots}).
      \end{itemize}
\item \textbf{Relecture :} Contrôlez majuscules (noms), Umlaute/ß, cas (accusatif/datif après prépositions), particule séparable en fin de proposition.
\end{enumerate}
\end{methode}

\begin{experience}[Production orale guidée : « Mon matin numérique »]
En binôme : l'élève A raconte une matinée de cours à distance difficile (panne de courant, mot de passe oublié, connexion lente) en utilisant \textbf{au moins 3 subordonnées temporelles} (\all{nachdem, als, während, bevor}) et le \textbf{Plusquamperfekt}. L'élève B note les temps employés et vérifie l'ordre des mots. Échangez les rôles. Durée : 5 min par élève.
\end{experience}

\begin{savaistu}[Plusquamperfekt à l'oral et au Bac]
À l'oral, beaucoup d'Allemands remplacent le Plusquamperfekt par le Perfekt après \all{nachdem} : on entend souvent \all{Nachdem ich gegessen habe, bin ich gegangen}. C'est toléré dans la conversation, mais \textbf{à l'écrit — et surtout au Bac — la concordance correcte est exigée} : \all{nachdem} + Plusquamperfekt quand la principale est au passé. Sois donc plus royaliste que le roi : à l'écrit, respecte toujours la règle !
\end{savaistu}

\begin{aretenir}[L'essentiel de la section]
\begin{itemize}
\item Plusquamperfekt = \all{hatte/war} + participe passé ; il exprime une action achevée \textbf{avant} une autre action passée.
\item Auxiliaire : \all{sein} (mouvement, changement d'état) ou \all{haben} (tous les autres verbes), comme au Perfekt.
\item \all{nachdem} + Plusquamperfekt $\rightarrow$ principale au Präteritum ; \all{als} + Präteritum (action unique) ; \all{während} = simultanéité ; \all{bevor/ehe} = action postérieure.
\item Ordre des mots : verbe conjugué à la fin de la subordonnée ; la principale qui suit commence par son verbe.
\end{itemize}
\end{aretenir}

\coursec{Méthode du commentaire et production guidée}

Dans la section précédente, nous avons appris à situer les événements dans le temps grâce au plus-que-parfait (\all{Nachdem der Schüler drei Monate online gelernt hatte, verstand er, dass die Schule wichtig ist.}). Ces outils grammaticaux — le plus-que-parfait, les propositions temporelles, l'opposition avec \all{einerseits … andererseits} — ne sont pas des exercices isolés : ce sont précisément les instruments dont vous aurez besoin pour rédiger un \cle{commentaire de texte} au baccalauréat et pour produire un résumé ou un paragraphe d'opinion. C'est l'objet de cette dernière section : passer de la compréhension à la \cle{production écrite et orale}.

\courssub{La méthode du commentaire de texte au baccalauréat}

\begin{definition}[Qu'est-ce qu'un commentaire de texte ?]
Le \cle{commentaire de texte} (\all{die Textanalyse} ou \all{der Kommentar}) est un exercice écrit dans lequel le candidat présente un texte, en dégage les idées principales, analyse les moyens d'expression de l'auteur et donne son \cle{avis personnel argumenté}. Au baccalauréat, il s'agit de montrer que vous avez compris le texte ET que vous savez en parler avec vos propres mots, en allemand correct.
\end{definition}

\begin{methode}[Les cinq étapes du commentaire]
\begin{enumerate}
\item \textbf{L'introduction} (\all{die Einleitung}) : présenter le texte en deux ou trois phrases — le thème, le type de texte (article, récit, lettre), la \cle{problématique} (la question que le texte pose).
\item \textbf{Le résumé structuré} (\all{die Zusammenfassung}) : reformuler les idées principales dans l'ordre du texte, au présent, sans recopier les phrases.
\item \textbf{L'analyse} (\all{die Analyse}) : montrer COMMENT l'auteur exprime ses idées — le lexique choisi, les procédés (opposition, exemples, chiffres), la structure du texte.
\item \textbf{L'avis personnel} (\all{die eigene Meinung}) : donner votre opinion avec des arguments (\all{Meiner Meinung nach …}, \all{Ich finde, dass …}).
\item \textbf{La conclusion} (\all{der Schluss}) : une phrase de bilan qui répond à la problématique.
\end{enumerate}
\end{methode}

\begin{popfigure}
\begin{center}
\begin{tikzpicture}[
  node distance=0.55cm,
  etape/.style={rectangle, rounded corners=4pt, draw=popblue, fill=popblueL, thick, align=center, text width=4.6cm, minimum height=0.9cm, font=\small},
  note/.style={rectangle, draw=poporange, fill=poporangeL, rounded corners=4pt, align=center, text width=4.2cm, font=\scriptsize},
  fleche/.style={-{Stealth[length=3mm]}, thick, popdark}
]
\node[etape, align=center] (e1){\textbf{1. Einleitung}\\ Thema + Problemfrage};
\node[etape, below=of e1, align=center] (e2){\textbf{2. Zusammenfassung}\\ Hauptideen, eigene Worte};
\node[etape, below=of e2, align=center] (e3){\textbf{3. Analyse}\\ Wortschatz, Stilmittel};
\node[etape, below=of e3, align=center] (e4){\textbf{4. Eigene Meinung}\\ Meinung + Argumente};
\node[etape, below=of e4, align=center] (e5){\textbf{5. Schluss}\\ Antwort auf die Problemfrage};
\draw[fleche] (e1) -- (e2);
\draw[fleche] (e2) -- (e3);
\draw[fleche] (e3) -- (e4);
\draw[fleche] (e4) -- (e5);
\node[note, right=1.1cm of e2, align=center] (n2){Präsens verwenden\\ nicht abschreiben!};
\node[note, right=1.1cm of e4, align=center] (n4){\all{Meiner Meinung nach …}\\ \all{Ich bin der Meinung, dass …}};
\draw[fleche, poporange, dashed] (n2) -- (e2);
\draw[fleche, poporange, dashed] (n4) -- (e4);
\end{tikzpicture}
\end{center}
\legende{Organigramme des cinq étapes du commentaire de texte}
\end{popfigure}

\begin{attention}[Les trois interdits du commentaire]
\begin{itemize}
\item \textbf{Ne jamais recopier le texte.} Le correcteur connaît le texte par cœur : recopier prouve que vous n'avez rien compris. Il faut \cle{reformuler} avec vos propres mots.
\item On peut \cle{citer} de courts extraits (deux à cinq mots) entre guillemets allemands : \all{Der Autor spricht von „der schlechten Verbindung“.}
\item \textbf{Toujours donner un avis personnel justifié} : une opinion sans argument ne vaut rien. Utilisez \all{weil}, \all{denn}, \all{deshalb} pour justifier.
\end{itemize}
\end{attention}

\courssub{La banque de phrases du commentaire}

Pour éviter les répétitions et gagner en correction, voici les expressions indispensables, classées par étape du commentaire.

\begin{center}
\begin{tabularx}{\linewidth}{|X|X|X|}
\hline
\rowcolor{popblueL} \textbf{Étape} & \textbf{Deutsch} & \textbf{Français} \\
\hline
Introduction & Der Text handelt von … & Le texte traite de … \\
\hline
Introduction & Es geht um die Frage, ob … & Il s'agit de la question de savoir si … \\
\hline
Résumé & Der Autor zeigt, dass … & L'auteur montre que … \\
\hline
Résumé & Er beschreibt die Situation … & Il décrit la situation … \\
\hline
Analyse & Der Autor betont die Probleme. & L'auteur souligne les problèmes. \\
\hline
Analyse & Er benutzt die Gegenüberstellung \all{einerseits … andererseits}. & Il utilise l'opposition « d'une part … d'autre part ». \\
\hline
Opinion & Meiner Meinung nach … & À mon avis … \\
\hline
Opinion & Ich bin der Meinung, dass … & Je suis d'avis que … \\
\hline
Conclusion & Zusammenfassend kann man sagen, dass … & En résumé, on peut dire que … \\
\hline
\end{tabularx}
\end{center}

\begin{exemplebox}[Phrases modèles traduites]
\all{Der Text handelt von den Vor- und Nachteilen des Fernunterrichts.} « Le texte traite des avantages et des inconvénients des cours à distance. »

\all{Meiner Meinung nach kann der Fernunterricht die Schule nicht ersetzen.} « À mon avis, les cours à distance ne peuvent pas remplacer l'école. »

\all{Zusammenfassend kann man sagen, dass das Digitale Chancen und Risiken hat.} « En résumé, on peut dire que le numérique présente des chances et des risques. »
\end{exemplebox}

\begin{attention}[Variez les verbes de parole !]
L'erreur la plus fréquente des candidats francophones est de répéter \all{der Text sagt} à chaque phrase. Un texte ne « dit » rien : il \textbf{montre}, \textbf{explique}, \textbf{décrit}, \textbf{souligne}. Variez : \all{zeigt} (montre), \all{erklärt} (explique), \all{beschreibt} (décrit), \all{betont} (souligne), \all{nennt} (cite), \all{kritisiert} (critique). Autre piège : après \all{Meiner Meinung nach}, le verbe est en deuxième position : \all{Meiner Meinung nach \textbf{kann} der Fernunterricht …} et non \all{Meiner Meinung nach der Fernunterricht kann …}.
\end{attention}

\courssub{Un commentaire modèle}

\begin{textebox}[Début de commentaire modèle]
Der Text beschreibt die Erfahrung eines Schülers mit dem Fernunterricht während der Corona-Pandemie. Einerseits zeigt er die Vorteile: Man spart Zeit und Geld. Andererseits nennt er viele Probleme: die schlechte Verbindung, den Stromausfall und den fehlenden Kontakt zu den Freunden. Nachdem der Schüler drei Monate online gelernt hatte, verstand er, dass die Schule wichtig ist. Meiner Meinung nach ist der Fernunterricht eine gute Lösung in einer Krise, aber er kann den normalen Unterricht nicht ersetzen.
\end{textebox}

\textbf{Glossaire :} \all{die Erfahrung, -en} (l'expérience) ; \all{der Stromausfall, die Stromausfälle} (la coupure de courant) ; \all{fehlend} (manquant) ; \all{die Krise, -n} (la crise) ; \all{ersetzen} (remplacer).

\textbf{Analyse des connecteurs du modèle :}
\begin{itemize}
\item \all{Einerseits … andererseits} : structure d'opposition qui organise le résumé (avantages / problèmes) — c'est la grammaire de la section 4 mise en pratique.
\item \all{Nachdem der Schüler drei Monate online gelernt hatte, verstand er …} : proposition temporelle au plus-que-parfait (antériorité) — c'est la grammaire de la section 5.
\item \all{Meiner Meinung nach} : introduit l'avis personnel, avec le verbe en position 2.
\item Les verbes de parole variés : \all{beschreibt}, \all{zeigt}, \all{nennt}.
\end{itemize}

\courssub{Production guidée 1 : le résumé en quatre phrases}

\begin{exoresolu}[Résumé du texte support]
Énoncé : Résumez le texte \all{Fernunterricht – Chancen und Probleme} en quatre phrases. Consignes obligatoires : utilisez au moins une fois \all{nachdem} avec le plus-que-parfait et une fois la structure \all{einerseits … andererseits}. N'utilisez pas les mêmes phrases que le texte.
\tcblower
Solution détaillée :

\all{Der Text handelt von einem Schüler, der während der Pandemie zu Hause lernen musste. Einerseits war der Fernunterricht praktisch, weil er Zeit und Geld sparte; andererseits gab es große Probleme mit dem Strom und dem Internet. Nachdem der Schüler mehrere Wochen allein vor dem Computer gesessen hatte, vermisste er seine Freunde und Lehrer. Zusammenfassend kann man sagen, dass der Fernunterricht nützlich, aber kein Ersatz für die Schule ist.}

Méthode : (1) la première phrase présente le thème (\all{Der Text handelt von …}) ; (2) la deuxième oppose avantages et inconvénients avec la structure imposée ; (3) la troisième exprime l'antériorité avec \all{nachdem} + plus-que-parfait et le verbe principal en position 2 ; (4) la dernière conclut avec \all{zusammenfassend}.
\end{exoresolu}

\courssub{Production guidée 2 : le paragraphe d'opinion}

\begin{exercice}[Fernunterricht in Kamerun : Vorteile und Nachteile]
Rédigez un paragraphe de 8 à 10 lignes sur le thème : \all{Fernunterricht in Kamerun : Vorteile und Nachteile}. Respectez le plan imposé :
\begin{enumerate}
\item \textbf{Einleitung} (1-2 phrases) : présentez le sujet avec \all{Es geht um die Frage, ob …} ou \all{In Kamerun ist der Fernunterricht …}.
\item \textbf{Vorteile} (2-3 phrases) : commencez par \all{Einerseits …} et donnez deux avantages avec \all{weil} ou \all{denn}.
\item \textbf{Nachteile} (2-3 phrases) : commencez par \all{Andererseits …} et citez deux problèmes concrets du contexte camerounais (\all{der Stromausfall}, \all{die schlechte Verbindung}, \all{das teure Internet}).
\item \textbf{Meinung + Schluss} (2 phrases) : \all{Meiner Meinung nach …} puis \all{Zusammenfassend kann man sagen, dass …}.
\end{enumerate}
\end{exercice}

\begin{corrige}[Production guidée 2 — modèle de rédaction]
\all{In Kamerun ist der Fernunterricht seit der Corona-Pandemie ein wichtiges Thema. Es geht um die Frage, ob er eine gute Lösung für unsere Schüler ist. Einerseits hat der Fernunterricht klare Vorteile: Die Schüler müssen nicht zur Schule fahren, und sie sparen Geld für den Transport. Auch kranke Schüler können weiter lernen. Andererseits gibt es in Kamerun große Probleme: Der Strom fällt oft aus, die Verbindung ist schlecht und das Internet ist für viele Familien zu teuer. Außerdem vermissen die Schüler den Kontakt zu ihren Freunden und Lehrern. Meiner Meinung nach kann der Fernunterricht die Schule nicht ersetzen, weil das Lernen in der Gruppe wichtiger ist. Zusammenfassend kann man sagen, dass der Fernunterricht eine gute Hilfe in einer Krise ist, aber keine dauerhafte Lösung.}

Points forts de ce modèle : plan respecté, opposition \all{einerseits … andererseits}, arguments justifiés avec \all{weil}, vocabulaire du thème (\all{der Stromausfall}, \all{die Verbindung}), verbes toujours en position correcte.
\end{corrige}

\courssub{Production orale : le débat en binômes}

\begin{experience}[Befürworter oder Gegner des Fernunterrichts ?]
Par binômes, préparez un débat de quatre minutes. L'élève A est \all{Befürworter} (partisan) des cours à distance, l'élève B est \all{Gegner} (adversaire).
\begin{enumerate}
\item Chacun note trois arguments avec le vocabulaire du chapitre (\all{der Vorteil}, \all{der Nachteil}, \all{die Verbindung}, \all{das Netz}, \all{der Strom}).
\item Le débat : chaque élève présente sa position en une minute, puis répond aux arguments de l'autre.
\item \textbf{Structure d'opposition obligatoire} : chaque élève doit utiliser au moins une fois \all{Einerseits … andererseits …} et une fois \all{Meiner Meinung nach …}.
\item La classe vote : qui a le mieux argumenté ?
\end{enumerate}
Exemples de répliques : \all{Einerseits sparst du Zeit, andererseits lernst du allein nicht so gut.} — \all{Ich bin anderer Meinung, denn …} (Je ne suis pas de ton avis, car …) — \all{Das stimmt, aber …} (C'est vrai, mais …).
\end{experience}

\begin{aretenir}[L'essentiel de la section]
\begin{itemize}
\item Un commentaire suit cinq étapes : \all{Einleitung → Zusammenfassung → Analyse → eigene Meinung → Schluss}.
\item On reformule, on ne recopie jamais ; on cite seulement de courts extraits entre „ … “.
\item On varie les verbes de parole : \all{zeigt, erklärt, beschreibt, betont, nennt}.
\item Les outils grammaticaux de la leçon — \all{einerseits … andererseits} et \all{nachdem} + plus-que-parfait — sont des connecteurs de rédaction à réemployer dans toutes vos productions écrites et orales.
\end{itemize}
\end{aretenir}

\newpage

\coursec{Activité d'intégration}

\coursec{Activité d'intégration : Fernunterricht – Fluch oder Segen?}

\courssub{Objectifs de l'activité}

À l'issue de cette activité d'intégration, l'élève sera capable de :
\begin{itemize}
\item mobiliser dans un échange oral argumenté le vocabulaire thématique des cours à distance et du numérique (\cle{der Fernunterricht}, \cle{die Verbindung}, \cle{das Netz}, \cle{der Strom}, \cle{der Vorteil}, \cle{der Nachteil}...) ;
\item construire une argumentation équilibrée grâce à la structure d'opposition \all{einerseits ... andererseits} ;
\item employer correctement le plus-que-parfait dans une proposition temporelle introduite par \all{nachdem} pour situer les faits dans le passé (la fermeture des écoles, la coupure de courant...) ;
\item rédiger individuellement un court article de presse scolaire (10 à 12 lignes) avec introduction, arguments pour et contre, et conclusion personnelle ;
\item prendre la parole en public, écouter et réagir aux arguments d'autrui dans un débat réglé.
\end{itemize}

\courssub{Situation}

La rédaction du journal de ton lycée à Yaoundé, \emph{Die Stimme der Jugend} (La Voix de la jeunesse), prépare son prochain numéro consacré aux nouvelles technologies. La rédactrice en chef, une élève de Terminale A, lance un appel à contributions après la lecture d'une lettre de lecteur parue dans le numéro précédent. Le thème est d'actualité : depuis quelques années, les cours à distance se sont développés au Cameroun, d'abord pendant la fermeture des écoles, puis avec les plateformes numériques du MINESEC et les cours en ligne des universités. Mais les problèmes restent nombreux : coupures de courant, connexion lente, coût des données, manque d'ordinateurs dans les familles.

Voici la lettre qui a déclenché le débat :

\begin{textebox}[Brief eines Lesers an die Schülerzeitung]
Liebe Redaktion,\\
ich heiße Arnaud und besuche die Terminale in Douala. Während der Schulschließung hatten wir drei Monate lang Fernunterricht per Internet und Radio. Einerseits war das eine gute Lösung: Wir konnten zu Hause weiterlernen und haben das Schuljahr nicht verloren. Andererseits gab es große Probleme: Die Verbindung war oft schlecht, der Strom fiel manchmal stundenlang aus, und mein kleiner Bruder hatte kein Handy. Viele meiner Freunde konnten nicht am Unterricht teilnehmen, weil ihre Eltern kein Geld für Internetdaten hatten. Nachdem die Schulen wieder geöffnet worden waren, mussten wir sehr viel Stoff wiederholen. Meine Frage an euch: Ist der Fernunterricht für Kamerun ein Fluch oder ein Segen?\\
Mit freundlichen Grüßen\\
Arnaud
\end{textebox}

\textbf{Glossaire :} die Schulschließung, -en (la fermeture des écoles) ; der Stoff, -e (le programme, la matière) ; wiederholen (répéter, revoir) ; der Fluch, \emph{pl.} die Flüche (la malédiction) ; der Segen, - (la bénédiction) ; teilnehmen an + \emph{Dat.} (participer à) ; die Internetdaten, \emph{pl.} (les données Internet).

Ton rôle : tu participes au grand débat de la classe, puis tu rédiges un article pour le journal du lycée afin de répondre à la question d'Arnaud.

\courssub{Consignes}

\textbf{Phase 1 — Préparation du débat (travail en binômes, 15 minutes).}

\begin{enumerate}
\item Relis la lettre d'Arnaud. Souligne en rouge les mots du champ lexical du numérique, en bleu la structure d'opposition et en vert la proposition temporelle au plus-que-parfait.
\item La classe est divisée en deux camps : le camp A défend les cours à distance (\all{der Fernunterricht ist ein Segen}), le camp B les critique (\all{der Fernunterricht ist ein Fluch}). Avec ton binôme, note \textbf{trois arguments} pour ton camp. Chaque argument doit contenir au moins un mot du lexique thématique de la leçon.
\item Formule au moins \textbf{deux phrases avec \all{einerseits ... andererseits}} pour montrer que tu vois les deux faces du problème (c'est ce qui rend une argumentation crédible).
\item Rédige au moins \textbf{une phrase au plus-que-parfait} avec \all{nachdem} pour raconter un fait antérieur, par exemple : \all{Nachdem der Präsenzunterricht ausgefallen war, begann der Fernunterricht.} (« Après que le cours en présentiel eut été interrompu, les cours à distance commencèrent. »)
\end{enumerate}

\textbf{Phase 2 — Le débat (10 minutes, classe entière).}

\begin{enumerate}
\setcounter{enumi}{4}
\item L'enseignant modère le débat. Chaque camp présente ses arguments à tour de rôle. Écoute attentivement le camp adverse et note un argument que tu pourras réutiliser ou réfuter dans ton article.
\item Réagis aux arguments de l'autre camp avec des formules polies : \all{Ich bin anderer Meinung, weil ...} (« Je suis d'un autre avis parce que... »), \all{Das stimmt, aber ...} (« C'est vrai, mais... »), \all{Man darf nicht vergessen, dass ...} (« Il ne faut pas oublier que... »).
\end{enumerate}

\textbf{Phase 3 — Rédaction individuelle de l'article (20 minutes).}

\begin{enumerate}
\setcounter{enumi}{6}
\item Rédige un article de \textbf{10 à 12 lignes} intitulé \all{Fernunterricht : eine Chance für Kamerun ?} pour le journal du lycée. Plan obligatoire :
\begin{itemize}
\item une \textbf{introduction} qui présente le sujet et le contexte camerounais ;
\item \textbf{deux arguments pour} les cours à distance ;
\item \textbf{deux arguments contre} (utilise \all{einerseits ... andererseits} au moins une fois) ;
\item une \textbf{conclusion personnelle} avec ton avis justifié.
\end{itemize}
\item Vérifie ton texte : majuscule à tous les noms, verbe conjugué en deuxième position, particule séparable en fin de phrase, au moins un plus-que-parfait correctement formé (\all{hatte / war} + participe passé).
\end{enumerate}

\courssub{Ressources à mobiliser}

\begin{aretenir}[Boîte à outils linguistique]
\textbf{Lexique :} \cle{der Fernunterricht} (l'enseignement à distance) ; \cle{die Verbindung}, -en (la connexion) ; \cle{das Netz}, -e (le réseau) ; \cle{der Strom} (l'électricité, le courant) ; \cle{der Stromausfall}, die Stromausfälle (la coupure de courant) ; \cle{der Vorteil}, -e (l'avantage) ; \cle{der Nachteil}, -e (l'inconvénient) ; \cle{die Lösung}, -en (la solution) ; \cle{das Handy}, -s (le téléphone portable) ; \cle{teilnehmen} (nimmt teil, nahm teil, hat teilgenommen) an + \emph{Dat.}.

\textbf{Grammaire 1 — l'opposition :} \all{einerseits} + phrase 1, \all{andererseits} + phrase 2. Les deux éléments occupent la première place, donc le verbe suit immédiatement : \all{Einerseits spart man Zeit, andererseits braucht man eine gute Verbindung.}

\textbf{Grammaire 2 — le plus-que-parfait :} auxiliaire \all{haben} ou \all{sein} au \emph{Präteritum} (\all{hatte}, \all{war}) + participe passé. Avec \all{nachdem}, le plus-que-parfait dans la subordonnée s'accorde avec le \emph{Präteritum} dans la principale : \all{Nachdem der Strom ausgefallen war, konnte niemand mehr arbeiten.} (« Après que le courant eut été coupé, plus personne ne pouvait travailler. »)

\textbf{Formules du débat :} \all{Meiner Meinung nach ...} (à mon avis) ; \all{Ich bin dafür / dagegen, weil ...} (je suis pour / contre parce que) ; \all{Ein wichtiger Vorteil ist, dass ...} (un avantage important est que).
\end{aretenir}

\begin{attention}[Pièges à éviter dans ta production]
\begin{itemize}
\item \all{einerseits} et \all{andererseits} ne sont jamais séparés du verbe par le sujet : écris \all{Einerseits ist der Fernunterricht flexibel}, et non \emph{Einerseits der Fernunterricht ist flexibel}.
\item Le plus-que-parfait avec les verbes de mouvement ou de changement d'état se construit avec \all{sein} : \all{der Strom war ausgefallen}, \all{die Schüler waren nach Hause gegangen}. Pour \all{schließen} (fermer) au sens passif « les écoles ont été fermées », utilise le passif : \all{Nachdem die Schulen geschlossen worden waren...}
\item Faux ami : \all{das Netz} signifie « le réseau / le filet » ; pour « Internet », on dit \all{das Internet}.
\item N'oublie pas la majuscule : \all{der fernunterricht} est une faute, il faut \all{der Fernunterricht}.
\item Ponctuation : en allemand, pas d'espace avant le point d'interrogation ou d'exclamation : \all{War das eine gute Erfahrung?}
\end{itemize}
\end{attention}

\courssub{Proposition de production et corrigé commenté}

\begin{exoresolu}[Article modèle : Fernunterricht — eine Chance für Kamerun?]
Rédige un article de 10 à 12 lignes selon le plan imposé.
\tcblower
\textbf{Production modèle :}

\all{Fernunterricht : eine Chance für Kamerun?}

\all{Seit einigen Jahren gibt es auch in Kamerun Fernunterricht, besonders in den Großstädten wie Yaoundé und Douala. Nachdem die Schulen während der Krise geschlossen worden waren, mussten Millionen Schüler zu Hause per Internet, Radio und Fernsehen lernen. War das eine gute Erfahrung?}

\all{Einerseits hat der Fernunterricht klare Vorteile: Die Schüler konnten weiterlernen und haben das Schuljahr nicht verloren. Außerdem ist er flexibel: Man kann die Lektionen wiederholen, wann man will. Andererseits gibt es in Kamerun große Probleme: Die Verbindung ist oft langsam, und der Strom fällt manchmal stundenlang aus. Nachdem der Strom ausgefallen war, konnten viele Schüler nicht mehr am Unterricht teilnehmen. Hinzu kommt: Nicht jede Familie hat ein Handy oder Geld für Internetdaten.}

\all{Meiner Meinung nach ist der Fernunterricht weder ein Fluch noch ein Segen, sondern ein nützliches Werkzeug. Er kann eine echte Chance für Kamerun sein — aber nur, wenn der Staat in Strom, Netze und Computer für alle investiert.}

\textbf{Commentaire des choix de langue :}
\begin{itemize}
\item \textbf{Plus-que-parfait} : \all{Nachdem die Schulen ... geschlossen worden waren} (passif : auxiliaire \all{werden} au Präteritum \all{wurden} + participe passé \all{geworden} + participe passé \all{geschlossen} — ou plus simplement \all{waren geschlossen worden}) et \all{Nachdem der Strom ausgefallen war} (auxiliaire \all{sein}, car \all{ausfallen} exprime un changement d'état). La principale est au Präteritum (\all{mussten}, \all{konnten}, \all{gab}, \all{war}), conformément à la concordance des temps.
\item \textbf{Opposition} : la structure \all{Einerseits ... Andererseits ...} organise le corps de l'article et équilibre les deux points de vue ; le verbe suit directement chaque connecteur (place 2).
\item \textbf{Lexique thématique} : \all{die Verbindung}, \all{der Strom}, \all{das Netz}, \all{teilnehmen an}, \all{die Internetdaten} sont réemployés dans des contextes nouveaux.
\item \textbf{Verbes à particule} : \all{ausfallen} (\all{der Strom fällt ... aus}), \all{weiterlernen}, \all{teilnehmen} — la particule est bien rejetée en fin de phrase à l'indicatif présent.
\item \textbf{Conclusion nuancée} : la formule \all{weder ... noch ...} (« ni... ni... ») permet de dépasser l'alternative Fluch/Segen et de formuler un avis personnel conditionnel, ce qui est attendu au niveau Terminale.
\end{itemize}
\end{exoresolu}

\courssub{Bilan de l'activité}

Cette activité a permis de réinvestir l'ensemble des acquis de la leçon dans une tâche finale de communication authentique. Sur le plan de la \textbf{compréhension}, tu as exploité une lettre de lecteur pour en dégager la problématique et les arguments. Sur le plan du \textbf{lexique}, tu as réemployé activement le vocabulaire des cours à distance et du numérique, à l'oral comme à l'écrit. Sur le plan de la \textbf{grammaire}, tu as manipulé la structure d'opposition \all{einerseits ... andererseits} pour construire une argumentation équilibrée, et le plus-que-parfait dans des propositions temporelles introduites par \all{nachdem} pour établir l'antériorité des faits (avec le choix correct de l'auxiliaire \all{haben}/\all{sein} et la forme passive si nécessaire). Sur le plan de l'\textbf{expression}, tu as pris la parole dans un débat réglé, puis produit un article de presse scolaire structuré (introduction, pour, contre, conclusion personnelle), genre d'écrit fréquent aux évaluations de Terminale. L'auto-vérification finale (majuscules, place du verbe, particules séparables, formation du plus-que-parfait, ponctuation) t'a entraîné à relire et corriger tes propres productions, compétence essentielle pour l'épreuve écrite.

\newpage

\coursec{Méthodes et exercices résolus}

\courssub{Fiches méthodes et exercices résolus}

\begin{definition}[Vocabulaire de base : \all{der Fernunterricht und das Digitale}]
\begin{center}
\begin{tabularx}{\linewidth}{|X|X|X|X|}
\hline
\rowcolor{popblueL}
\textbf{Article} & \textbf{Singulier} & \textbf{Pluriel} & \textbf{Français} \\
\hline
der & Fernunterricht & (kein Plural / die Fernunterrichte) & l'enseignement à distance \\
die & Verbindung & -en & la connexion \\
der & Strom & die Ströme & le courant (électrique) \\
das & Netz & -e & le réseau (Internet) \\
der & Vorteil & -e & l'avantage \\
der & Nachteil & -e & l'inconvénient \\
der & Bildschirm & -e & l'écran \\
die & Datei & -en & le fichier \\
der & Computer & - & l'ordinateur \\
das & Internet & (kein Plural) & Internet \\
die & Schule & -n & l'école \\
der & Lehrer & - & le professeur \\
\hline
\end{tabularx}
\end{center}
\emph{Apprenez chaque nom avec son article, son pluriel et son genre. Les verbes à particule du thème : \all{ausfallen} (le courant coupe : \all{der Strom fällt aus}), \all{anschalten} (allumer), \all{herunterladen} (télécharger), \all{hochladen} (mettre en ligne), \all{teilnehmen an} (participer à + datif).}
\end{definition}

\begin{methode}[Apprendre et utiliser le vocabulaire thématique]
\begin{enumerate}
\item \textbf{Nom + article + pluriel} : ne mémorisez jamais un nom seul. Ex. : \all{die Verbindung, -en} (féminin), \all{der Strom, die Ströme} (masculin, Umlaut au pluriel).
\item \textbf{Familles de mots (Wortfamilien)} : \all{verbinden} (verbe) $\rightarrow$ \all{die Verbindung} (nom) $\rightarrow$ \all{verbunden} (participe/adjectif) ; \all{unterrichten} $\rightarrow$ \all{der Unterricht} $\rightarrow$ \all{der Fernunterricht} (nom composé).
\item \textbf{Verbes à particule séparable} : identifiez la particule (\all{aus-, an-, herunter-, hoch-}). Au présent/prétérit : particule en fin (\all{der Strom fällt aus}). Au participe passé : \all{ge} entre particule et radical (\all{ausgefallen, heruntergeladen, angeschaltet, hochgeladen}).
\item \textbf{Collocations (associations figées)} : \all{eine (gute/schlechte) Verbindung haben}, \all{am Fernunterricht teilnehmen} (datif !), \all{im Internet surfen}, \all{Dateien herunterladen/hochladen}.
\item \textbf{Majuscule systématique} : tout nom allemand prend une majuscule, même au milieu de la phrase (\all{die Verbindung}, \all{der Strom}).
\end{enumerate}
\end{methode}

\begin{exoresolu}[Lexique : phrase à compléter, reformuler et dériver]
Complétez avec le nom correct (article + nominatif) : \all{Während des Fernunterrichts konnten viele Schüler nicht am Unterricht teilnehmen, weil \_\_\_\_\_\_\_\_ (la connexion) zu schlecht war und \_\_\_\_\_\_\_\_ (le courant) oft ausfiel.} Puis donnez le nom dérivé (article + nominatif singulier) de \all{verbinden} et de \all{unterrichten}.
\tcblower
\textbf{Raisonnement} :
\begin{itemize}
\item « la connexion » = sujet de \all{war} $\rightarrow$ nominatif féminin $\rightarrow$ \all{die Verbindung}.
\item « le courant » = sujet de \all{fiel ... aus} $\rightarrow$ nominatif masculin $\rightarrow$ \all{der Strom}.
\item \all{ausfallen} est séparable : prétérit \all{fiel ... aus} (particule \all{aus} en fin).
\end{itemize}
\textbf{Phrase complète} : \all{Während des Fernunterrichts konnten viele Schüler nicht am Unterricht teilnehmen, weil die Verbindung zu schlecht war und der Strom oft ausfiel.}
\textbf{Traduction} : « Pendant les cours à distance, beaucoup d'élèves ne pouvaient pas participer aux cours car la connexion était trop mauvaise et le courant coupait souvent. »
\textbf{Familles de mots} :
\begin{itemize}
\item \all{verbinden} (verbe) $\rightarrow$ \all{die Verbindung} (nom, fém.).
\item \all{unterrichten} (verbe) $\rightarrow$ \all{der Unterricht} (nom, masc.) $\rightarrow$ \all{der Fernunterricht} (nom composé).
\end{itemize}
\textbf{Point de vigilance} : toujours vérifier le genre (article), le cas (ici nominatif sujet) et la place de la particule séparable en fin de proposition principale.
\end{exoresolu}

\begin{propriete}[L'opposition avec \all{einerseits … andererseits}]
\begin{itemize}
\item \textbf{Fonction} : présente les deux faces d'une même réalité (« d'une part … d'autre part »), sans hiérarchie ni concession. Structure \emph{parallèle} recommandée.
\item \textbf{Nature grammaticale} : adverbes de liaison (Konjunktionaladverbien). Ils occupent la \textbf{position 1} (Vorfeld) ou le \textbf{Mittelfeld} (après le verbe).
\item \textbf{Ordre des mots} : verbe en \textbf{deuxième position} après \all{einerseits} / \all{andererseits} en position 1.
\all{Einerseits \textbf{ist} der Fernunterricht praktisch, andererseits \textbf{ist} er anstrengend.}
\item \textbf{Variantes formelles} : \all{zwar … aber} (certes … mais, concession) ; \all{auf der einen Seite … auf der anderen Seite} (plus long, insistant).
\item \textbf{Interdiction} : ne jamais traduire « d'une part » par \all{*auf einer Hand} (calque fautif du français).
\end{itemize}
\end{propriete}

\begin{methode}[{Construire une phrase avec \all{einerseits … andererseits}]}]
\begin{enumerate}
\item Identifiez les deux arguments contraires (avantage / inconvénient).
\item Placez \all{einerseits} au début de la première proposition (ou après le verbe). Conjuguez le verbe en position 2.
\item Placez \all{andererseits} au début de la seconde proposition. Conjuguez le verbe en position 2 (inversion si sujet après verbe).
\item Vérifiez la parallélisme : même temps, même structure de compléments.
\item Relisez : majuscules des noms, cas (accusatif/datif selon le verbe), particule séparable en fin.
\end{enumerate}
\end{methode}

\begin{exoresolu}[Thème : traduction avec \all{einerseits … andererseits}]
Traduisez en allemand : « D'une part, les cours à distance permettent aux élèves de rester en sécurité ; d'autre part, beaucoup n'ont pas d'ordinateur à la maison. »
\tcblower
\textbf{Raisonnement} :
\begin{itemize}
\item « Permettre à qqn de faire » = \all{jmdm. (Dat.) erlauben, zu + Infinitiv} (ou \all{ermöglichen} + Acc. + zu). Choix : \all{erlauben} + datif.
\item « Rester en sécurité » = \all{sicher bleiben} (verbe \all{bleiben} $\rightarrow$ auxiliaire \all{sein} aux temps composés).
\item « Ne pas avoir d'ordinateur » = \all{keinen Computer haben} (accusatif masculin négatif \all{keinen}).
\item Structure parallèle : deux principales liées par point-virgule ; verbe en position 2 des deux côtés.
\end{itemize}
\textbf{Traduction} : \all{Einerseits erlaubt der Fernunterricht den Schülern, sicher zu bleiben; andererseits haben viele keinen Computer zu Hause.}
\textbf{Vérifications} :
\begin{itemize}
\item \all{der Fernunterricht} (sujet sg.) $\rightarrow$ \all{erlaubt}.
\item \all{den Schülern} (datif pl. après \all{erlauben}).
\item \all{keinen Computer} (accusatif sg. masc. négatif).
\item Après \all{andererseits} : inversion \all{haben viele} (verbe position 2).
\end{itemize}
\end{exoresolu}

\begin{propriete}[Le plus-que-parfait (Plusquamperfekt) et les propositions temporelles]
\begin{itemize}
\item \textbf{Formation} : Auxiliaire \all{haben} ou \all{sein} au \textbf{Präteritum} + \textbf{Participe II}.
\begin{center}
\begin{tabular}{|l|l|l|}
\hline
\rowcolor{popblueL}
Verbe & Auxiliaire Präteritum & Participe II / Forme 1ère pers. sg. \\
\hline
\all{lernen} (faible) & \all{hatte} & \all{hatte gelernt} \\
\all{hochladen} (séparable) & \all{hatte} & \all{hatte hochgeladen} \\
\all{anschalten} (séparable) & \all{hatte} & \all{hatte angeschaltet} \\
\all{kommen / bleiben / ausfallen} & \all{war} & \all{war gekommen / geblieben / ausgefallen} \\
\hline
\end{tabular}
\end{center}
\item \textbf{Choix de l'auxiliaire} : Même règle qu'au Perfekt. \all{sein} pour : mouvement (\all{kommen, gehen}), changement d'état (\all{ausfallen, einschlafen, sterben}), \all{sein, bleiben, werden}. \all{haben} pour le reste (transitifs, réflexifs, la plupart intransitifs).
\item \textbf{Emploi} : Antériorité dans le passé (raconter un passé antérieur à un autre passé).
\item \textbf{Propositions temporelles clés} :
\begin{center}
\begin{tabularx}{\linewidth}{|X|X|X|X|}
\hline
\rowcolor{popblueL}
Conjonction & Temps subordonnée & Temps principale (récit) & Sens \\
\hline
\all{nachdem} & \textbf{Plusquamperfekt} & Präteritum & Après que (antériorité stricte) \\
\all{bevor} / \all{ehe} & Plusquamperfekt (ou Présent) & Präteritum & Avant que \\
\all{als} & Präteritum / Plusquamperfekt & Präteritum & Quand (ponctuel passé) \\
\all{während} & Präteritum / Plusquamperfekt & Präteritum & Pendant que \\
\all{seitdem} / \all{seit} & Présent / Präteritum / Perfekt & Présent / Präteritum & Depuis que \\
\hline
\end{tabularx}
\end{center}
\item \textbf{Place du verbe} : En fin de subordonnée (\all{nachdem die Schulen \textbf{geschlossen hatten}}). Si subordonnée en tête : inversion en principale (\all{\textbf{begann} der Fernunterricht}).
\end{itemize}
\end{propriete}

\begin{methode}[Réussir la version et la transformation au plus-que-parfait]
\begin{enumerate}
\item \textbf{Repérez l'antériorité} : deux actions passées, l'une avant l'autre. La plus ancienne $\rightarrow$ Plusquamperfekt.
\item \textbf{Choisissez l'auxiliaire} : \all{haben} ou \all{sein} (règle du mouvement/changement d'état). Attention à \all{ausfallen, bleiben, kommen} $\rightarrow$ \all{sein}.
\item \textbf{Construisez le Participe II} : Verbe faible \all{ge...t} ; fort \all{ge...en} ; séparable \all{Particule+ge+Radikal+t/en} (\all{ausgefallen, angeschaltet, hochgeladen, heruntergeladen}).
\item \textbf{Version (Allemand $\rightarrow$ Français)} : \all{nachdem} + Plusquamperfekt $\rightarrow$ « Après que + plus-que-parfait » ou « Après avoir + participe passé » (plus fluide).
\item \textbf{Transformation (Coordination $\rightarrow$ Subordination)} : Reliez les deux verbes par \all{nachdem} + Plusquamperfekt (1ère action) ; principale au Präteritum (2nde action). Verbe conjugué en fin de subordonnée ; inversion en principale si subordonnée initiale.
\end{enumerate}
\end{methode}

\begin{exoresolu}[Version et transformation : le plus-que-parfait]
1) Traduisez en français : \all{Nachdem der Lehrer den Kurs hochgeladen hatte, konnten die Schüler die Dateien herunterladen.}
2) Transformez en subordonnée temporelle avec \all{nachdem} + Plusquamperfekt : \all{Die Schüler schalteten den Computer an und begannen die Arbeit.}
\tcblower
\textbf{1) Version} :
\textbf{Raisonnement} : \all{hatte ... hochgeladen} = Plusquamperfekt (\all{haben} + Part. II \all{hochgeladen}, particule \all{hoch} soudée). \all{nachdem} marque l'antériorité. \all{konnten} = Präteritum de \all{können}.
\textbf{Traduction} : « Après que le professeur eut mis le cours en ligne (ou : Après avoir mis le cours en ligne), les élèves purent (ou : ont pu) télécharger les fichiers. »

\textbf{2) Transformation} :
\textbf{Raisonnement} : Action 1 (antérieure) : \all{anschalten} (transitif, séparable) $\rightarrow$ \all{hatten ... angeschaltet}. Action 2 : \all{beginnen} (fort, intransitif mais \all{haben} car pas de déplacement) $\rightarrow$ \all{begannen}. Liaison par \all{nachdem}.
\textbf{Réponse} : \all{Nachdem die Schüler den Computer angeschaltet hatten, begannen sie die Arbeit.}
\textit{Alternative (double Plusquamperfekt en coordination, moins élégant) :} \all{Die Schüler hatten den Computer angeschaltet und hatten die Arbeit begonnen.}
\textbf{Vérifications} :
\begin{itemize}
\item \all{hatten} (aux. \all{haben} Präteritum) en fin de subordonnée.
\item \all{angeschaltet} (Part. II correct : \all{an} + \all{ge} + \all{schalt} + \all{et}).
\item \all{begannen sie} (inversion en principale, Präteritum de \all{beginnen} : \all{a} $\rightarrow$ \all{a}).
\end{itemize}
\end{exoresolu}

\begin{methode}[Comprendre un texte sur le numérique (global puis détaillé)]
Pour réussir les questions de compréhension au Bac, procédez en cinq étapes rigoureuses :
\begin{enumerate}
\item \textbf{Lecture globale} : lisez une première fois sans dictionnaire. Repérez le thème (\all{der Fernunterricht}), le type de texte (article, rapport, témoignage) et la thèse générale (avantages ? problèmes ? les deux ?).
\item \textbf{Repérage des mots-clés} : soulignez le champ lexical : \all{die Verbindung, das Netz/Internet, der Strom, der Vorteil, der Nachteil, der Bildschirm, die Datei, herunterladen, hochladen, ausfallen}. Ce sont les porteurs de sens.
\item \textbf{Exploitation des connecteurs logiques} : identifiez \all{einerseits … andererseits} (opposition), \all{aber / jedoch} (concession), \all{deshalb / darum} (conséquence), \all{weil / denn} (cause), \all{nachdem / als / während} (temporalité). Ils structurent l'argumentation.
\item \textbf{Lecture détaillée ciblée} : relisez phrase par phrase en reliant \textbf{chaque question} à un passage précis. Ne répondez jamais « de mémoire » ou par culture générale.
\item \textbf{Rédaction de la réponse en allemand} : phrase complète, verbe à la bonne place (pos. 2 principale, fin subordonnée), reformulation (pas de copier-coller), majuscules aux noms.
\end{enumerate}
\textbf{Réflexes types} : \all{Warum?} $\rightarrow$ réponse avec \all{weil} (verbe fin) ou \all{denn} (verbe pos. 2) ; \all{Wie?} $\rightarrow$ description manière/moyen ; \all{Was für Folgen?} $\rightarrow$ \all{deshalb / darum} + conséquence.
\end{methode}

\begin{exoresolu}[Questions de compréhension détaillée]
\textbf{Texte support} : \all{Viele Schüler in Kamerun hatten während des Fernunterrichts große Probleme. Einerseits konnten sie zu Hause sicher lernen, andererseits war die Internetverbindung oft schlecht, und der Strom fiel manchmal aus. Nachdem die Schulen geschlossen hatten, mussten die Lehrer ihre Kurse online geben.}
\textbf{Questions} :
\begin{enumerate}
\item \all{Welche zwei Seiten des Fernunterrichts nennt der Text?}
\item \all{Warum war das Lernen zu Hause schwierig?}
\item \all{Was mussten die Lehrer tun, nachdem die Schulen geschlossen hatten?}
\end{enumerate}
\tcblower
\textbf{Raisonnement} :
\begin{itemize}
\item Q1 vise la structure \all{einerseits … andererseits} (opposition avantages/inconvénients).
\item Q2 (\all{Warum}) appelle une cause $\rightarrow$ \all{weil} + verbe en fin.
\item Q3 porte sur la proposition temporelle \all{nachdem} + Plusquamperfekt (\all{hatten ... geschlossen}) $\rightarrow$ réponse au Präteritum (\all{mussten}).
\end{itemize}
\textbf{Réponses rédigées} :
\begin{enumerate}
\item \all{Der Text nennt einerseits den Vorteil, dass die Schüler zu Hause sicher lernen konnten, und andererseits den Nachteil, dass die Internetverbindung oft schlecht war und der Strom manchmal ausfiel.}
\item \all{Das Lernen zu Hause war schwierig, weil die Internetverbindung oft schlecht war und der Strom manchmal ausfiel.}
\item \all{Nachdem die Schulen geschlossen hatten, mussten die Lehrer ihre Kurse online geben.}
\end{enumerate}
\textbf{Conclusion} : chaque réponse est une phrase autonome, grammaticalement correcte (place du verbe, cas, majuscules), et reformule le texte sans le recopier.
\end{exoresolu}

\begin{methode}[Produire un court paragraphe d'opinion (expression écrite guidée)]
\begin{enumerate}
\item \textbf{Accroche / Sujet} : une phrase d'introduction. \all{Der Fernunterricht spielt heute eine große Rolle.}
\item \textbf{Développement structuré} : utilisez \all{einerseits … andererseits} pour équilibrer. Citez un avantage concret (sécurité, gain de temps) et un problème concret (connexion, électricité, matériel).
\item \textbf{Exemple personnel / illustration} : au passé. Si antériorité $\rightarrow$ \all{nachdem} + Plusquamperfekt. Ex. : \all{Nachdem ich die Datei heruntergeladen hatte, fiel der Strom aus.}
\item \textbf{Opinion personnelle} : \all{Meiner Meinung nach …} (inversion) ou \all{Ich finde, dass …} (verbe en fin après \all{dass}).
\item \textbf{Relecture finale (check-list)} : Majuscules (noms) ? Verbe position 2 (principale) / fin (subordonnée) ? Particules séparables en fin ? Articles/Adjectifs accordés ? Orthographe \all{ß / ss} / Umlaute ?
\end{enumerate}
\textbf{Plan type} : Thèse $\rightarrow$ \all{einerseits} $\rightarrow$ \all{andererseits} $\rightarrow$ Exemple vécu (\all{nachdem} + Plusquamperfekt) $\rightarrow$ \all{Meiner Meinung nach}.
\end{methode}

\begin{exoresolu}[Rédaction guidée (environ 60--80 mots)]
\textbf{Consigne} : Rédigez un court texte en allemand : « Les cours à distance au Cameroun : avantages et problèmes ». Utilisez \all{einerseits … andererseits}, au moins un plus-que-parfait avec \all{nachdem}, et donnez votre avis personnel.
\tcblower
\textbf{Raisonnement} : Application du plan en 5 étapes. Contrôle de chaque verbe. Choix de l'auxiliaire \all{sein} pour \all{ausfallen} (changement d'état) $\rightarrow$ \all{ist ausgefallen} (Perfekt) $\rightarrow$ \all{war ausgefallen} (Plusquamperfekt).
\textbf{Production modèle} :
\all{Der Fernunterricht in Kamerun hat Vor- und Nachteile. Einerseits können die Schüler zu Hause sicher lernen und sparen lange Wege. Andererseits ist die Internetverbindung oft instabil, und viele Familien besitzen keinen eigenen Computer. Nachdem ich einmal eine Stunde vergeblich auf eine stabile Verbindung gewartet hatte, war der Strom komplett ausgefallen. Meiner Meinung nach ist der Fernunterricht eine Chance, aber der Staat muss das Stromnetz und das Internet dringend verbessern.}
\textbf{Traduction} : « L'enseignement à distance au Cameroun a des avantages et des inconvénients. D'une part, les élèves peuvent apprendre en sécurité à la maison et économisent de longs trajets. D'autre part, la connexion Internet est souvent instable et beaucoup de familles ne possèdent pas d'ordinateur personnel. Après avoir attendu en vain une connexion stable pendant une heure, le courant avait complètement coupé. À mon avis, l'enseignement à distance est une chance, mais l'État doit améliorer d'urgence le réseau électrique et Internet. »
\textbf{Vérifications ciblées} :
\begin{itemize}
\item \all{war ausgefallen} (Plusquamperfekt \all{sein} + Part. II \all{ausgefallen} -- correct).
\item \all{Meiner Meinung nach ist ...} (inversion après complément de phrase).
\item \all{dass}-phrase absente ici (opinion directe), mais \all{dass} aurait mis le verbe en fin.
\item Tous les noms en majuscules : \all{Vor- und Nachteile, Schüler, Wege, Internetverbindung, Familien, Computer, Stunde, Verbindung, Strom, Meinung, Chance, Staat, Stromnetz, Internet}.
\end{itemize}
\end{exoresolu}

\begin{attention}[Les 5 erreurs qui coûtent des points au Bac]
\begin{enumerate}
\item \textbf{Calque de « d'une part »} : \all{*auf einer Hand} n'existe pas. La seule formule correcte est \all{einerseits … andererseits} (ou \all{auf der einen Seite … auf der anderen Seite}).
\item \textbf{Place du verbe après \all{nachdem}} : Subordonnée = verbe conjugué \textbf{en fin}. \all{Nachdem die Schulen geschlossen hatten, …} (pas \all{*hatten geschlossen}).
\item \textbf{Mauvais auxiliaire au Plusquamperfekt} : \all{ausfallen, kommen, gehen, bleiben, werden, sein} prennent \all{sein}. \all{Der Strom war ausgefallen} (pas \all{*hatte ausgefallen}). \all{Er war gekommen} (pas \all{*hatte gekommen}).
\item \textbf{Oubli de la majuscule et de l'article} : \all{das Netz, die Verbindung, der Strom, der Fernunterricht}. Un nom sans majuscule ou avec un article incorrect (ex. \all{*die Strom}) est pénalisé. Apprenez \textbf{Article + Singulier + Pluriel} par cœur.
\item \textbf{Faux amis et particules séparables} :
\begin{itemize}
\item \all{bekommen} = recevoir (pas « devenir» $\rightarrow$ \all{werden}).
\item \all{ausfallen, anschalten, herunterladen, hochladen, teilnehmen} : particule \textbf{détachée en fin} de proposition principale (\all{der Strom fällt aus}, \all{ich lade die Datei herunter}) ; \textbf{soudée avec \all{ge}} au Participe II (\all{ausgefallen, heruntergeladen}).
\end{itemize}
\end{enumerate}
\end{attention}

\newpage

\coursec{Exercices}

\courssub{Je vérifie mes connaissances}

\begin{exercice}[QCM : le vocabulaire du numérique]
Entoure la bonne réponse.
\begin{enumerate}
\item \all{der Fernunterricht} signifie : a) le cours particulier \quad b) les cours à distance \quad c) la télévision scolaire
\item \all{die Verbindung} est : a) le courant électrique \quad b) la connexion \quad c) le mot de passe
\item Le pluriel de \all{der Vorteil} est : a) \all{die Vorteilen} \quad b) \all{die Vorteils} \quad c) \all{die Vorteile}
\item \all{das Netz} signifie : a) le réseau \quad b) l'écran \quad c) le clavier
\item Quand il n'y a plus d'électricité, on dit : a) \all{der Strom ist ausgefallen} \quad b) \all{der Strom hat gelernt} \quad c) \all{der Strom ist gelaufen}
\end{enumerate}
\end{exercice}

\begin{exercice}[Vrai ou faux ? Justifie ta réponse]
Réponds par \all{richtig} ou \all{falsch} et justifie en une phrase complète en allemand, d'après le texte étudié \all{Fernunterricht – Chancen und Probleme}.
\begin{enumerate}
\item \all{Beim Fernunterricht müssen die Schüler nicht zur Schule fahren.}
\item \all{Der Fernunterricht hat nur Vorteile und keine Nachteile.}
\item \all{Eine schlechte Internetverbindung ist ein großes Problem für den Fernunterricht.}
\item \all{Nach dem Stromausfall konnten die Schüler den Kurs ohne Probleme fortsetzen.}
\end{enumerate}
\end{exercice}

\begin{exercice}[Vocabulaire : articles et familles de mots]
Complète le tableau avec l'article et le pluriel de chaque nom, puis forme une phrase complète avec deux des mots.
\begin{center}
\begin{tabularx}{\linewidth}{|l|X|X|}
\hline
\rowcolor{popblueL} Nom & Article & Pluriel \\
\hline
Fernunterricht & & \\
\hline
Verbindung & & \\
\hline
Strom & & \\
\hline
Nachteil & & \\
\hline
Netz & & \\
\hline
\end{tabularx}
\end{center}
\end{exercice}

\begin{exercice}[Conjugaison : le plus-que-parfait]
Mets les verbes entre parenthèses au \all{Plusquamperfekt}.
\begin{enumerate}
\item \all{Nachdem der Lehrer die E-Mail (schicken), begann der Kurs.}
\item \all{Die Schüler (aufstehen) schon um sechs Uhr.}
\item \all{Ich (vergessen) mein Passwort, deshalb konnte ich mich nicht anmelden.}
\item \all{Nachdem wir die Hausaufgaben (machen), gingen wir ins Bett.}
\item \all{Der Strom (ausfallen), bevor der Unterricht zu Ende war.}
\end{enumerate}
\end{exercice}

\courssub{Je m'entraîne}

\begin{exercice}[Thème : traduis en allemand]
Traduis les phrases suivantes en allemand. Attention à la place du verbe et au temps.
\begin{enumerate}
\item D'une part les cours en ligne sont pratiques, d'autre part la connexion est souvent mauvaise.
\item Après que le courant eut été coupé, nous ne pûmes plus suivre le cours.
\item Quand le professeur avait envoyé les documents, les élèves pouvaient commencer.
\item D'une part ma sœur aime travailler sur l'ordinateur, d'autre part elle préfère le contact avec ses camarades.
\item Après que j'eus trouvé mon mot de passe, je me connectai à la plateforme.
\end{enumerate}
\end{exercice}

\begin{exercice}[Transformation : construis l'opposition]
À partir des deux arguments donnés, construis une seule phrase avec \all{einerseits … andererseits}. Attention à l'inversion du sujet après \all{andererseits}.

Exemple : \all{Der Fernunterricht spart Zeit. / Er ist oft einsam.} → \all{Einerseits spart der Fernunterricht Zeit, andererseits ist er oft einsam.}
\begin{enumerate}
\item \all{Die Schüler müssen nicht zur Schule fahren. / Sie vermissen ihre Freunde.}
\item \all{Man kann zu Hause lernen. / Das Internet funktioniert nicht immer.}
\item \all{Der Computer ist sehr praktisch. / Er kostet viel Geld.}
\item \all{Die Lehrer schicken die Aufgaben per E-Mail. / Viele Schüler haben keinen Strom.}
\end{enumerate}
\end{exercice}

\begin{exercice}[Texte à trous : la journée d'un élève en ligne]
Complète le texte avec les mots de la liste suivante : \all{Verbindung – Stromausfall – Passwort – Netz – Nachteil – Fernunterricht}.

\begin{textebox}[Mein Tag mit dem Computer]
Heute hatte ich \all{\_\_\_\_\_\_\_\_\_\_}. Um acht Uhr wollte ich mich anmelden, aber ich hatte mein \all{\_\_\_\_\_\_\_\_\_\_} vergessen. Nach zehn Minuten fand ich es endlich. Leider war die \all{\_\_\_\_\_\_\_\_\_\_} sehr schlecht, und das \all{\_\_\_\_\_\_\_\_\_\_} war oft unterbrochen. Um elf Uhr gab es sogar einen \all{\_\_\_\_\_\_\_\_\_\_} : kein Licht, kein Internet ! Das ist ein großer \all{\_\_\_\_\_\_\_\_\_\_} des Lernens zu Hause.
\end{textebox}
\end{exercice}

\begin{exercice}[Appariements : mots et traductions]
Relie chaque mot allemand à sa traduction française.
\begin{center}
\begin{tabularx}{\linewidth}{|X|X|}
\hline
\rowcolor{popblueL} Deutsch & Français \\
\hline
1. \all{der Stromausfall} & a. manquer, faire défaut \\
\hline
2. \all{die Störung} & b. remplacer \\
\hline
3. \all{das Passwort} & c. la panne de courant \\
\hline
4. \all{ersetzen} & d. le mot de passe \\
\hline
5. \all{fehlen} & e. la perturbation, la panne \\
\hline
\end{tabularx}
\end{center}
\end{exercice}

\courssub{Je me prépare au Bac}

\begin{exercice}[Compréhension de l'écrit : texte et questions]
Lis le texte suivant, puis réponds aux questions \textbf{en allemand, par des phrases complètes}.

\begin{textebox}[Digitaler Unterricht an deutschen Schulen]
An vielen deutschen Schulen gehört der Computer heute zum Alltag. Die Schüler lernen mit Tablets, schicken ihre Hausaufgaben per E-Mail und benutzen digitale Plattformen. Einerseits bietet diese Entwicklung große Chancen : die Schüler lernen schneller, Informationen zu finden, und sie können auch zu Hause arbeiten, wenn sie krank sind. Andererseits gibt es ernste Probleme. Nicht jede Familie hat einen guten Computer oder eine schnelle Internetverbindung. Manche Schüler sitzen viele Stunden vor dem Bildschirm und bewegen sich zu wenig. Nachdem eine Schule in Bayern neue Tablets gekauft hatte, stellte man fest, dass viele Geräte nach einem Jahr schon kaputt waren. Auch die Lehrer brauchen eine gute Ausbildung : nachdem sie gelernt hatten, die neuen Programme zu benutzen, konnten sie ihren Unterricht modern gestalten. Experten sagen : die Technik darf den Lehrer nicht ersetzen, sie muss ihm helfen. Der Kontakt zwischen Lehrern und Schülern bleibt das Wichtigste in der Schule.
\end{textebox}

Questions :
\begin{enumerate}
\item \all{Womit lernen die Schüler an deutschen Schulen heute?} (2 éléments)
\item \all{Nennen Sie zwei Chancen des digitalen Unterrichts.}
\item \all{Welche Probleme werden im Text genannt?} (2 éléments)
\item \all{Was stellte man an der Schule in Bayern fest?}
\item \all{Was sagen die Experten über die Rolle der Technik?}
\item \textbf{Question de langue} : souligne dans le texte une proposition avec \all{nachdem} et explique quel temps on utilise dans la proposition subordonnée et dans la proposition principale.
\item \textbf{Question de langue} : relève dans le texte la structure d'opposition et recopie-la.
\end{enumerate}
\end{exercice}

\begin{exercice}[Thème et version]
\textbf{A. Thème.} Traduis en allemand les phrases suivantes.
\begin{enumerate}
\item D'une part les cours à distance permettent d'économiser du temps, d'autre part beaucoup d'élèves se sentent seuls.
\item Après que le professeur eut expliqué la leçon, les élèves posèrent des questions.
\item Quand la connexion avait été rétablie, nous pûmes continuer le cours.
\item D'une part mon frère travaille volontiers sur la tablette, d'autre part il a souvent mal aux yeux.
\item Après que le courant eut été coupé, je ne pus plus envoyer mes devoirs.
\end{enumerate}

\textbf{B. Version.} Traduis en français le paragraphe suivant.

\begin{textebox}[Version]
Einerseits war der Fernunterricht für mich eine gute Erfahrung, andererseits hatte ich viele technische Probleme. Nachdem der Strom ausgefallen war, konnte ich zwei Stunden lang nicht arbeiten. Nachdem ich endlich eine gute Verbindung gefunden hatte, schickte ich meine Hausaufgaben an den Lehrer.
\end{textebox}
\end{exercice}

\begin{exercice}[Rédaction guidée type Bac]
\textbf{Sujet} : \all{Fernunterricht in meinem Land : Vorteile und Nachteile}

Tu es élève en Terminale au Cameroun. Pendant une période difficile, ton école a organisé des cours à distance. Rédige un texte de \textbf{80 à 100 mots} en allemand dans lequel tu :
\begin{itemize}
\item présentes deux avantages et deux inconvénients des cours à distance dans ton pays ;
\item utilises \textbf{obligatoirement} les deux structures suivantes :
\begin{enumerate}
\item \all{einerseits … andererseits} (avec l'inversion du sujet) ;
\item une proposition temporelle avec \all{nachdem} + \all{Plusquamperfekt} ;
\end{enumerate}
\item donnes ton avis personnel dans la conclusion.
\end{itemize}
\textbf{Conseil} : fais d'abord un plan en trois parties (introduction, avantages/inconvénients, conclusion) et vérifie la place du verbe dans chaque phrase.
\end{exercice}

\newpage

\coursec{Corrigés des exercices}

\begin{corrige}[Exercice 1 — QCM : le vocabulaire du numérique]
\begin{enumerate}
\item \textbf{b)} \all{der Fernunterricht} = les cours à distance (mot composé : \all{fern} « à distance » + \all{der Unterricht} « le cours, l'enseignement »).
\item \textbf{b)} \all{die Verbindung} = la connexion (du verbe \all{verbinden} « relier »). Le courant électrique se dit \all{der Strom}.
\item \textbf{c)} \all{der Vorteil, die Vorteile} : pluriel en \all{-e}, sans \all{-n} ni \all{-s}. Même formation pour le contraire : \all{der Nachteil, die Nachteile}.
\item \textbf{a)} \all{das Netz} = le réseau. L'écran se dit \all{der Bildschirm}, le clavier \all{die Tastatur}.
\item \textbf{a)} \all{Der Strom ist ausgefallen.} = « Le courant est tombé en panne. » \all{ausfallen} est un verbe de changement d'état : il se conjugue avec l'auxiliaire \all{sein} au \all{Perfekt} (\all{ist ausgefallen}).
\end{enumerate}
\end{corrige}

\begin{corrige}[Exercice 2 — Vrai ou faux ? Justifie ta réponse]
\begin{enumerate}
\item \all{Richtig. Beim Fernunterricht lernen die Schüler zu Hause, deshalb müssen sie nicht zur Schule fahren.}\\
« Vrai. Avec les cours à distance, les élèves apprennent à la maison, donc ils n'ont pas besoin d'aller à l'école. »
\item \all{Falsch. Der Fernunterricht hat nicht nur Vorteile, sondern auch Nachteile, zum Beispiel technische Probleme und Einsamkeit.}\\
« Faux. Les cours à distance n'ont pas seulement des avantages, mais aussi des inconvénients, par exemple des problèmes techniques et la solitude. »
\item \all{Richtig. Wenn die Internetverbindung schlecht ist, können die Schüler dem Kurs nicht gut folgen.}\\
« Vrai. Quand la connexion Internet est mauvaise, les élèves ne peuvent pas bien suivre le cours. »
\item \all{Falsch. Nach dem Stromausfall konnten die Schüler den Kurs nicht fortsetzen, weil es kein Internet gab.}\\
« Faux. Après la panne de courant, les élèves n'ont pas pu continuer le cours, car il n'y avait pas d'Internet. »
\end{enumerate}
\textbf{Points de vigilance} : attention à la place du verbe conjugué. Après \all{deshalb} (mot en tête de phrase) : inversion, \all{deshalb müssen sie...}. Dans la subordonnée en \all{weil} : verbe conjugué en fin de phrase, \all{weil es kein Internet gab}. Remarque aussi la construction \all{dem Kurs folgen} « suivre le cours » : \all{folgen} se construit avec le datif.
\end{corrige}

\begin{corrige}[Exercice 3 — Vocabulaire : articles et familles de mots]
\begin{center}
\begin{tabularx}{\linewidth}{|l|X|X|}
\hline
\rowcolor{popblueL} Nom & Article & Pluriel \\
\hline
Fernunterricht & der & (pluriel rare : die Fernunterrichte) \\
\hline
Verbindung & die & die Verbindungen \\
\hline
Strom & der & die Ströme (souvent sans pluriel) \\
\hline
Nachteil & der & die Nachteile \\
\hline
Netz & das & die Netze \\
\hline
\end{tabularx}
\end{center}
\textbf{Phrases possibles} :
\begin{itemize}
\item \all{Eine schlechte Verbindung ist ein großer Nachteil des Fernunterrichts.} « Une mauvaise connexion est un grand inconvénient des cours à distance. »
\item \all{Ohne Strom funktioniert das Netz nicht.} « Sans courant, le réseau ne fonctionne pas. »
\end{itemize}
\textbf{Rappels} :
\begin{itemize}
\item Les noms en \all{-ung} sont toujours féminins (\all{die Verbindung, die Störung, die Lösung}) et font leur pluriel en \all{-en}.
\item \all{des Fernunterrichts} : génitif masculin en \all{-(e)s}, sans article devant ce nom abstrait.
\item Familles de mots : \all{verbinden} (verbe) → \all{die Verbindung} (nom) ; \all{stören} (verbe) → \all{die Störung} (nom) ; \all{der Vorteil} ↔ \all{der Nachteil} (antonymes).
\end{itemize}
\end{corrige}

\begin{corrige}[Exercice 4 — Conjugaison : le plus-que-parfait]
\begin{enumerate}
\item \all{Nachdem der Lehrer die E-Mail \textbf{geschickt hatte}, begann der Kurs.}\\
\all{schicken} : verbe faible, participe \all{geschickt} + auxiliaire \all{haben} au \all{Präteritum} : \all{hatte}.
\item \all{Die Schüler \textbf{waren} schon um sechs Uhr \textbf{aufgestanden}.}\\
\all{aufstehen} : verbe de mouvement → auxiliaire \all{sein} ; participe \all{aufgestanden} (verbe fort à particule séparable : \all{auf-} + \all{ge-} + \all{standen}).
\item \all{Ich \textbf{hatte} mein Passwort \textbf{vergessen}, deshalb konnte ich mich nicht anmelden.}\\
\all{vergessen} : verbe fort, participe \all{vergessen} ; le préfixe inséparable \all{ver-} supprime le \all{ge-} du participe.
\item \all{Nachdem wir die Hausaufgaben \textbf{gemacht hatten}, gingen wir ins Bett.}\\
\all{machen} : verbe faible → \all{gemacht} + \all{hatten}.
\item \all{Der Strom \textbf{war ausgefallen}, bevor der Unterricht zu Ende war.}\\
\all{ausfallen} : changement d'état → auxiliaire \all{sein} : \all{war ausgefallen}.
\end{enumerate}
\textbf{Règle} : \all{Plusquamperfekt} = \all{hatte/n} ou \all{war/en} + participe passé. On choisit \all{sein} pour les verbes de mouvement ou de changement d'état (\all{gehen, aufstehen, ausfallen}), \all{haben} pour tous les autres (y compris les verbes pronominaux : \all{ich hatte mich angemeldet}).
\begin{center}
\begin{tabular}{|l|l|l|}
\hline
\rowcolor{popblueL} Personne & \all{haben} & \all{sein} \\
\hline
ich & hatte & war \\
\hline
du & hattest & warst \\
\hline
er/sie/es & hatte & war \\
\hline
wir & hatten & waren \\
\hline
ihr & hattet & wart \\
\hline
sie/Sie & hatten & waren \\
\hline
\end{tabular}
\end{center}
\end{corrige}

\begin{corrige}[Exercice 5 — Thème : traduis en allemand]
\begin{enumerate}
\item \all{Einerseits sind die Online-Kurse praktisch, andererseits ist die Verbindung oft schlecht.}\\
Après \all{einerseits} et \all{andererseits} en tête de phrase : inversion (verbe conjugué en 2\textsuperscript{e} position, sujet après le verbe).
\item \all{Nachdem der Strom abgeschaltet worden war, konnten wir den Kurs nicht mehr verfolgen.}\\
Passif au \all{Plusquamperfekt} : \all{worden war} ; « ne... plus » = \all{nicht mehr} ; « suivre un cours » = \all{einen Kurs verfolgen}.
\item \all{Als der Lehrer die Dokumente geschickt hatte, konnten die Schüler anfangen.}\\
Passé unique au français → \all{als} (et non \all{wenn}) ; subordonnée au \all{Plusquamperfekt} (antériorité), principale au \all{Präteritum} avec inversion après la subordonnée en tête.
\item \all{Einerseits arbeitet meine Schwester gern am Computer, andererseits zieht sie den Kontakt mit ihren Kameraden vor.}\\
« Préférer » = \all{vorziehen} (verbe fort à particule séparable : \all{zieht ... vor} ; participe \all{vorgezogen}).
\item \all{Nachdem ich mein Passwort gefunden hatte, meldete ich mich bei der Plattform an.}\\
\all{sich anmelden} : verbe pronominal à particule séparable : \all{meldete ... an} ; « se connecter à » = \all{sich bei + Dativ anmelden}.
\end{enumerate}
\end{corrige}

\begin{corrige}[Exercice 6 — Transformation : construis l'opposition]
\begin{enumerate}
\item \all{Einerseits müssen die Schüler nicht zur Schule fahren, andererseits vermissen sie ihre Freunde.}
\item \all{Einerseits kann man zu Hause lernen, andererseits funktioniert das Internet nicht immer.}
\item \all{Einerseits ist der Computer sehr praktisch, andererseits kostet er viel Geld.}
\item \all{Einerseits schicken die Lehrer die Aufgaben per E-Mail, andererseits haben viele Schüler keinen Strom.}
\end{enumerate}
\textbf{Points de vigilance} :
\begin{itemize}
\item \all{Einerseits} et \all{andererseits} occupent la 1\textsuperscript{re} position : le verbe conjugué vient juste après, puis le sujet (\all{andererseits vermissen \textbf{sie}}, et non \all{andererseits sie vermissen}).
\item On peut aussi placer ces mots à l'intérieur de la phrase, sans inversion : \all{Die Schüler müssen einerseits nicht zur Schule fahren, vermissen aber andererseits ihre Freunde.}
\item Ne pas calquer le français « d'une part... d'autre part », qui ne provoque pas d'inversion.
\end{itemize}
\end{corrige}

\begin{corrige}[Exercice 7 — Texte à trous : la journée d'un élève en ligne]
\textbf{Texte complété} :\\
\all{Heute hatte ich \textbf{Fernunterricht}. Um acht Uhr wollte ich mich anmelden, aber ich hatte mein \textbf{Passwort} vergessen. Nach zehn Minuten fand ich es endlich. Leider war die \textbf{Verbindung} sehr schlecht, und das \textbf{Netz} war oft unterbrochen. Um elf Uhr gab es sogar einen \textbf{Stromausfall} : kein Licht, kein Internet ! Das ist ein großer \textbf{Nachteil} des Lernens zu Hause.}

\textbf{Traduction} : « Aujourd'hui j'avais cours à distance. À huit heures, je voulais me connecter, mais j'avais oublié mon mot de passe. Au bout de dix minutes, je l'ai enfin trouvé. Malheureusement, la connexion était très mauvaise et le réseau était souvent interrompu. À onze heures, il y a même eu une panne de courant : pas de lumière, pas d'Internet ! C'est un grand inconvénient de l'apprentissage à la maison. »

\textbf{Indices de résolution} : \all{mein} exige un nom neutre (\all{das Passwort}) ; l'article \all{die} exige un nom féminin (\all{die Verbindung}) ; \all{einen} (accusatif masculin) exige \all{der Stromausfall} ; \all{das} exige un neutre (\all{das Netz}) ; \all{ein großer} (nominatif masculin) exige \all{der Nachteil}.
\end{corrige}

\begin{corrige}[Exercice 8 — Appariements : mots et traductions]
\begin{center}
\begin{tabularx}{\linewidth}{|X|X|}
\hline
\rowcolor{popblueL} Deutsch & Français \\
\hline
1. \all{der Stromausfall} & \textbf{c.} la panne de courant \\
\hline
2. \all{die Störung} & \textbf{e.} la perturbation, la panne \\
\hline
3. \all{das Passwort} & \textbf{d.} le mot de passe \\
\hline
4. \all{ersetzen} & \textbf{b.} remplacer \\
\hline
5. \all{fehlen} & \textbf{a.} manquer, faire défaut \\
\hline
\end{tabularx}
\end{center}
\textbf{Attention au faux ami} : \all{fehlen} signifie « manquer » au sens de « faire défaut » (\all{Mir fehlt das Geld.} « L'argent me manque »), et non « rater ». « Manquer quelqu'un » se dit \all{jemanden vermissen} ; « rater le bus » se dit \all{den Bus verpassen}.
\end{corrige}

\begin{corrige}[Exercice 9 — Compréhension de l'écrit : texte et questions]
\textbf{Réponses attendues} (formulations équivalentes acceptées) :
\begin{enumerate}
\item \all{Die Schüler lernen heute mit Tablets und benutzen digitale Plattformen. Sie schicken auch ihre Hausaufgaben per E-Mail.}\\
« Les élèves apprennent aujourd'hui avec des tablettes et utilisent des plateformes numériques. Ils envoient aussi leurs devoirs par courriel. »
\item \all{Eine Chance ist, dass die Schüler schneller Informationen finden. Eine andere Chance ist, dass sie auch zu Hause arbeiten können, wenn sie krank sind.}\\
« Une chance : les élèves trouvent plus vite des informations. Une autre : ils peuvent travailler à la maison quand ils sont malades. »
\item \all{Ein Problem ist, dass nicht jede Familie einen guten Computer oder eine schnelle Internetverbindung hat. Ein anderes Problem ist, dass manche Schüler viele Stunden vor dem Bildschirm sitzen und sich zu wenig bewegen.}\\
« Un problème : toutes les familles n'ont pas un bon ordinateur ni une connexion rapide. Un autre : certains élèves passent de longues heures devant l'écran et ne bougent pas assez. »
\item \all{Man stellte fest, dass viele Geräte nach einem Jahr schon kaputt waren.}\\
« On a constaté que beaucoup d'appareils étaient déjà en panne au bout d'un an. »
\item \all{Die Experten sagen, dass die Technik den Lehrer nicht ersetzen darf, sondern ihm helfen muss.}\\
« Les experts disent que la technique ne doit pas remplacer le professeur, mais qu'elle doit l'aider. »
\item Proposition relevée : \all{Nachdem eine Schule in Bayern neue Tablets gekauft hatte, stellte man fest, dass...}\\
Dans la subordonnée introduite par \all{nachdem}, on emploie le \all{Plusquamperfekt} (\all{gekauft hatte}) ; dans la principale, le \all{Präteritum} (\all{stellte ... fest}). Le \all{Plusquamperfekt} marque l'antériorité : l'action de la subordonnée est achevée avant celle de la principale.
\item La structure d'opposition : \all{Einerseits bietet diese Entwicklung große Chancen, ... Andererseits gibt es ernste Probleme.}
\end{enumerate}
\textbf{Barème indicatif} (sur 20) : questions 1 à 5 : 2 points chacune (1 point pour l'information, 1 point pour la correction de la langue) ; questions 6 et 7 : 3 points chacune ; compréhension globale et qualité de l'expression : 4 points.
\end{corrige}

\begin{corrige}[Exercice 10 — Thème et version]
\textbf{A. Thème.}
\begin{enumerate}
\item \all{Einerseits ermöglicht der Fernunterricht, Zeit zu sparen, andererseits fühlen sich viele Schüler einsam.}\\
« Permettre de » = \all{ermöglichen + Infinitivsatz} ; « se sentir seul » = \all{sich einsam fühlen}.
\item \all{Nachdem der Lehrer die Lektion erklärt hatte, stellten die Schüler Fragen.}\\
« Poser une question » = \all{eine Frage stellen} (et non \all{fragen} seul, qui signifie « demander »).
\item \all{Als die Verbindung wiederhergestellt worden war, konnten wir den Kurs fortsetzen.}\\
Passé unique → \all{als} ; passif au \all{Plusquamperfekt} : \all{worden war} ; « continuer » = \all{fortsetzen} (particule séparable).
\item \all{Einerseits arbeitet mein Bruder gern auf dem Tablet, andererseits tun ihm oft die Augen weh.}\\
« Avoir mal aux yeux » = \all{die Augen tun ihm weh} (construction avec le datif de la personne).
\item \all{Nachdem der Strom abgeschaltet worden war, konnte ich meine Hausaufgaben nicht mehr schicken.}
\end{enumerate}
\textbf{B. Version.}\\
« D'une part, les cours à distance ont été pour moi une bonne expérience, d'autre part j'ai eu beaucoup de problèmes techniques. Après que le courant eut été coupé, je n'ai pas pu travailler pendant deux heures. Après que j'eus enfin trouvé une bonne connexion, j'ai envoyé mes devoirs au professeur. »\\
\textbf{Points de vigilance} : \all{war ausgefallen} se traduit par un plus-que-parfait (« eut été coupé / était tombé en panne ») ; \all{zwei Stunden lang} = « pendant deux heures » ; \all{an den Lehrer schicken} = « envoyer au professeur ».\\
\textbf{Barème indicatif} (sur 20) : thème 12 points (2 à 3 points par phrase, pénalité pour faute de place du verbe, de temps ou de cas) ; version 8 points (fidélité du sens 5 points, qualité du français 3 points).
\end{corrige}

\begin{corrige}[Exercice 11 — Rédaction guidée type Bac]
\textbf{Proposition de rédaction modèle} (environ 95 mots) :

\all{In meinem Land, Kamerun, hatten viele Schulen während einer schwierigen Zeit Fernunterricht. Einerseits war das eine gute Lösung : die Schüler mussten nicht zur Schule fahren und konnten zu Hause weiterlernen. Andererseits gab es große Probleme : viele Familien haben keinen Computer, und der Strom fällt oft aus. Nachdem der Strom eines Tages ausgefallen war, konnte ich drei Stunden lang nicht am Kurs teilnehmen. Meiner Meinung nach ist der Fernunterricht nützlich, aber er kann den Kontakt mit dem Lehrer und den Freunden nicht ersetzen. Die Schule bleibt für mich der beste Ort zum Lernen.}

\textbf{Traduction du modèle} : « Dans mon pays, le Cameroun, beaucoup d'écoles ont fait des cours à distance pendant une période difficile. D'une part, c'était une bonne solution : les élèves n'avaient pas besoin d'aller à l'école et pouvaient continuer à apprendre à la maison. D'autre part, il y avait de grands problèmes : beaucoup de familles n'ont pas d'ordinateur et le courant tombe souvent en panne. Un jour, après que le courant eut été coupé, je n'ai pas pu participer au cours pendant trois heures. À mon avis, les cours à distance sont utiles, mais ils ne peuvent pas remplacer le contact avec le professeur et les amis. L'école reste pour moi le meilleur endroit pour apprendre. »

\textbf{Barème indicatif} (sur 20) :
\begin{itemize}
\item Respect de la tâche (2 avantages, 2 inconvénients, avis personnel) : 5 points ;
\item Emploi correct de \all{einerseits … andererseits} avec inversion : 3 points ;
\item Emploi correct de \all{nachdem} + \all{Plusquamperfekt} : 3 points ;
\item Correction grammaticale (place du verbe, cas, conjugaisons) : 4 points ;
\item Vocabulaire du thème et orthographe (majuscules, Umlaute) : 3 points ;
\item Cohérence et longueur (80–100 mots) : 2 points.
\end{itemize}
\textbf{Points de vigilance} :
\begin{itemize}
\item Après \all{einerseits / andererseits} en tête : verbe en 2\textsuperscript{e} position (\all{Einerseits war das...}, jamais \all{Einerseits das war...}).
\item \all{Nachdem} renvoie le verbe conjugué en fin de subordonnée : \all{Nachdem der Strom ausgefallen war, ...} ; la principale suit avec inversion : \all{..., konnte ich...}.
\item Auxiliaire \all{sein} au \all{Plusquamperfekt} avec \all{ausfallen} : \all{war ausgefallen}.
\item « Participer à » = \all{an + Dativ teilnehmen} : \all{am Kurs teilnehmen} (particule séparable).
\item Majuscule à tous les noms : \all{der Strom, die Schule, das Problem, der Fernunterricht}.
\item « À mon avis » = \all{meiner Meinung nach} (tournure la plus idiomatique).
\end{itemize}
\end{corrige}

\newpage

\coursec{Fiche bilan}

\begin{popfigure}
\begin{tikzpicture}[
  mindmap,
  grow cyclic,
  every node/.style={concept, align=center, font=\small},
  root concept/.append style={concept color=popblue, font=\bfseries\small, text=white},
  level 1 concept/.append style={level distance=3.6cm, sibling angle=60, font=\footnotesize},
  level 2 concept/.append style={level distance=2.4cm, font=\scriptsize},
]
\node[root concept, align=center]{Fernunterricht\\ Chancen und\\ Probleme}
  child[concept color=poporange]{ node[align=center]{Lexique\\ der Fernunterricht}
    child { node[concept color=poporangeL, text=popdark, align=center]{die Verbindung,\\ das Netz,\\ der Strom} }
    child { node[concept color=poporangeL, text=popdark, align=center]{der Vorteil /\\ der Nachteil} }
  }
  child[concept color=popgreen]{ node[align=center]{Grammaire 1\\ Opposition}
    child { node[concept color=popgreenL, text=popdark, align=center]{einerseits \ldots\\ andererseits} }
  }
  child[concept color=poppurple]{ node[align=center]{Grammaire 2\\ Plus-que-parfait}
    child { node[concept color=poppurpleL, text=popdark, align=center]{hatte / war\\ + Partizip II} }
    child { node[concept color=poppurpleL, text=popdark, align=center]{nachdem,\\ bevor, als} }
  }
  child[concept color=popgold]{ node[align=center]{Compréhension\\ du texte}
    child { node[concept color=popgoldL, text=popdark, align=center]{global puis\\ détaillé} }
  }
  child[concept color=poppink]{ node[align=center]{Expression\\ guidée}
    child { node[concept color=poppinkL, text=popdark, align=center]{résumé,\\ questions,\\ paragraphe} }
  }
  child[concept color=popblue]{ node[align=center]{Méthode\\ commentaire}
    child { node[concept color=popblueL, text=popdark, align=center]{lire -- relever\\ -- analyser\\ -- rédiger} }
  };
\end{tikzpicture}
\legende{Carte mentale de la leçon : les cours à distance et les problèmes du numérique}
\end{popfigure}

\begin{bilan}
\textbf{1. Lexique clé (à connaître avec article et pluriel)}
\begin{itemize}\setlength{\itemsep}{1pt}
  \item \all{der Fernunterricht} (--) : les cours à distance ; \all{die Verbindung, -en} : la connexion ; \all{das Netz, -e} : le réseau ; \all{der Strom} (--) : l'électricité, le courant ; \all{der Stromausfall, "-fälle} : la coupure de courant.
  \item \all{der Vorteil, -e} : l'avantage ; \all{der Nachteil, -e} : l'inconvénient ; \all{der Bildschirm, -e} : l'écran ; \all{das Handy, -s} : le téléphone portable ; \all{die E-Mail, -s} : le courriel.
  \item Verbes : \all{sich anmelden} (se connecter), \all{herunterladen} (télécharger), \all{teilen} (partager), \all{funktionieren} (fonctionner), \all{stören} (déranger).
  \item Expressions : \all{an einem Kurs teilnehmen} (suivre un cours), \all{Zugang zum Internet haben} (avoir accès à Internet), \all{zu Hause lernen} (apprendre à la maison).
\end{itemize}

\textbf{2. Grammaire à connaître par c\oe ur}
\begin{center}
\begin{tabularx}{\linewidth}{|l|X|X|}
\hline
\rowcolor{popblueL} \textbf{Point} & \textbf{Règle} & \textbf{Exemple} \\
\hline
Opposition & \all{einerseits \ldots\ andererseits} encadre deux aspects contraires ; place libre, verbe en 2\textsuperscript{e} position & \all{Einerseits spart man Zeit, andererseits braucht man ein gutes Netz.} \\
\hline
Plus-que-parfait & \all{hatte / war} (Präteritum de \all{haben/sein}) + Partizip II à la fin & \all{Ich hatte die E-Mail geschrieben.} / \all{Er war nach Hause gegangen.} \\
\hline
Auxiliaire & \all{sein} avec les verbes de mouvement ou de changement d'état ; \all{haben} ailleurs & \all{Sie war eingeschlafen.} mais \all{Sie hatte gelernt.} \\
\hline
Proposition temporelle & conjonction (\all{nachdem, bevor, als, während}) $\rightarrow$ verbe conjugué à la fin & \all{Nachdem der Strom ausgefallen war, konnte ich nicht mehr lernen.} \\
\hline
Concordance & \all{nachdem} + Plus-que-parfait $\rightarrow$ principale au Präteritum (antériorité) & \all{Nachdem er sich angemeldet hatte, begann der Kurs.} \\
\hline
\end{tabularx}
\end{center}

\textbf{3. Méthodes express}
\begin{itemize}\setlength{\itemsep}{1pt}
  \item Compréhension : lire le titre et la source $\rightarrow$ lecture globale (idée générale) $\rightarrow$ lecture détaillée (surligner le lexique du thème) $\rightarrow$ répondre en phrases complètes.
  \item Résumé : une phrase par paragraphe, au Präsens ou au Perfekt, sans exemple ni citation.
  \item Production : introduire (\all{In diesem Text geht es um \ldots}), opposer avec \all{einerseits \ldots\ andererseits}, conclure avec son avis (\all{Meiner Meinung nach \ldots}).
\end{itemize}
\end{bilan}

\begin{aretenir}[Checklist avant le Bac]
\begin{itemize}\setlength{\itemsep}{2pt}
  \item[$\square$] Je connais l'article et le pluriel de \all{der Vorteil, der Nachteil, die Verbindung, das Netz}.
  \item[$\square$] Je sais employer \all{einerseits \ldots\ andererseits} avec le verbe en deuxième position.
  \item[$\square$] Je forme le plus-que-parfait : \all{hatte/war} + Partizip II placé à la fin.
  \item[$\square$] Je choisis le bon auxiliaire : \all{sein} pour les mouvements et changements d'état.
  \item[$\square$] Je mets le verbe conjugué à la fin après \all{nachdem, bevor, als, während}.
  \item[$\square$] J'applique la concordance : \all{nachdem} + plus-que-parfait, principale au Präteritum.
  \item[$\square$] Je n'oublie pas la majuscule à tous les noms allemands (\all{das Netz, der Strom}).
  \item[$\square$] Je distingue les faux amis : \all{der Strom} = électricité, pas « courant d'eau » au sens figuré.
  \item[$\square$] Je sais résumer un texte en trois ou quatre phrases sans recopier.
  \item[$\square$] Je sais exprimer mon avis : \all{Meiner Meinung nach hat der Fernunterricht viele Nachteile.}
  \item[$\square$] Je relis ma production : place du verbe, cas, orthographe des Umlaute.
\end{itemize}
\end{aretenir}

\findecours

\end{document}
